% hecke-catalog.tex — COMPLETE reference catalog of the irreducible
% representations of the Hecke algebras in catalogue.pdf.
% One section per graph: Sabidussi pair (G,H), Wedderburn decomposition,
% table of all 1-dimensional characters, and for every block of dimension
% >= 2: generators, complete relations, explicit matrices for ALL orbital
% generators, and the expansion of every rho(T_i) in the block generators.
% Data blocks are machine-generated: hecke_catalog_full.g -> catalog-data/.
% Do not edit catalog-data/*.tex by hand.
\documentclass[11pt]{article}
\usepackage[margin=2.4cm]{geometry}
\usepackage{amsmath,amssymb,amsthm}
\usepackage[hidelinks]{hyperref}
\allowdisplaybreaks

\newcommand{\C}{\mathbb{C}}
\newcommand{\Q}{\mathbb{Q}}
\newcommand{\Z}{\mathbb{Z}}
\newcommand{\HH}{\mathcal{H}}
\newcommand{\mat}[1]{\begin{pmatrix}#1\end{pmatrix}}
\newtheorem{remark}{Remark}

\title{A complete catalog of Hecke-algebra representations\\
\large for the graphs of \texttt{catalogue.pdf}}
\author{Fable (Claude), with Robin Langer}
\date{June 2026}

\begin{document}
\maketitle

\noindent
For each vertex-transitive graph $\Gamma$ in \texttt{catalogue.pdf} whose
Hecke algebra is noncommutative, this catalog records:
\begin{itemize}
\item the \emph{Sabidussi pair} $(G,H)$: $G=\operatorname{Aut}\Gamma$ (or the
  chosen transitive subgroup) and $H=\operatorname{Stab}_G(v)$, so that
  $\Gamma$ is a coset graph on $G/H$;
\item the Wedderburn decomposition of
  $\HH=\HH(G,H)=\operatorname{End}_G\C[G/H]$;
\item the table of \emph{all} one-dimensional characters
  $T_i\mapsto\lambda_i$ (columns labelled by the $G$-irreducible
  $\chi$ they sit on, with its degree);
\item for \emph{every} irreducible representation of dimension $m\ge2$:
  \begin{itemize}
  \item a minimal generating set $A,B,\dots$ chosen among the orbital
    images $\rho(T_i)$, as explicit matrices;
  \item \emph{complete} defining relations: the minimal polynomial of each
    generator, together with rewriting rules expressing $w\cdot g$ in the
    monomial basis for every basis word $w$ and generator $g$ --- a
    presentation of $M_m(\C)$ from which any word in the generators
    reduces;
  \item the explicit matrix of \emph{every} orbital generator
    $\rho(T_1),\dots,\rho(T_r)$, each accompanied by its expansion as a
    linear combination of words in the block generators.
  \end{itemize}
\end{itemize}
Throughout, $T_1=I,T_2,\dots,T_r$ is the orbital basis of $\HH$, sorted by
suborbit length $k_i$ (ties among equal valencies are ordered by the orbit
enumeration of the generating GAP session; the numbering is therefore fixed
unambiguously by the explicit matrices below, but may differ from earlier
write-ups within a tied valency class). All matrices were computed exactly
over cyclotomic fields by isotypic restriction
(\texttt{hecke\_catalog\_full.g}) and verified by reconstructing the
$m^2$-dimensional matrix algebra, checking $\rho(\sum_i T_i)=0$ on
non-principal blocks, and re-multiplying every word expansion.

\paragraph{Notation.}
$\eta_a=\sum_{s\in H_p}\zeta_p^{as}$ denotes the sextic Gauss period of
$\Q(\zeta_p)$ over the index-$6$ subgroup $H_p\le(\Z/p)^\times$;
$c_k=\zeta_7^k+\zeta_7^{-k}$ (so $c_1+c_2+c_3=-1$);
$\alpha=\frac{-1+\sqrt{-7}}{2}$ (so $\alpha\bar\alpha=2$,
$\alpha+\bar\alpha=-1$); $\omega=e^{2\pi i/3}$ (so
$\omega\bar\omega=1$, $\omega+\bar\omega=-1$).

\tableofcontents

%──────────────────────────────────────────────────────────────────────
\section{Dodecahedron, refined: $\bigl(\mathrm{PSL}(2,5),\,C_3\bigr)$}

\textbf{Graph.} The dodecahedron on $20=|{\mathrm{PSL}(2,5)}|/3$ vertices.
Under the full group $C_2\times A_5$ the pair is Gelfand; the
representation theory below belongs to the refined pair
$G=\mathrm{PSL}(2,5)\cong A_5$, $H=C_3$.

\textbf{Algebra.} Rank $8$, $\HH\cong\C^4\oplus M_2(\C)$.

\textbf{Highlights.} The valency-3 generator satisfies $B^2+B+3I=0$
($\det B=3=$ valency); the golden $\pm\sqrt5$ of the dodecahedron spectrum
lives entirely in the 1-dimensional characters below. The ``why-$N$''
constant of this 20-point Meinhold miniature is $4$:
$(\rho(T_6)-I)^2=4I$.

% auto-generated by hecke_catalog_full.g — do not edit
\paragraph{One-dimensional characters.}
\[
\begin{array}{c|cccc}
 & \chi_{1}\,(1) & \chi_{2}\,(3) & \chi_{3}\,(3) & \chi_{5}\,(5) \\\hline
T_{1} & 1 & 1 & 1 & 1 \\
T_{2} & 1 & -1 & -1 & 1 \\
T_{3} & 3 & -\sqrt5 & \sqrt5 & 1 \\
T_{4} & 3 & \sqrt5 & -\sqrt5 & 1 \\
T_{5} & 3 & 1 & 1 & -1 \\
T_{6} & 3 & -1 & -1 & -1 \\
T_{7} & 3 & 1 & 1 & -1 \\
T_{8} & 3 & -1 & -1 & -1 \\
\end{array}
\]

\subsubsection*{Block 1: $M_2(\C)$ on $\chi_{4}$ (degree 4)}
\noindent\textbf{Generators.}
\[
A = \rho(T_{2}) = \begin{pmatrix} 1 & 0 \\ 0 & -1 \end{pmatrix}, \qquad
B = \rho(T_{5}) = \begin{pmatrix} \tfrac{1}{2} & 1 \\ -\tfrac{15}{4} & -\tfrac{3}{2} \end{pmatrix} .
\]
\textbf{Complete relations} (monomial basis $I,\;A,\;B,\;AB$ of $M_2(\C)$):
\begin{gather*}
A^2-I = 0 ,\\
B^2+B+3I = 0 ,\\
BA = 2I-A-AB ,\\
ABA = -I+2\,A-B ,\\
ABB = -3\,A-AB .
\end{gather*}
\textbf{All orbital images} ($\rho(T_i)$, valency $k_i$, word in the generators):
\begin{align*}
\rho(T_{1}) &= \begin{pmatrix} 1 & 0 \\ 0 & 1 \end{pmatrix} = I & (k_{1} &= 1) \\
\rho(T_{2}) &= \begin{pmatrix} 1 & 0 \\ 0 & -1 \end{pmatrix} = A & (k_{2} &= 1) \\
\rho(T_{3}) &= \begin{pmatrix} -2 & 0 \\ 0 & 0 \end{pmatrix} = -I-A & (k_{3} &= 3) \\
\rho(T_{4}) &= \begin{pmatrix} -2 & 0 \\ 0 & 0 \end{pmatrix} = -I-A & (k_{4} &= 3) \\
\rho(T_{5}) &= \begin{pmatrix} \tfrac{1}{2} & 1 \\ -\tfrac{15}{4} & -\tfrac{3}{2} \end{pmatrix} = B & (k_{5} &= 3) \\
\rho(T_{6}) &= \begin{pmatrix} \tfrac{1}{2} & -1 \\ -\tfrac{15}{4} & \tfrac{3}{2} \end{pmatrix} = 2I-A-AB & (k_{6} &= 3) \\
\rho(T_{7}) &= \begin{pmatrix} \tfrac{1}{2} & -1 \\ \tfrac{15}{4} & -\tfrac{3}{2} \end{pmatrix} = -I+2\,A-B & (k_{7} &= 3) \\
\rho(T_{8}) &= \begin{pmatrix} \tfrac{1}{2} & 1 \\ \tfrac{15}{4} & \tfrac{3}{2} \end{pmatrix} = AB & (k_{8} &= 3) .
\end{align*}



%──────────────────────────────────────────────────────────────────────
\section{F26A: $\bigl(C_{13}{:}C_6,\,C_3\bigr)$}

\textbf{Graph.} F26A, the unique cubic arc-transitive graph on 26 vertices;
$G=C_{13}{:}C_6$ ($|G|=78$) is \emph{sharply} arc-transitive, $H=C_3$.

\textbf{Algebra.} Rank $10$,
$\HH\cong\C^2\oplus M_2(\C)\oplus M_2(\C)$, the two blocks Galois conjugate
over $\Q(\sqrt{13})$ (the substitution $\sqrt{13}\mapsto-\sqrt{13}$ maps
Block~1 to Block~2).

\textbf{Highlights.} With $\eta_0,\eta_1=\frac{-1\pm\sqrt{13}}{2}$ the
sextic Gauss periods of $\Q(\zeta_{13})$, the traceless generators satisfy
\[
B^2=\frac{5-\sqrt{13}}{2}\,I \quad(\text{Block 1}),\qquad
B^2=\frac{5+\sqrt{13}}{2}\,I \quad(\text{Block 2}),\qquad
N_+N_-=3=\text{valency},
\]
and $\frac{5\mp\sqrt{13}}{2}=3+\eta$: \emph{valency plus a Gauss period}
(see the Remark after CR(31)).

% auto-generated by hecke_catalog_full.g — do not edit
\paragraph{One-dimensional characters.}
\[
\begin{array}{c|cc}
 & \chi_{1}\,(1) & \chi_{2}\,(1) \\\hline
T_{1} & 1 & 1 \\
T_{2} & 1 & -1 \\
T_{3} & 3 & -3 \\
T_{4} & 3 & 3 \\
T_{5} & 3 & -3 \\
T_{6} & 3 & 3 \\
T_{7} & 3 & -3 \\
T_{8} & 3 & 3 \\
T_{9} & 3 & -3 \\
T_{10} & 3 & 3 \\
\end{array}
\]

\subsubsection*{Block 1: $M_2(\C)$ on $\chi_{7}$ (degree 6)}
\noindent\textbf{Generators.}
\[
A = \rho(T_{2}) = \begin{pmatrix} 1 & 0 \\ 0 & -1 \end{pmatrix}, \qquad
B = \rho(T_{3}) = \begin{pmatrix} -\tfrac{1}{4}+\tfrac{1}{4}\sqrt{13} & 1 \\ \tfrac{13}{8}-\tfrac{3}{8}\sqrt{13} & \tfrac{1}{4}-\tfrac{1}{4}\sqrt{13} \end{pmatrix} .
\]
\textbf{Complete relations} (monomial basis $I,\;A,\;B,\;AB$ of $M_2(\C)$):
\begin{gather*}
A^2-I = 0 ,\\
B^2+\bigl(-\tfrac{5}{2}+\tfrac{1}{2}\sqrt{13}\bigr)I = 0 ,\\
BA = \bigl(-\tfrac{1}{2}+\tfrac{1}{2}\sqrt{13}\bigr)I-AB ,\\
ABA = \bigl(-\tfrac{1}{2}+\tfrac{1}{2}\sqrt{13}\bigr)A-B ,\\
ABB = \bigl(\tfrac{5}{2}-\tfrac{1}{2}\sqrt{13}\bigr)A .
\end{gather*}
\textbf{All orbital images} ($\rho(T_i)$, valency $k_i$, word in the generators):
\begin{align*}
\rho(T_{1}) &= \begin{pmatrix} 1 & 0 \\ 0 & 1 \end{pmatrix} = I & (k_{1} &= 1) \\
\rho(T_{2}) &= \begin{pmatrix} 1 & 0 \\ 0 & -1 \end{pmatrix} = A & (k_{2} &= 1) \\
\rho(T_{3}) &= \begin{pmatrix} -\tfrac{1}{4}+\tfrac{1}{4}\sqrt{13} & 1 \\ \tfrac{13}{8}-\tfrac{3}{8}\sqrt{13} & \tfrac{1}{4}-\tfrac{1}{4}\sqrt{13} \end{pmatrix} = B & (k_{3} &= 3) \\
\rho(T_{4}) &= \begin{pmatrix} -\tfrac{1}{4}-\tfrac{1}{4}\sqrt{13} & \tfrac{3}{2}+\tfrac{1}{2}\sqrt{13} \\ -\tfrac{1}{4}\sqrt{13} & -\tfrac{1}{4}-\tfrac{1}{4}\sqrt{13} \end{pmatrix} & (k_{4} &= 3) \\
 &= \bigl(-\tfrac{3}{2}-\tfrac{1}{2}\sqrt{13}\bigr)I+\bigl(\tfrac{3}{2}+\tfrac{1}{2}\sqrt{13}\bigr)AB && \\
\rho(T_{5}) &= \begin{pmatrix} -\tfrac{1}{4}+\tfrac{1}{4}\sqrt{13} & -1 \\ -\tfrac{13}{8}+\tfrac{3}{8}\sqrt{13} & \tfrac{1}{4}-\tfrac{1}{4}\sqrt{13} \end{pmatrix} = \bigl(-\tfrac{1}{2}+\tfrac{1}{2}\sqrt{13}\bigr)A-B & (k_{5} &= 3) \\
\rho(T_{6}) &= \begin{pmatrix} -\tfrac{1}{4}+\tfrac{1}{4}\sqrt{13} & -1 \\ \tfrac{13}{8}-\tfrac{3}{8}\sqrt{13} & -\tfrac{1}{4}+\tfrac{1}{4}\sqrt{13} \end{pmatrix} = \bigl(-\tfrac{1}{2}+\tfrac{1}{2}\sqrt{13}\bigr)I-AB & (k_{6} &= 3) \\
\rho(T_{7}) &= \begin{pmatrix} -\tfrac{1}{4}-\tfrac{1}{4}\sqrt{13} & -\tfrac{3}{2}-\tfrac{1}{2}\sqrt{13} \\ -\tfrac{1}{4}\sqrt{13} & \tfrac{1}{4}+\tfrac{1}{4}\sqrt{13} \end{pmatrix} = A+\bigl(-\tfrac{3}{2}-\tfrac{1}{2}\sqrt{13}\bigr)B & (k_{7} &= 3) \\
\rho(T_{8}) &= \begin{pmatrix} -\tfrac{1}{4}-\tfrac{1}{4}\sqrt{13} & -\tfrac{3}{2}-\tfrac{1}{2}\sqrt{13} \\ \tfrac{1}{4}\sqrt{13} & -\tfrac{1}{4}-\tfrac{1}{4}\sqrt{13} \end{pmatrix} = I+\bigl(-\tfrac{3}{2}-\tfrac{1}{2}\sqrt{13}\bigr)AB & (k_{8} &= 3) \\
\rho(T_{9}) &= \begin{pmatrix} -\tfrac{1}{4}-\tfrac{1}{4}\sqrt{13} & \tfrac{3}{2}+\tfrac{1}{2}\sqrt{13} \\ \tfrac{1}{4}\sqrt{13} & \tfrac{1}{4}+\tfrac{1}{4}\sqrt{13} \end{pmatrix} & (k_{9} &= 3) \\
 &= \bigl(-\tfrac{3}{2}-\tfrac{1}{2}\sqrt{13}\bigr)A+\bigl(\tfrac{3}{2}+\tfrac{1}{2}\sqrt{13}\bigr)B && \\
\rho(T_{10}) &= \begin{pmatrix} -\tfrac{1}{4}+\tfrac{1}{4}\sqrt{13} & 1 \\ -\tfrac{13}{8}+\tfrac{3}{8}\sqrt{13} & -\tfrac{1}{4}+\tfrac{1}{4}\sqrt{13} \end{pmatrix} = AB & (k_{10} &= 3) .
\end{align*}

\subsubsection*{Block 2: $M_2(\C)$ on $\chi_{8}$ (degree 6)}
\noindent\textbf{Generators.}
\[
A = \rho(T_{2}) = \begin{pmatrix} 1 & 0 \\ 0 & -1 \end{pmatrix}, \qquad
B = \rho(T_{3}) = \begin{pmatrix} -\tfrac{1}{4}-\tfrac{1}{4}\sqrt{13} & 1 \\ \tfrac{13}{8}+\tfrac{3}{8}\sqrt{13} & \tfrac{1}{4}+\tfrac{1}{4}\sqrt{13} \end{pmatrix} .
\]
\textbf{Complete relations} (monomial basis $I,\;A,\;B,\;AB$ of $M_2(\C)$):
\begin{gather*}
A^2-I = 0 ,\\
B^2+\bigl(-\tfrac{5}{2}-\tfrac{1}{2}\sqrt{13}\bigr)I = 0 ,\\
BA = \bigl(-\tfrac{1}{2}-\tfrac{1}{2}\sqrt{13}\bigr)I-AB ,\\
ABA = \bigl(-\tfrac{1}{2}-\tfrac{1}{2}\sqrt{13}\bigr)A-B ,\\
ABB = \bigl(\tfrac{5}{2}+\tfrac{1}{2}\sqrt{13}\bigr)A .
\end{gather*}
\textbf{All orbital images} ($\rho(T_i)$, valency $k_i$, word in the generators):
\begin{align*}
\rho(T_{1}) &= \begin{pmatrix} 1 & 0 \\ 0 & 1 \end{pmatrix} = I & (k_{1} &= 1) \\
\rho(T_{2}) &= \begin{pmatrix} 1 & 0 \\ 0 & -1 \end{pmatrix} = A & (k_{2} &= 1) \\
\rho(T_{3}) &= \begin{pmatrix} -\tfrac{1}{4}-\tfrac{1}{4}\sqrt{13} & 1 \\ \tfrac{13}{8}+\tfrac{3}{8}\sqrt{13} & \tfrac{1}{4}+\tfrac{1}{4}\sqrt{13} \end{pmatrix} = B & (k_{3} &= 3) \\
\rho(T_{4}) &= \begin{pmatrix} -\tfrac{1}{4}+\tfrac{1}{4}\sqrt{13} & \tfrac{3}{2}-\tfrac{1}{2}\sqrt{13} \\ \tfrac{1}{4}\sqrt{13} & -\tfrac{1}{4}+\tfrac{1}{4}\sqrt{13} \end{pmatrix} & (k_{4} &= 3) \\
 &= \bigl(-\tfrac{3}{2}+\tfrac{1}{2}\sqrt{13}\bigr)I+\bigl(\tfrac{3}{2}-\tfrac{1}{2}\sqrt{13}\bigr)AB && \\
\rho(T_{5}) &= \begin{pmatrix} -\tfrac{1}{4}-\tfrac{1}{4}\sqrt{13} & -1 \\ -\tfrac{13}{8}-\tfrac{3}{8}\sqrt{13} & \tfrac{1}{4}+\tfrac{1}{4}\sqrt{13} \end{pmatrix} = \bigl(-\tfrac{1}{2}-\tfrac{1}{2}\sqrt{13}\bigr)A-B & (k_{5} &= 3) \\
\rho(T_{6}) &= \begin{pmatrix} -\tfrac{1}{4}-\tfrac{1}{4}\sqrt{13} & -1 \\ \tfrac{13}{8}+\tfrac{3}{8}\sqrt{13} & -\tfrac{1}{4}-\tfrac{1}{4}\sqrt{13} \end{pmatrix} = \bigl(-\tfrac{1}{2}-\tfrac{1}{2}\sqrt{13}\bigr)I-AB & (k_{6} &= 3) \\
\rho(T_{7}) &= \begin{pmatrix} -\tfrac{1}{4}+\tfrac{1}{4}\sqrt{13} & -\tfrac{3}{2}+\tfrac{1}{2}\sqrt{13} \\ \tfrac{1}{4}\sqrt{13} & \tfrac{1}{4}-\tfrac{1}{4}\sqrt{13} \end{pmatrix} = A+\bigl(-\tfrac{3}{2}+\tfrac{1}{2}\sqrt{13}\bigr)B & (k_{7} &= 3) \\
\rho(T_{8}) &= \begin{pmatrix} -\tfrac{1}{4}+\tfrac{1}{4}\sqrt{13} & -\tfrac{3}{2}+\tfrac{1}{2}\sqrt{13} \\ -\tfrac{1}{4}\sqrt{13} & -\tfrac{1}{4}+\tfrac{1}{4}\sqrt{13} \end{pmatrix} = I+\bigl(-\tfrac{3}{2}+\tfrac{1}{2}\sqrt{13}\bigr)AB & (k_{8} &= 3) \\
\rho(T_{9}) &= \begin{pmatrix} -\tfrac{1}{4}+\tfrac{1}{4}\sqrt{13} & \tfrac{3}{2}-\tfrac{1}{2}\sqrt{13} \\ -\tfrac{1}{4}\sqrt{13} & \tfrac{1}{4}-\tfrac{1}{4}\sqrt{13} \end{pmatrix} & (k_{9} &= 3) \\
 &= \bigl(-\tfrac{3}{2}+\tfrac{1}{2}\sqrt{13}\bigr)A+\bigl(\tfrac{3}{2}-\tfrac{1}{2}\sqrt{13}\bigr)B && \\
\rho(T_{10}) &= \begin{pmatrix} -\tfrac{1}{4}-\tfrac{1}{4}\sqrt{13} & 1 \\ -\tfrac{13}{8}-\tfrac{3}{8}\sqrt{13} & -\tfrac{1}{4}-\tfrac{1}{4}\sqrt{13} \end{pmatrix} = AB & (k_{10} &= 3) .
\end{align*}



%──────────────────────────────────────────────────────────────────────
\section{CR(19): $\bigl(C_{19}{:}C_6,\,C_3\bigr)$}

\textbf{Graph.} The 38-vertex cubic arc-transitive Frobenius graph;
$G=C_{19}{:}C_6$ ($|G|=114$), $H=C_3$. Here
$H_{19}=\langle 8\rangle=\{1,7,8,11,12,18\}\le(\Z/19)^\times$ and the three
sextic periods are $\eta_1,\eta_2,\eta_4$ (indexed by coset representatives
$1,2,4$), satisfying $\eta_1+\eta_2+\eta_4=-1$; entries below are written
in the basis $\{1,\eta_a,\eta_b\}$ that gives the sparsest expression.

\textbf{Algebra.} Rank $14$, $\HH\cong\C^2\oplus M_2(\C)^{\oplus3}$.

\textbf{Highlights.} The three blocks are permuted cyclically by the Galois
action $a\mapsto2a$ on period indices; in each, the traceless valency-3
orbital squares to $(3+\eta)I$.

% auto-generated by hecke_catalog_full.g — do not edit
\paragraph{One-dimensional characters.}
\[
\begin{array}{c|cc}
 & \chi_{1}\,(1) & \chi_{2}\,(1) \\\hline
T_{1} & 1 & 1 \\
T_{2} & 1 & -1 \\
T_{3} & 3 & 3 \\
T_{4} & 3 & -3 \\
T_{5} & 3 & 3 \\
T_{6} & 3 & -3 \\
T_{7} & 3 & -3 \\
T_{8} & 3 & 3 \\
T_{9} & 3 & -3 \\
T_{10} & 3 & 3 \\
T_{11} & 3 & -3 \\
T_{12} & 3 & 3 \\
T_{13} & 3 & -3 \\
T_{14} & 3 & 3 \\
\end{array}
\]

\subsubsection*{Block 1: $M_2(\C)$ on $\chi_{7}$ (degree 6)}
\noindent\textbf{Generators.}
\[
A = \rho(T_{2}) = \begin{pmatrix} 1 & 0 \\ 0 & -1 \end{pmatrix}, \qquad
B = \rho(T_{3}) = \begin{pmatrix} \tfrac{1}{2}\eta_4 & 1 \\ -2-\tfrac{1}{4}\eta_1-\eta_2 & \tfrac{1}{2}\eta_4 \end{pmatrix} .
\]
\textbf{Complete relations} (monomial basis $I,\;A,\;B,\;AB$ of $M_2(\C)$):
\begin{gather*}
A^2-I = 0 ,\\
B^2-\eta_4\,B+\bigl(3+\eta_2\bigr)I = 0 ,\\
BA = \eta_4\,A-AB ,\\
ABA = \eta_4I-B ,\\
ABB = \bigl(-3-\eta_2\bigr)A+\eta_4\,AB .
\end{gather*}
\textbf{All orbital images} ($\rho(T_i)$, valency $k_i$, word in the generators):
\begin{align*}
\rho(T_{1}) &= \begin{pmatrix} 1 & 0 \\ 0 & 1 \end{pmatrix} = I & (k_{1} &= 1) \\
\rho(T_{2}) &= \begin{pmatrix} 1 & 0 \\ 0 & -1 \end{pmatrix} = A & (k_{2} &= 1) \\
\rho(T_{3}) &= \begin{pmatrix} \tfrac{1}{2}\eta_4 & 1 \\ -2-\tfrac{1}{4}\eta_1-\eta_2 & \tfrac{1}{2}\eta_4 \end{pmatrix} = B & (k_{3} &= 3) \\
\rho(T_{4}) &= \begin{pmatrix} \tfrac{1}{2}\eta_4 & -1 \\ -2-\tfrac{1}{4}\eta_1-\eta_2 & -\tfrac{1}{2}\eta_4 \end{pmatrix} = \eta_4\,A-AB & (k_{4} &= 3) \\
\rho(T_{5}) &= \begin{pmatrix} \tfrac{1}{2}\eta_1 & 2+\eta_4 \\ \tfrac{1}{4}\eta_2-\tfrac{1}{4}\eta_4 & \tfrac{1}{2}\eta_1 \end{pmatrix} = \bigl(-1+2\eta_1+\eta_2\bigr)I+\bigl(2+\eta_4\bigr)B & (k_{5} &= 3) \\
\rho(T_{6}) &= \begin{pmatrix} \tfrac{1}{2}\eta_1 & -2-\eta_4 \\ \tfrac{1}{4}\eta_2-\tfrac{1}{4}\eta_4 & -\tfrac{1}{2}\eta_1 \end{pmatrix} = \bigl(2+\eta_4\bigr)A+\bigl(-2-\eta_4\bigr)AB & (k_{6} &= 3) \\
\rho(T_{7}) &= \begin{pmatrix} \tfrac{1}{2}\eta_4 & 1 \\ 2+\tfrac{1}{4}\eta_1+\eta_2 & -\tfrac{1}{2}\eta_4 \end{pmatrix} = AB & (k_{7} &= 3) \\
\rho(T_{8}) &= \begin{pmatrix} \tfrac{1}{2}\eta_4 & -1 \\ 2+\tfrac{1}{4}\eta_1+\eta_2 & \tfrac{1}{2}\eta_4 \end{pmatrix} = \eta_4I-B & (k_{8} &= 3) \\
\rho(T_{9}) &= \begin{pmatrix} \tfrac{1}{2}\eta_1 & 2+\eta_4 \\ -\tfrac{1}{4}\eta_2+\tfrac{1}{4}\eta_4 & -\tfrac{1}{2}\eta_1 \end{pmatrix} & (k_{9} &= 3) \\
 &= \bigl(-1+2\eta_1+\eta_2\bigr)A+\bigl(2+\eta_4\bigr)AB && \\
\rho(T_{10}) &= \begin{pmatrix} \tfrac{1}{2}\eta_1 & -2-\eta_4 \\ -\tfrac{1}{4}\eta_2+\tfrac{1}{4}\eta_4 & \tfrac{1}{2}\eta_1 \end{pmatrix} = \bigl(2+\eta_4\bigr)I+\bigl(-2-\eta_4\bigr)B & (k_{10} &= 3) \\
\rho(T_{11}) &= \begin{pmatrix} \tfrac{1}{2}\eta_2 & -1-\eta_1 \\ -\tfrac{1}{4}\eta_1+\tfrac{1}{4}\eta_2 & -\tfrac{1}{2}\eta_2 \end{pmatrix} = \bigl(-1+\eta_2\bigr)A+\bigl(-1-\eta_1\bigr)AB & (k_{11} &= 3) \\
\rho(T_{12}) &= \begin{pmatrix} \tfrac{1}{2}\eta_2 & 1+\eta_1 \\ -\tfrac{1}{4}\eta_1+\tfrac{1}{4}\eta_2 & \tfrac{1}{2}\eta_2 \end{pmatrix} = I+\bigl(1+\eta_1\bigr)B & (k_{12} &= 3) \\
\rho(T_{13}) &= \begin{pmatrix} \tfrac{1}{2}\eta_2 & 1+\eta_1 \\ \tfrac{1}{4}\eta_1-\tfrac{1}{4}\eta_2 & -\tfrac{1}{2}\eta_2 \end{pmatrix} = A+\bigl(1+\eta_1\bigr)AB & (k_{13} &= 3) \\
\rho(T_{14}) &= \begin{pmatrix} \tfrac{1}{2}\eta_2 & -1-\eta_1 \\ \tfrac{1}{4}\eta_1-\tfrac{1}{4}\eta_2 & \tfrac{1}{2}\eta_2 \end{pmatrix} = \bigl(-1+\eta_2\bigr)I+\bigl(-1-\eta_1\bigr)B & (k_{14} &= 3) .
\end{align*}

\subsubsection*{Block 2: $M_2(\C)$ on $\chi_{8}$ (degree 6)}
\noindent\textbf{Generators.}
\[
A = \rho(T_{2}) = \begin{pmatrix} 1 & 0 \\ 0 & -1 \end{pmatrix}, \qquad
B = \rho(T_{3}) = \begin{pmatrix} \tfrac{1}{2}\eta_1 & 1 \\ -1+\eta_1+\tfrac{3}{4}\eta_2 & \tfrac{1}{2}\eta_1 \end{pmatrix} .
\]
\textbf{Complete relations} (monomial basis $I,\;A,\;B,\;AB$ of $M_2(\C)$):
\begin{gather*}
A^2-I = 0 ,\\
B^2-\eta_1\,B+\bigl(3+\eta_4\bigr)I = 0 ,\\
BA = \eta_1\,A-AB ,\\
ABA = \eta_1I-B ,\\
ABB = \bigl(-3-\eta_4\bigr)A+\eta_1\,AB .
\end{gather*}
\textbf{All orbital images} ($\rho(T_i)$, valency $k_i$, word in the generators):
\begin{align*}
\rho(T_{1}) &= \begin{pmatrix} 1 & 0 \\ 0 & 1 \end{pmatrix} = I & (k_{1} &= 1) \\
\rho(T_{2}) &= \begin{pmatrix} 1 & 0 \\ 0 & -1 \end{pmatrix} = A & (k_{2} &= 1) \\
\rho(T_{3}) &= \begin{pmatrix} \tfrac{1}{2}\eta_1 & 1 \\ -1+\eta_1+\tfrac{3}{4}\eta_2 & \tfrac{1}{2}\eta_1 \end{pmatrix} = B & (k_{3} &= 3) \\
\rho(T_{4}) &= \begin{pmatrix} \tfrac{1}{2}\eta_1 & -1 \\ -1+\eta_1+\tfrac{3}{4}\eta_2 & -\tfrac{1}{2}\eta_1 \end{pmatrix} = \eta_1\,A-AB & (k_{4} &= 3) \\
\rho(T_{5}) &= \begin{pmatrix} \tfrac{1}{2}\eta_2 & 2+\eta_1 \\ -\tfrac{1}{4}\eta_1+\tfrac{1}{4}\eta_4 & \tfrac{1}{2}\eta_2 \end{pmatrix} = \bigl(-2-\eta_1+\eta_2\bigr)I+\bigl(2+\eta_1\bigr)B & (k_{5} &= 3) \\
\rho(T_{6}) &= \begin{pmatrix} \tfrac{1}{2}\eta_2 & -2-\eta_1 \\ -\tfrac{1}{4}\eta_1+\tfrac{1}{4}\eta_4 & -\tfrac{1}{2}\eta_2 \end{pmatrix} = \bigl(2+\eta_1\bigr)A+\bigl(-2-\eta_1\bigr)AB & (k_{6} &= 3) \\
\rho(T_{7}) &= \begin{pmatrix} \tfrac{1}{2}\eta_1 & 1 \\ 1-\eta_1-\tfrac{3}{4}\eta_2 & -\tfrac{1}{2}\eta_1 \end{pmatrix} = AB & (k_{7} &= 3) \\
\rho(T_{8}) &= \begin{pmatrix} \tfrac{1}{2}\eta_1 & -1 \\ 1-\eta_1-\tfrac{3}{4}\eta_2 & \tfrac{1}{2}\eta_1 \end{pmatrix} = \eta_1I-B & (k_{8} &= 3) \\
\rho(T_{9}) &= \begin{pmatrix} \tfrac{1}{2}\eta_2 & 2+\eta_1 \\ \tfrac{1}{4}\eta_1-\tfrac{1}{4}\eta_4 & -\tfrac{1}{2}\eta_2 \end{pmatrix} = \bigl(-2-\eta_1+\eta_2\bigr)A+\bigl(2+\eta_1\bigr)AB & (k_{9} &= 3) \\
\rho(T_{10}) &= \begin{pmatrix} \tfrac{1}{2}\eta_2 & -2-\eta_1 \\ \tfrac{1}{4}\eta_1-\tfrac{1}{4}\eta_4 & \tfrac{1}{2}\eta_2 \end{pmatrix} = \bigl(2+\eta_1\bigr)I+\bigl(-2-\eta_1\bigr)B & (k_{10} &= 3) \\
\rho(T_{11}) &= \begin{pmatrix} \tfrac{1}{2}\eta_4 & -1-\eta_2 \\ -\tfrac{1}{4}\eta_2+\tfrac{1}{4}\eta_4 & -\tfrac{1}{2}\eta_4 \end{pmatrix} = \bigl(-1+\eta_4\bigr)A+\bigl(-1-\eta_2\bigr)AB & (k_{11} &= 3) \\
\rho(T_{12}) &= \begin{pmatrix} \tfrac{1}{2}\eta_4 & 1+\eta_2 \\ -\tfrac{1}{4}\eta_2+\tfrac{1}{4}\eta_4 & \tfrac{1}{2}\eta_4 \end{pmatrix} = I+\bigl(1+\eta_2\bigr)B & (k_{12} &= 3) \\
\rho(T_{13}) &= \begin{pmatrix} \tfrac{1}{2}\eta_4 & 1+\eta_2 \\ \tfrac{1}{4}\eta_2-\tfrac{1}{4}\eta_4 & -\tfrac{1}{2}\eta_4 \end{pmatrix} = A+\bigl(1+\eta_2\bigr)AB & (k_{13} &= 3) \\
\rho(T_{14}) &= \begin{pmatrix} \tfrac{1}{2}\eta_4 & -1-\eta_2 \\ \tfrac{1}{4}\eta_2-\tfrac{1}{4}\eta_4 & \tfrac{1}{2}\eta_4 \end{pmatrix} = \bigl(-1+\eta_4\bigr)I+\bigl(-1-\eta_2\bigr)B & (k_{14} &= 3) .
\end{align*}

\subsubsection*{Block 3: $M_2(\C)$ on $\chi_{9}$ (degree 6)}
\noindent\textbf{Generators.}
\[
A = \rho(T_{2}) = \begin{pmatrix} 1 & 0 \\ 0 & -1 \end{pmatrix}, \qquad
B = \rho(T_{3}) = \begin{pmatrix} \tfrac{1}{2}\eta_2 & 1 \\ -\tfrac{7}{4}-\tfrac{3}{4}\eta_1+\tfrac{1}{4}\eta_2 & \tfrac{1}{2}\eta_2 \end{pmatrix} .
\]
\textbf{Complete relations} (monomial basis $I,\;A,\;B,\;AB$ of $M_2(\C)$):
\begin{gather*}
A^2-I = 0 ,\\
B^2-\eta_2\,B+\bigl(3+\eta_1\bigr)I = 0 ,\\
BA = \eta_2\,A-AB ,\\
ABA = \eta_2I-B ,\\
ABB = \bigl(-3-\eta_1\bigr)A+\eta_2\,AB .
\end{gather*}
\textbf{All orbital images} ($\rho(T_i)$, valency $k_i$, word in the generators):
\begin{align*}
\rho(T_{1}) &= \begin{pmatrix} 1 & 0 \\ 0 & 1 \end{pmatrix} = I & (k_{1} &= 1) \\
\rho(T_{2}) &= \begin{pmatrix} 1 & 0 \\ 0 & -1 \end{pmatrix} = A & (k_{2} &= 1) \\
\rho(T_{3}) &= \begin{pmatrix} \tfrac{1}{2}\eta_2 & 1 \\ -\tfrac{7}{4}-\tfrac{3}{4}\eta_1+\tfrac{1}{4}\eta_2 & \tfrac{1}{2}\eta_2 \end{pmatrix} = B & (k_{3} &= 3) \\
\rho(T_{4}) &= \begin{pmatrix} \tfrac{1}{2}\eta_2 & -1 \\ -\tfrac{7}{4}-\tfrac{3}{4}\eta_1+\tfrac{1}{4}\eta_2 & -\tfrac{1}{2}\eta_2 \end{pmatrix} = \eta_2\,A-AB & (k_{4} &= 3) \\
\rho(T_{5}) &= \begin{pmatrix} \tfrac{1}{2}\eta_4 & 2+\eta_2 \\ \tfrac{1}{4}\eta_1-\tfrac{1}{4}\eta_2 & \tfrac{1}{2}\eta_4 \end{pmatrix} = \bigl(-3-\eta_1-2\eta_2\bigr)I+\bigl(2+\eta_2\bigr)B & (k_{5} &= 3) \\
\rho(T_{6}) &= \begin{pmatrix} \tfrac{1}{2}\eta_4 & -2-\eta_2 \\ \tfrac{1}{4}\eta_1-\tfrac{1}{4}\eta_2 & -\tfrac{1}{2}\eta_4 \end{pmatrix} = \bigl(2+\eta_2\bigr)A+\bigl(-2-\eta_2\bigr)AB & (k_{6} &= 3) \\
\rho(T_{7}) &= \begin{pmatrix} \tfrac{1}{2}\eta_2 & 1 \\ \tfrac{7}{4}+\tfrac{3}{4}\eta_1-\tfrac{1}{4}\eta_2 & -\tfrac{1}{2}\eta_2 \end{pmatrix} = AB & (k_{7} &= 3) \\
\rho(T_{8}) &= \begin{pmatrix} \tfrac{1}{2}\eta_2 & -1 \\ \tfrac{7}{4}+\tfrac{3}{4}\eta_1-\tfrac{1}{4}\eta_2 & \tfrac{1}{2}\eta_2 \end{pmatrix} = \eta_2I-B & (k_{8} &= 3) \\
\rho(T_{9}) &= \begin{pmatrix} \tfrac{1}{2}\eta_4 & 2+\eta_2 \\ -\tfrac{1}{4}\eta_1+\tfrac{1}{4}\eta_2 & -\tfrac{1}{2}\eta_4 \end{pmatrix} & (k_{9} &= 3) \\
 &= \bigl(-3-\eta_1-2\eta_2\bigr)A+\bigl(2+\eta_2\bigr)AB && \\
\rho(T_{10}) &= \begin{pmatrix} \tfrac{1}{2}\eta_4 & -2-\eta_2 \\ -\tfrac{1}{4}\eta_1+\tfrac{1}{4}\eta_2 & \tfrac{1}{2}\eta_4 \end{pmatrix} = \bigl(2+\eta_2\bigr)I+\bigl(-2-\eta_2\bigr)B & (k_{10} &= 3) \\
\rho(T_{11}) &= \begin{pmatrix} \tfrac{1}{2}\eta_1 & \eta_1+\eta_2 \\ \tfrac{1}{4}\eta_1-\tfrac{1}{4}\eta_4 & -\tfrac{1}{2}\eta_1 \end{pmatrix} = \bigl(-1+\eta_1\bigr)A+\bigl(\eta_1+\eta_2\bigr)AB & (k_{11} &= 3) \\
\rho(T_{12}) &= \begin{pmatrix} \tfrac{1}{2}\eta_1 & -\eta_1-\eta_2 \\ \tfrac{1}{4}\eta_1-\tfrac{1}{4}\eta_4 & \tfrac{1}{2}\eta_1 \end{pmatrix} = I+\bigl(-\eta_1-\eta_2\bigr)B & (k_{12} &= 3) \\
\rho(T_{13}) &= \begin{pmatrix} \tfrac{1}{2}\eta_1 & -\eta_1-\eta_2 \\ -\tfrac{1}{4}\eta_1+\tfrac{1}{4}\eta_4 & -\tfrac{1}{2}\eta_1 \end{pmatrix} = A+\bigl(-\eta_1-\eta_2\bigr)AB & (k_{13} &= 3) \\
\rho(T_{14}) &= \begin{pmatrix} \tfrac{1}{2}\eta_1 & \eta_1+\eta_2 \\ -\tfrac{1}{4}\eta_1+\tfrac{1}{4}\eta_4 & \tfrac{1}{2}\eta_1 \end{pmatrix} = \bigl(-1+\eta_1\bigr)I+\bigl(\eta_1+\eta_2\bigr)B & (k_{14} &= 3) .
\end{align*}



%──────────────────────────────────────────────────────────────────────
\section{CR(31): $\bigl(C_{31}{:}C_6,\,C_3\bigr)$}

{\sloppy
\textbf{Graph.} The 62-vertex sibling; $G=C_{31}{:}C_6$ ($|G|=186$),
$H=C_3$. Here $H_{31}=\langle 26\rangle=\{1,5,6,25,26,30\}$, and there
are five sextic periods $\eta_1,\eta_2,\eta_3,\eta_4,\eta_8$ (coset
representatives $1,2,3,4,8$), summing to $-1$.\par}

\textbf{Algebra.} Rank $22$, $\HH\cong\C^2\oplus M_2(\C)^{\oplus5}$.

\textbf{Highlights.} Five Galois-conjugate blocks, cycled by $a\mapsto2a$;
same $(3+\eta)$-law as CR(19).

% auto-generated by hecke_catalog_full.g — do not edit
\paragraph{One-dimensional characters.}
\[
\begin{array}{c|cc}
 & \chi_{1}\,(1) & \chi_{2}\,(1) \\\hline
T_{1} & 1 & 1 \\
T_{2} & 1 & -1 \\
T_{3} & 3 & 3 \\
T_{4} & 3 & -3 \\
T_{5} & 3 & 3 \\
T_{6} & 3 & -3 \\
T_{7} & 3 & -3 \\
T_{8} & 3 & 3 \\
T_{9} & 3 & 3 \\
T_{10} & 3 & -3 \\
T_{11} & 3 & -3 \\
T_{12} & 3 & 3 \\
T_{13} & 3 & -3 \\
T_{14} & 3 & 3 \\
T_{15} & 3 & 3 \\
T_{16} & 3 & -3 \\
T_{17} & 3 & -3 \\
T_{18} & 3 & 3 \\
T_{19} & 3 & -3 \\
T_{20} & 3 & 3 \\
T_{21} & 3 & 3 \\
T_{22} & 3 & -3 \\
\end{array}
\]

\subsubsection*{Block 1: $M_2(\C)$ on $\chi_{7}$ (degree 6)}
\noindent\textbf{Generators.}
\[
A = \rho(T_{2}) = \begin{pmatrix} 1 & 0 \\ 0 & -1 \end{pmatrix}, \qquad
B = \rho(T_{3}) = \begin{pmatrix} \tfrac{1}{2}\eta_3 & 1 \\ -\tfrac{3}{2}+\tfrac{1}{4}\eta_1-\tfrac{1}{2}\eta_2+\tfrac{1}{2}\eta_3 & \tfrac{1}{2}\eta_3 \end{pmatrix} .
\]
\textbf{Complete relations} (monomial basis $I,\;A,\;B,\;AB$ of $M_2(\C)$):
\begin{gather*}
A^2-I = 0 ,\\
B^2-\eta_3\,B+\bigl(3+\eta_2\bigr)I = 0 ,\\
BA = \eta_3\,A-AB ,\\
ABA = \eta_3I-B ,\\
ABB = \bigl(-3-\eta_2\bigr)A+\eta_3\,AB .
\end{gather*}
\textbf{All orbital images} ($\rho(T_i)$, valency $k_i$, word in the generators):
\begin{align*}
\rho(T_{1}) &= \begin{pmatrix} 1 & 0 \\ 0 & 1 \end{pmatrix} = I & (k_{1} &= 1) \\
\rho(T_{2}) &= \begin{pmatrix} 1 & 0 \\ 0 & -1 \end{pmatrix} = A & (k_{2} &= 1) \\
\rho(T_{3}) &= \begin{pmatrix} \tfrac{1}{2}\eta_3 & 1 \\ -\tfrac{3}{2}+\tfrac{1}{4}\eta_1-\tfrac{1}{2}\eta_2+\tfrac{1}{2}\eta_3 & \tfrac{1}{2}\eta_3 \end{pmatrix} = B & (k_{3} &= 3) \\
\rho(T_{4}) &= \begin{pmatrix} \tfrac{1}{2}\eta_3 & -1 \\ -\tfrac{3}{2}+\tfrac{1}{4}\eta_1-\tfrac{1}{2}\eta_2+\tfrac{1}{2}\eta_3 & -\tfrac{1}{2}\eta_3 \end{pmatrix} = \eta_3\,A-AB & (k_{4} &= 3) \\
\rho(T_{5}) &= \begin{pmatrix} \tfrac{1}{2}\eta_1 & 2+\eta_3 \\ \tfrac{1}{2}\eta_2-\tfrac{1}{4}\eta_3-\tfrac{1}{4}\eta_4 & \tfrac{1}{2}\eta_1 \end{pmatrix} & (k_{5} &= 3) \\
 &= \bigl(-3-\eta_2-2\eta_3\bigr)I+\bigl(2+\eta_3\bigr)B && \\
\rho(T_{6}) &= \begin{pmatrix} \tfrac{1}{2}\eta_1 & -2-\eta_3 \\ \tfrac{1}{2}\eta_2-\tfrac{1}{4}\eta_3-\tfrac{1}{4}\eta_4 & -\tfrac{1}{2}\eta_1 \end{pmatrix} & (k_{6} &= 3) \\
 &= \bigl(3+\eta_1+\eta_2+2\eta_3\bigr)A+\bigl(-2-\eta_3\bigr)AB && \\
\rho(T_{7}) &= \begin{pmatrix} \tfrac{1}{2}\eta_3 & 1 \\ \tfrac{3}{2}-\tfrac{1}{4}\eta_1+\tfrac{1}{2}\eta_2-\tfrac{1}{2}\eta_3 & -\tfrac{1}{2}\eta_3 \end{pmatrix} = AB & (k_{7} &= 3) \\
\rho(T_{8}) &= \begin{pmatrix} \tfrac{1}{2}\eta_3 & -1 \\ \tfrac{3}{2}-\tfrac{1}{4}\eta_1+\tfrac{1}{2}\eta_2-\tfrac{1}{2}\eta_3 & \tfrac{1}{2}\eta_3 \end{pmatrix} = \eta_3I-B & (k_{8} &= 3) \\
\rho(T_{9}) &= \begin{pmatrix} \tfrac{1}{2}\eta_8 & 3+\eta_1+\eta_2+2\eta_3 \\ -\tfrac{1}{2}\eta_1+\tfrac{1}{4}\eta_2+\tfrac{1}{4}\eta_8 & \tfrac{1}{2}\eta_8 \end{pmatrix} & (k_{9} &= 3) \\
 &= \bigl(-5-\eta_1-2\eta_2-4\eta_3\bigr)I+\bigl(3+\eta_1+\eta_2+2\eta_3\bigr)B && \\
\rho(T_{10}) &= \begin{pmatrix} \tfrac{1}{2}\eta_8 & -3-\eta_1-\eta_2-2\eta_3 \\ -\tfrac{1}{2}\eta_1+\tfrac{1}{4}\eta_2+\tfrac{1}{4}\eta_8 & -\tfrac{1}{2}\eta_8 \end{pmatrix} & (k_{10} &= 3) \\
 &= \bigl(4+\eta_2+3\eta_3-\eta_4\bigr)A+\bigl(-3-\eta_1-\eta_2-2\eta_3\bigr)AB && \\
\rho(T_{11}) &= \begin{pmatrix} \tfrac{1}{2}\eta_1 & 2+\eta_3 \\ -\tfrac{1}{2}\eta_2+\tfrac{1}{4}\eta_3+\tfrac{1}{4}\eta_4 & -\tfrac{1}{2}\eta_1 \end{pmatrix} & (k_{11} &= 3) \\
 &= \bigl(-3-\eta_2-2\eta_3\bigr)A+\bigl(2+\eta_3\bigr)AB && \\
\rho(T_{12}) &= \begin{pmatrix} \tfrac{1}{2}\eta_1 & -2-\eta_3 \\ -\tfrac{1}{2}\eta_2+\tfrac{1}{4}\eta_3+\tfrac{1}{4}\eta_4 & \tfrac{1}{2}\eta_1 \end{pmatrix} & (k_{12} &= 3) \\
 &= \bigl(3+\eta_1+\eta_2+2\eta_3\bigr)I+\bigl(-2-\eta_3\bigr)B && \\
\rho(T_{13}) &= \begin{pmatrix} \tfrac{1}{2}\eta_2 & 2+\eta_3-\eta_4 \\ -\tfrac{1}{4}\eta_4+\tfrac{1}{4}\eta_8 & -\tfrac{1}{2}\eta_2 \end{pmatrix} & (k_{13} &= 3) \\
 &= \bigl(-4-\eta_1-\eta_2-3\eta_3\bigr)A+\bigl(2+\eta_3-\eta_4\bigr)AB && \\
\rho(T_{14}) &= \begin{pmatrix} \tfrac{1}{2}\eta_2 & -2-\eta_3+\eta_4 \\ -\tfrac{1}{4}\eta_4+\tfrac{1}{4}\eta_8 & \tfrac{1}{2}\eta_2 \end{pmatrix} & (k_{14} &= 3) \\
 &= \bigl(4+\eta_1+2\eta_2+3\eta_3\bigr)I+\bigl(-2-\eta_3+\eta_4\bigr)B && \\
\rho(T_{15}) &= \begin{pmatrix} \tfrac{1}{2}\eta_2 & 2+\eta_3-\eta_4 \\ \tfrac{1}{4}\eta_4-\tfrac{1}{4}\eta_8 & \tfrac{1}{2}\eta_2 \end{pmatrix} & (k_{15} &= 3) \\
 &= \bigl(-4-\eta_1-\eta_2-3\eta_3\bigr)I+\bigl(2+\eta_3-\eta_4\bigr)B && \\
\rho(T_{16}) &= \begin{pmatrix} \tfrac{1}{2}\eta_2 & -2-\eta_3+\eta_4 \\ \tfrac{1}{4}\eta_4-\tfrac{1}{4}\eta_8 & -\tfrac{1}{2}\eta_2 \end{pmatrix} & (k_{16} &= 3) \\
 &= \bigl(4+\eta_1+2\eta_2+3\eta_3\bigr)A+\bigl(-2-\eta_3+\eta_4\bigr)AB && \\
\rho(T_{17}) &= \begin{pmatrix} \tfrac{1}{2}\eta_8 & 3+\eta_1+\eta_2+2\eta_3 \\ \tfrac{1}{2}\eta_1-\tfrac{1}{4}\eta_2-\tfrac{1}{4}\eta_8 & -\tfrac{1}{2}\eta_8 \end{pmatrix} & (k_{17} &= 3) \\
 &= \bigl(-5-\eta_1-2\eta_2-4\eta_3\bigr)A+\bigl(3+\eta_1+\eta_2+2\eta_3\bigr)AB && \\
\rho(T_{18}) &= \begin{pmatrix} \tfrac{1}{2}\eta_8 & -3-\eta_1-\eta_2-2\eta_3 \\ \tfrac{1}{2}\eta_1-\tfrac{1}{4}\eta_2-\tfrac{1}{4}\eta_8 & \tfrac{1}{2}\eta_8 \end{pmatrix} & (k_{18} &= 3) \\
 &= \bigl(4+\eta_2+3\eta_3-\eta_4\bigr)I+\bigl(-3-\eta_1-\eta_2-2\eta_3\bigr)B && \\
\rho(T_{19}) &= \begin{pmatrix} \tfrac{1}{2}\eta_4 & 5+\eta_1+2\eta_2+4\eta_3 \\ -\tfrac{1}{4}\eta_1+\tfrac{1}{4}\eta_2 & -\tfrac{1}{2}\eta_4 \end{pmatrix} & (k_{19} &= 3) \\
 &= \bigl(-10-2\eta_1-3\eta_2-7\eta_3+\eta_4\bigr)A+\bigl(5+\eta_1+2\eta_2+4\eta_3\bigr)AB && \\
\rho(T_{20}) &= \begin{pmatrix} \tfrac{1}{2}\eta_4 & -5-\eta_1-2\eta_2-4\eta_3 \\ -\tfrac{1}{4}\eta_1+\tfrac{1}{4}\eta_2 & \tfrac{1}{2}\eta_4 \end{pmatrix} & (k_{20} &= 3) \\
 &= \bigl(10+2\eta_1+3\eta_2+7\eta_3\bigr)I+\bigl(-5-\eta_1-2\eta_2-4\eta_3\bigr)B && \\
\rho(T_{21}) &= \begin{pmatrix} \tfrac{1}{2}\eta_4 & 5+\eta_1+2\eta_2+4\eta_3 \\ \tfrac{1}{4}\eta_1-\tfrac{1}{4}\eta_2 & \tfrac{1}{2}\eta_4 \end{pmatrix} & (k_{21} &= 3) \\
 &= \bigl(-10-2\eta_1-3\eta_2-7\eta_3+\eta_4\bigr)I+\bigl(5+\eta_1+2\eta_2+4\eta_3\bigr)B && \\
\rho(T_{22}) &= \begin{pmatrix} \tfrac{1}{2}\eta_4 & -5-\eta_1-2\eta_2-4\eta_3 \\ \tfrac{1}{4}\eta_1-\tfrac{1}{4}\eta_2 & -\tfrac{1}{2}\eta_4 \end{pmatrix} & (k_{22} &= 3) \\
 &= \bigl(10+2\eta_1+3\eta_2+7\eta_3\bigr)A+\bigl(-5-\eta_1-2\eta_2-4\eta_3\bigr)AB && .
\end{align*}

\subsubsection*{Block 2: $M_2(\C)$ on $\chi_{8}$ (degree 6)}
\noindent\textbf{Generators.}
\[
A = \rho(T_{2}) = \begin{pmatrix} 1 & 0 \\ 0 & -1 \end{pmatrix}, \qquad
B = \rho(T_{3}) = \begin{pmatrix} \tfrac{1}{2}\eta_1 & 1 \\ -\tfrac{3}{2}+\tfrac{1}{2}\eta_1+\tfrac{1}{4}\eta_2-\tfrac{1}{2}\eta_4 & \tfrac{1}{2}\eta_1 \end{pmatrix} .
\]
\textbf{Complete relations} (monomial basis $I,\;A,\;B,\;AB$ of $M_2(\C)$):
\begin{gather*}
A^2-I = 0 ,\\
B^2-\eta_1\,B+\bigl(3+\eta_4\bigr)I = 0 ,\\
BA = \eta_1\,A-AB ,\\
ABA = \eta_1I-B ,\\
ABB = \bigl(-3-\eta_4\bigr)A+\eta_1\,AB .
\end{gather*}
\textbf{All orbital images} ($\rho(T_i)$, valency $k_i$, word in the generators):
\begin{align*}
\rho(T_{1}) &= \begin{pmatrix} 1 & 0 \\ 0 & 1 \end{pmatrix} = I & (k_{1} &= 1) \\
\rho(T_{2}) &= \begin{pmatrix} 1 & 0 \\ 0 & -1 \end{pmatrix} = A & (k_{2} &= 1) \\
\rho(T_{3}) &= \begin{pmatrix} \tfrac{1}{2}\eta_1 & 1 \\ -\tfrac{3}{2}+\tfrac{1}{2}\eta_1+\tfrac{1}{4}\eta_2-\tfrac{1}{2}\eta_4 & \tfrac{1}{2}\eta_1 \end{pmatrix} = B & (k_{3} &= 3) \\
\rho(T_{4}) &= \begin{pmatrix} \tfrac{1}{2}\eta_1 & -1 \\ -\tfrac{3}{2}+\tfrac{1}{2}\eta_1+\tfrac{1}{4}\eta_2-\tfrac{1}{2}\eta_4 & -\tfrac{1}{2}\eta_1 \end{pmatrix} = \eta_1\,A-AB & (k_{4} &= 3) \\
\rho(T_{5}) &= \begin{pmatrix} \tfrac{1}{2}\eta_2 & 2+\eta_1 \\ -\tfrac{1}{4}\eta_1+\tfrac{1}{2}\eta_4-\tfrac{1}{4}\eta_8 & \tfrac{1}{2}\eta_2 \end{pmatrix} & (k_{5} &= 3) \\
 &= \bigl(-3-2\eta_1-\eta_4\bigr)I+\bigl(2+\eta_1\bigr)B && \\
\rho(T_{6}) &= \begin{pmatrix} \tfrac{1}{2}\eta_2 & -2-\eta_1 \\ -\tfrac{1}{4}\eta_1+\tfrac{1}{2}\eta_4-\tfrac{1}{4}\eta_8 & -\tfrac{1}{2}\eta_2 \end{pmatrix} & (k_{6} &= 3) \\
 &= \bigl(3+2\eta_1+\eta_2+\eta_4\bigr)A+\bigl(-2-\eta_1\bigr)AB && \\
\rho(T_{7}) &= \begin{pmatrix} \tfrac{1}{2}\eta_1 & 1 \\ \tfrac{3}{2}-\tfrac{1}{2}\eta_1-\tfrac{1}{4}\eta_2+\tfrac{1}{2}\eta_4 & -\tfrac{1}{2}\eta_1 \end{pmatrix} = AB & (k_{7} &= 3) \\
\rho(T_{8}) &= \begin{pmatrix} \tfrac{1}{2}\eta_1 & -1 \\ \tfrac{3}{2}-\tfrac{1}{2}\eta_1-\tfrac{1}{4}\eta_2+\tfrac{1}{2}\eta_4 & \tfrac{1}{2}\eta_1 \end{pmatrix} = \eta_1I-B & (k_{8} &= 3) \\
\rho(T_{9}) &= \begin{pmatrix} \tfrac{1}{2}\eta_3 & 3+2\eta_1+\eta_2+\eta_4 \\ -\tfrac{1}{2}\eta_2+\tfrac{1}{4}\eta_3+\tfrac{1}{4}\eta_4 & \tfrac{1}{2}\eta_3 \end{pmatrix} & (k_{9} &= 3) \\
 &= \bigl(-5-4\eta_1-\eta_2-2\eta_4\bigr)I+\bigl(3+2\eta_1+\eta_2+\eta_4\bigr)B && \\
\rho(T_{10}) &= \begin{pmatrix} \tfrac{1}{2}\eta_3 & -3-2\eta_1-\eta_2-\eta_4 \\ -\tfrac{1}{2}\eta_2+\tfrac{1}{4}\eta_3+\tfrac{1}{4}\eta_4 & -\tfrac{1}{2}\eta_3 \end{pmatrix} & (k_{10} &= 3) \\
 &= \bigl(4+3\eta_1+\eta_4-\eta_8\bigr)A+\bigl(-3-2\eta_1-\eta_2-\eta_4\bigr)AB && \\
\rho(T_{11}) &= \begin{pmatrix} \tfrac{1}{2}\eta_2 & 2+\eta_1 \\ \tfrac{1}{4}\eta_1-\tfrac{1}{2}\eta_4+\tfrac{1}{4}\eta_8 & -\tfrac{1}{2}\eta_2 \end{pmatrix} & (k_{11} &= 3) \\
 &= \bigl(-3-2\eta_1-\eta_4\bigr)A+\bigl(2+\eta_1\bigr)AB && \\
\rho(T_{12}) &= \begin{pmatrix} \tfrac{1}{2}\eta_2 & -2-\eta_1 \\ \tfrac{1}{4}\eta_1-\tfrac{1}{2}\eta_4+\tfrac{1}{4}\eta_8 & \tfrac{1}{2}\eta_2 \end{pmatrix} & (k_{12} &= 3) \\
 &= \bigl(3+2\eta_1+\eta_2+\eta_4\bigr)I+\bigl(-2-\eta_1\bigr)B && \\
\rho(T_{13}) &= \begin{pmatrix} \tfrac{1}{2}\eta_4 & 2+\eta_1-\eta_8 \\ \tfrac{1}{4}\eta_3-\tfrac{1}{4}\eta_8 & -\tfrac{1}{2}\eta_4 \end{pmatrix} & (k_{13} &= 3) \\
 &= \bigl(-4-3\eta_1-\eta_2-\eta_4\bigr)A+\bigl(2+\eta_1-\eta_8\bigr)AB && \\
\rho(T_{14}) &= \begin{pmatrix} \tfrac{1}{2}\eta_4 & -2-\eta_1+\eta_8 \\ \tfrac{1}{4}\eta_3-\tfrac{1}{4}\eta_8 & \tfrac{1}{2}\eta_4 \end{pmatrix} & (k_{14} &= 3) \\
 &= \bigl(4+3\eta_1+\eta_2+2\eta_4\bigr)I+\bigl(-2-\eta_1+\eta_8\bigr)B && \\
\rho(T_{15}) &= \begin{pmatrix} \tfrac{1}{2}\eta_4 & 2+\eta_1-\eta_8 \\ -\tfrac{1}{4}\eta_3+\tfrac{1}{4}\eta_8 & \tfrac{1}{2}\eta_4 \end{pmatrix} & (k_{15} &= 3) \\
 &= \bigl(-4-3\eta_1-\eta_2-\eta_4\bigr)I+\bigl(2+\eta_1-\eta_8\bigr)B && \\
\rho(T_{16}) &= \begin{pmatrix} \tfrac{1}{2}\eta_4 & -2-\eta_1+\eta_8 \\ -\tfrac{1}{4}\eta_3+\tfrac{1}{4}\eta_8 & -\tfrac{1}{2}\eta_4 \end{pmatrix} & (k_{16} &= 3) \\
 &= \bigl(4+3\eta_1+\eta_2+2\eta_4\bigr)A+\bigl(-2-\eta_1+\eta_8\bigr)AB && \\
\rho(T_{17}) &= \begin{pmatrix} \tfrac{1}{2}\eta_3 & 3+2\eta_1+\eta_2+\eta_4 \\ \tfrac{1}{2}\eta_2-\tfrac{1}{4}\eta_3-\tfrac{1}{4}\eta_4 & -\tfrac{1}{2}\eta_3 \end{pmatrix} & (k_{17} &= 3) \\
 &= \bigl(-5-4\eta_1-\eta_2-2\eta_4\bigr)A+\bigl(3+2\eta_1+\eta_2+\eta_4\bigr)AB && \\
\rho(T_{18}) &= \begin{pmatrix} \tfrac{1}{2}\eta_3 & -3-2\eta_1-\eta_2-\eta_4 \\ \tfrac{1}{2}\eta_2-\tfrac{1}{4}\eta_3-\tfrac{1}{4}\eta_4 & \tfrac{1}{2}\eta_3 \end{pmatrix} & (k_{18} &= 3) \\
 &= \bigl(4+3\eta_1+\eta_4-\eta_8\bigr)I+\bigl(-3-2\eta_1-\eta_2-\eta_4\bigr)B && \\
\rho(T_{19}) &= \begin{pmatrix} \tfrac{1}{2}\eta_8 & 5+4\eta_1+\eta_2+2\eta_4 \\ -\tfrac{1}{4}\eta_2+\tfrac{1}{4}\eta_4 & -\tfrac{1}{2}\eta_8 \end{pmatrix} & (k_{19} &= 3) \\
 &= \bigl(-11-8\eta_1-3\eta_2-\eta_3-4\eta_4\bigr)A+\bigl(5+4\eta_1+\eta_2+2\eta_4\bigr)AB && \\
\rho(T_{20}) &= \begin{pmatrix} \tfrac{1}{2}\eta_8 & -5-4\eta_1-\eta_2-2\eta_4 \\ -\tfrac{1}{4}\eta_2+\tfrac{1}{4}\eta_4 & \tfrac{1}{2}\eta_8 \end{pmatrix} & (k_{20} &= 3) \\
 &= \bigl(10+7\eta_1+2\eta_2+3\eta_4\bigr)I+\bigl(-5-4\eta_1-\eta_2-2\eta_4\bigr)B && \\
\rho(T_{21}) &= \begin{pmatrix} \tfrac{1}{2}\eta_8 & 5+4\eta_1+\eta_2+2\eta_4 \\ \tfrac{1}{4}\eta_2-\tfrac{1}{4}\eta_4 & \tfrac{1}{2}\eta_8 \end{pmatrix} & (k_{21} &= 3) \\
 &= \bigl(-11-8\eta_1-3\eta_2-\eta_3-4\eta_4\bigr)I+\bigl(5+4\eta_1+\eta_2+2\eta_4\bigr)B && \\
\rho(T_{22}) &= \begin{pmatrix} \tfrac{1}{2}\eta_8 & -5-4\eta_1-\eta_2-2\eta_4 \\ \tfrac{1}{4}\eta_2-\tfrac{1}{4}\eta_4 & -\tfrac{1}{2}\eta_8 \end{pmatrix} & (k_{22} &= 3) \\
 &= \bigl(10+7\eta_1+2\eta_2+3\eta_4\bigr)A+\bigl(-5-4\eta_1-\eta_2-2\eta_4\bigr)AB && .
\end{align*}

\subsubsection*{Block 3: $M_2(\C)$ on $\chi_{9}$ (degree 6)}
\noindent\textbf{Generators.}
\[
A = \rho(T_{2}) = \begin{pmatrix} 1 & 0 \\ 0 & -1 \end{pmatrix}, \qquad
B = \rho(T_{3}) = \begin{pmatrix} \tfrac{1}{2}\eta_8 & 1 \\ -\tfrac{3}{2}-\tfrac{1}{2}\eta_1+\tfrac{1}{4}\eta_3+\tfrac{1}{2}\eta_8 & \tfrac{1}{2}\eta_8 \end{pmatrix} .
\]
\textbf{Complete relations} (monomial basis $I,\;A,\;B,\;AB$ of $M_2(\C)$):
\begin{gather*}
A^2-I = 0 ,\\
B^2-\eta_8\,B+\bigl(3+\eta_1\bigr)I = 0 ,\\
BA = \eta_8\,A-AB ,\\
ABA = \eta_8I-B ,\\
ABB = \bigl(-3-\eta_1\bigr)A+\eta_8\,AB .
\end{gather*}
\textbf{All orbital images} ($\rho(T_i)$, valency $k_i$, word in the generators):
\begin{align*}
\rho(T_{1}) &= \begin{pmatrix} 1 & 0 \\ 0 & 1 \end{pmatrix} = I & (k_{1} &= 1) \\
\rho(T_{2}) &= \begin{pmatrix} 1 & 0 \\ 0 & -1 \end{pmatrix} = A & (k_{2} &= 1) \\
\rho(T_{3}) &= \begin{pmatrix} \tfrac{1}{2}\eta_8 & 1 \\ -\tfrac{3}{2}-\tfrac{1}{2}\eta_1+\tfrac{1}{4}\eta_3+\tfrac{1}{2}\eta_8 & \tfrac{1}{2}\eta_8 \end{pmatrix} = B & (k_{3} &= 3) \\
\rho(T_{4}) &= \begin{pmatrix} \tfrac{1}{2}\eta_8 & -1 \\ -\tfrac{3}{2}-\tfrac{1}{2}\eta_1+\tfrac{1}{4}\eta_3+\tfrac{1}{2}\eta_8 & -\tfrac{1}{2}\eta_8 \end{pmatrix} = \eta_8\,A-AB & (k_{4} &= 3) \\
\rho(T_{5}) &= \begin{pmatrix} \tfrac{1}{2}\eta_3 & 2+\eta_8 \\ \tfrac{1}{2}\eta_1-\tfrac{1}{4}\eta_2-\tfrac{1}{4}\eta_8 & \tfrac{1}{2}\eta_3 \end{pmatrix} & (k_{5} &= 3) \\
 &= \bigl(-3-\eta_1-2\eta_8\bigr)I+\bigl(2+\eta_8\bigr)B && \\
\rho(T_{6}) &= \begin{pmatrix} \tfrac{1}{2}\eta_3 & -2-\eta_8 \\ \tfrac{1}{2}\eta_1-\tfrac{1}{4}\eta_2-\tfrac{1}{4}\eta_8 & -\tfrac{1}{2}\eta_3 \end{pmatrix} & (k_{6} &= 3) \\
 &= \bigl(3+\eta_1+\eta_3+2\eta_8\bigr)A+\bigl(-2-\eta_8\bigr)AB && \\
\rho(T_{7}) &= \begin{pmatrix} \tfrac{1}{2}\eta_8 & 1 \\ \tfrac{3}{2}+\tfrac{1}{2}\eta_1-\tfrac{1}{4}\eta_3-\tfrac{1}{2}\eta_8 & -\tfrac{1}{2}\eta_8 \end{pmatrix} = AB & (k_{7} &= 3) \\
\rho(T_{8}) &= \begin{pmatrix} \tfrac{1}{2}\eta_8 & -1 \\ \tfrac{3}{2}+\tfrac{1}{2}\eta_1-\tfrac{1}{4}\eta_3-\tfrac{1}{2}\eta_8 & \tfrac{1}{2}\eta_8 \end{pmatrix} = \eta_8I-B & (k_{8} &= 3) \\
\rho(T_{9}) &= \begin{pmatrix} \tfrac{1}{2}\eta_4 & 3+\eta_1+\eta_3+2\eta_8 \\ \tfrac{1}{4}\eta_1-\tfrac{1}{2}\eta_3+\tfrac{1}{4}\eta_4 & \tfrac{1}{2}\eta_4 \end{pmatrix} & (k_{9} &= 3) \\
 &= \bigl(-5-2\eta_1-\eta_3-4\eta_8\bigr)I+\bigl(3+\eta_1+\eta_3+2\eta_8\bigr)B && \\
\rho(T_{10}) &= \begin{pmatrix} \tfrac{1}{2}\eta_4 & -3-\eta_1-\eta_3-2\eta_8 \\ \tfrac{1}{4}\eta_1-\tfrac{1}{2}\eta_3+\tfrac{1}{4}\eta_4 & -\tfrac{1}{2}\eta_4 \end{pmatrix} & (k_{10} &= 3) \\
 &= \bigl(4+\eta_1-\eta_2+3\eta_8\bigr)A+\bigl(-3-\eta_1-\eta_3-2\eta_8\bigr)AB && \\
\rho(T_{11}) &= \begin{pmatrix} \tfrac{1}{2}\eta_3 & 2+\eta_8 \\ -\tfrac{1}{2}\eta_1+\tfrac{1}{4}\eta_2+\tfrac{1}{4}\eta_8 & -\tfrac{1}{2}\eta_3 \end{pmatrix} & (k_{11} &= 3) \\
 &= \bigl(-3-\eta_1-2\eta_8\bigr)A+\bigl(2+\eta_8\bigr)AB && \\
\rho(T_{12}) &= \begin{pmatrix} \tfrac{1}{2}\eta_3 & -2-\eta_8 \\ -\tfrac{1}{2}\eta_1+\tfrac{1}{4}\eta_2+\tfrac{1}{4}\eta_8 & \tfrac{1}{2}\eta_3 \end{pmatrix} & (k_{12} &= 3) \\
 &= \bigl(3+\eta_1+\eta_3+2\eta_8\bigr)I+\bigl(-2-\eta_8\bigr)B && \\
\rho(T_{13}) &= \begin{pmatrix} \tfrac{1}{2}\eta_1 & 2-\eta_2+\eta_8 \\ -\tfrac{1}{4}\eta_2+\tfrac{1}{4}\eta_4 & -\tfrac{1}{2}\eta_1 \end{pmatrix} & (k_{13} &= 3) \\
 &= \bigl(-4-\eta_1-\eta_3-3\eta_8\bigr)A+\bigl(2-\eta_2+\eta_8\bigr)AB && \\
\rho(T_{14}) &= \begin{pmatrix} \tfrac{1}{2}\eta_1 & -2+\eta_2-\eta_8 \\ -\tfrac{1}{4}\eta_2+\tfrac{1}{4}\eta_4 & \tfrac{1}{2}\eta_1 \end{pmatrix} & (k_{14} &= 3) \\
 &= \bigl(4+2\eta_1+\eta_3+3\eta_8\bigr)I+\bigl(-2+\eta_2-\eta_8\bigr)B && \\
\rho(T_{15}) &= \begin{pmatrix} \tfrac{1}{2}\eta_1 & 2-\eta_2+\eta_8 \\ \tfrac{1}{4}\eta_2-\tfrac{1}{4}\eta_4 & \tfrac{1}{2}\eta_1 \end{pmatrix} & (k_{15} &= 3) \\
 &= \bigl(-4-\eta_1-\eta_3-3\eta_8\bigr)I+\bigl(2-\eta_2+\eta_8\bigr)B && \\
\rho(T_{16}) &= \begin{pmatrix} \tfrac{1}{2}\eta_1 & -2+\eta_2-\eta_8 \\ \tfrac{1}{4}\eta_2-\tfrac{1}{4}\eta_4 & -\tfrac{1}{2}\eta_1 \end{pmatrix} & (k_{16} &= 3) \\
 &= \bigl(4+2\eta_1+\eta_3+3\eta_8\bigr)A+\bigl(-2+\eta_2-\eta_8\bigr)AB && \\
\rho(T_{17}) &= \begin{pmatrix} \tfrac{1}{2}\eta_4 & 3+\eta_1+\eta_3+2\eta_8 \\ -\tfrac{1}{4}\eta_1+\tfrac{1}{2}\eta_3-\tfrac{1}{4}\eta_4 & -\tfrac{1}{2}\eta_4 \end{pmatrix} & (k_{17} &= 3) \\
 &= \bigl(-5-2\eta_1-\eta_3-4\eta_8\bigr)A+\bigl(3+\eta_1+\eta_3+2\eta_8\bigr)AB && \\
\rho(T_{18}) &= \begin{pmatrix} \tfrac{1}{2}\eta_4 & -3-\eta_1-\eta_3-2\eta_8 \\ -\tfrac{1}{4}\eta_1+\tfrac{1}{2}\eta_3-\tfrac{1}{4}\eta_4 & \tfrac{1}{2}\eta_4 \end{pmatrix} & (k_{18} &= 3) \\
 &= \bigl(4+\eta_1-\eta_2+3\eta_8\bigr)I+\bigl(-3-\eta_1-\eta_3-2\eta_8\bigr)B && \\
\rho(T_{19}) &= \begin{pmatrix} \tfrac{1}{2}\eta_2 & 5+2\eta_1+\eta_3+4\eta_8 \\ \tfrac{1}{4}\eta_1-\tfrac{1}{4}\eta_3 & -\tfrac{1}{2}\eta_2 \end{pmatrix} & (k_{19} &= 3) \\
 &= \bigl(-3+4\eta_1+8\eta_2+5\eta_3+7\eta_4\bigr)A+\bigl(5+2\eta_1+\eta_3+4\eta_8\bigr)AB && \\
\rho(T_{20}) &= \begin{pmatrix} \tfrac{1}{2}\eta_2 & -5-2\eta_1-\eta_3-4\eta_8 \\ \tfrac{1}{4}\eta_1-\tfrac{1}{4}\eta_3 & \tfrac{1}{2}\eta_2 \end{pmatrix} & (k_{20} &= 3) \\
 &= \bigl(10+3\eta_1+2\eta_3+7\eta_8\bigr)I+\bigl(-5-2\eta_1-\eta_3-4\eta_8\bigr)B && \\
\rho(T_{21}) &= \begin{pmatrix} \tfrac{1}{2}\eta_2 & 5+2\eta_1+\eta_3+4\eta_8 \\ -\tfrac{1}{4}\eta_1+\tfrac{1}{4}\eta_3 & \tfrac{1}{2}\eta_2 \end{pmatrix} & (k_{21} &= 3) \\
 &= \bigl(-3+4\eta_1+8\eta_2+5\eta_3+7\eta_4\bigr)I+\bigl(5+2\eta_1+\eta_3+4\eta_8\bigr)B && \\
\rho(T_{22}) &= \begin{pmatrix} \tfrac{1}{2}\eta_2 & -5-2\eta_1-\eta_3-4\eta_8 \\ -\tfrac{1}{4}\eta_1+\tfrac{1}{4}\eta_3 & -\tfrac{1}{2}\eta_2 \end{pmatrix} & (k_{22} &= 3) \\
 &= \bigl(10+3\eta_1+2\eta_3+7\eta_8\bigr)A+\bigl(-5-2\eta_1-\eta_3-4\eta_8\bigr)AB && .
\end{align*}

\subsubsection*{Block 4: $M_2(\C)$ on $\chi_{10}$ (degree 6)}
\noindent\textbf{Generators.}
\[
A = \rho(T_{2}) = \begin{pmatrix} 1 & 0 \\ 0 & -1 \end{pmatrix}, \qquad
B = \rho(T_{3}) = \begin{pmatrix} \tfrac{1}{2}\eta_2 & 1 \\ -\tfrac{3}{2}+\tfrac{1}{2}\eta_2+\tfrac{1}{4}\eta_4-\tfrac{1}{2}\eta_8 & \tfrac{1}{2}\eta_2 \end{pmatrix} .
\]
\textbf{Complete relations} (monomial basis $I,\;A,\;B,\;AB$ of $M_2(\C)$):
\begin{gather*}
A^2-I = 0 ,\\
B^2-\eta_2\,B+\bigl(3+\eta_8\bigr)I = 0 ,\\
BA = \eta_2\,A-AB ,\\
ABA = \eta_2I-B ,\\
ABB = \bigl(-3-\eta_8\bigr)A+\eta_2\,AB .
\end{gather*}
\textbf{All orbital images} ($\rho(T_i)$, valency $k_i$, word in the generators):
\begin{align*}
\rho(T_{1}) &= \begin{pmatrix} 1 & 0 \\ 0 & 1 \end{pmatrix} = I & (k_{1} &= 1) \\
\rho(T_{2}) &= \begin{pmatrix} 1 & 0 \\ 0 & -1 \end{pmatrix} = A & (k_{2} &= 1) \\
\rho(T_{3}) &= \begin{pmatrix} \tfrac{1}{2}\eta_2 & 1 \\ -\tfrac{3}{2}+\tfrac{1}{2}\eta_2+\tfrac{1}{4}\eta_4-\tfrac{1}{2}\eta_8 & \tfrac{1}{2}\eta_2 \end{pmatrix} = B & (k_{3} &= 3) \\
\rho(T_{4}) &= \begin{pmatrix} \tfrac{1}{2}\eta_2 & -1 \\ -\tfrac{3}{2}+\tfrac{1}{2}\eta_2+\tfrac{1}{4}\eta_4-\tfrac{1}{2}\eta_8 & -\tfrac{1}{2}\eta_2 \end{pmatrix} = \eta_2\,A-AB & (k_{4} &= 3) \\
\rho(T_{5}) &= \begin{pmatrix} \tfrac{1}{2}\eta_4 & 2+\eta_2 \\ -\tfrac{1}{4}\eta_2-\tfrac{1}{4}\eta_3+\tfrac{1}{2}\eta_8 & \tfrac{1}{2}\eta_4 \end{pmatrix} & (k_{5} &= 3) \\
 &= \bigl(-3-2\eta_2-\eta_8\bigr)I+\bigl(2+\eta_2\bigr)B && \\
\rho(T_{6}) &= \begin{pmatrix} \tfrac{1}{2}\eta_4 & -2-\eta_2 \\ -\tfrac{1}{4}\eta_2-\tfrac{1}{4}\eta_3+\tfrac{1}{2}\eta_8 & -\tfrac{1}{2}\eta_4 \end{pmatrix} & (k_{6} &= 3) \\
 &= \bigl(2-\eta_1+\eta_2-\eta_3\bigr)A+\bigl(-2-\eta_2\bigr)AB && \\
\rho(T_{7}) &= \begin{pmatrix} \tfrac{1}{2}\eta_2 & 1 \\ \tfrac{3}{2}-\tfrac{1}{2}\eta_2-\tfrac{1}{4}\eta_4+\tfrac{1}{2}\eta_8 & -\tfrac{1}{2}\eta_2 \end{pmatrix} = AB & (k_{7} &= 3) \\
\rho(T_{8}) &= \begin{pmatrix} \tfrac{1}{2}\eta_2 & -1 \\ \tfrac{3}{2}-\tfrac{1}{2}\eta_2-\tfrac{1}{4}\eta_4+\tfrac{1}{2}\eta_8 & \tfrac{1}{2}\eta_2 \end{pmatrix} = \eta_2I-B & (k_{8} &= 3) \\
\rho(T_{9}) &= \begin{pmatrix} \tfrac{1}{2}\eta_1 & 2-\eta_1+\eta_2-\eta_3 \\ \tfrac{1}{4}\eta_1-\tfrac{1}{2}\eta_4+\tfrac{1}{4}\eta_8 & \tfrac{1}{2}\eta_1 \end{pmatrix} & (k_{9} &= 3) \\
 &= \bigl(-5-4\eta_2-\eta_4-2\eta_8\bigr)I+\bigl(2-\eta_1+\eta_2-\eta_3\bigr)B && \\
\rho(T_{10}) &= \begin{pmatrix} \tfrac{1}{2}\eta_1 & -2+\eta_1-\eta_2+\eta_3 \\ \tfrac{1}{4}\eta_1-\tfrac{1}{2}\eta_4+\tfrac{1}{4}\eta_8 & -\tfrac{1}{2}\eta_1 \end{pmatrix} & (k_{10} &= 3) \\
 &= \bigl(4+3\eta_2-\eta_3+\eta_8\bigr)A+\bigl(-2+\eta_1-\eta_2+\eta_3\bigr)AB && \\
\rho(T_{11}) &= \begin{pmatrix} \tfrac{1}{2}\eta_4 & 2+\eta_2 \\ \tfrac{1}{4}\eta_2+\tfrac{1}{4}\eta_3-\tfrac{1}{2}\eta_8 & -\tfrac{1}{2}\eta_4 \end{pmatrix} & (k_{11} &= 3) \\
 &= \bigl(-3-2\eta_2-\eta_8\bigr)A+\bigl(2+\eta_2\bigr)AB && \\
\rho(T_{12}) &= \begin{pmatrix} \tfrac{1}{2}\eta_4 & -2-\eta_2 \\ \tfrac{1}{4}\eta_2+\tfrac{1}{4}\eta_3-\tfrac{1}{2}\eta_8 & \tfrac{1}{2}\eta_4 \end{pmatrix} & (k_{12} &= 3) \\
 &= \bigl(2-\eta_1+\eta_2-\eta_3\bigr)I+\bigl(-2-\eta_2\bigr)B && \\
\rho(T_{13}) &= \begin{pmatrix} \tfrac{1}{2}\eta_8 & 2+\eta_2-\eta_3 \\ \tfrac{1}{4}\eta_1-\tfrac{1}{4}\eta_3 & -\tfrac{1}{2}\eta_8 \end{pmatrix} & (k_{13} &= 3) \\
 &= \bigl(-3+\eta_1-2\eta_2+\eta_3\bigr)A+\bigl(2+\eta_2-\eta_3\bigr)AB && \\
\rho(T_{14}) &= \begin{pmatrix} \tfrac{1}{2}\eta_8 & -2-\eta_2+\eta_3 \\ \tfrac{1}{4}\eta_1-\tfrac{1}{4}\eta_3 & \tfrac{1}{2}\eta_8 \end{pmatrix} & (k_{14} &= 3) \\
 &= \bigl(4+3\eta_2+\eta_4+2\eta_8\bigr)I+\bigl(-2-\eta_2+\eta_3\bigr)B && \\
\rho(T_{15}) &= \begin{pmatrix} \tfrac{1}{2}\eta_8 & 2+\eta_2-\eta_3 \\ -\tfrac{1}{4}\eta_1+\tfrac{1}{4}\eta_3 & \tfrac{1}{2}\eta_8 \end{pmatrix} & (k_{15} &= 3) \\
 &= \bigl(-3+\eta_1-2\eta_2+\eta_3\bigr)I+\bigl(2+\eta_2-\eta_3\bigr)B && \\
\rho(T_{16}) &= \begin{pmatrix} \tfrac{1}{2}\eta_8 & -2-\eta_2+\eta_3 \\ -\tfrac{1}{4}\eta_1+\tfrac{1}{4}\eta_3 & -\tfrac{1}{2}\eta_8 \end{pmatrix} & (k_{16} &= 3) \\
 &= \bigl(4+3\eta_2+\eta_4+2\eta_8\bigr)A+\bigl(-2-\eta_2+\eta_3\bigr)AB && \\
\rho(T_{17}) &= \begin{pmatrix} \tfrac{1}{2}\eta_1 & 2-\eta_1+\eta_2-\eta_3 \\ -\tfrac{1}{4}\eta_1+\tfrac{1}{2}\eta_4-\tfrac{1}{4}\eta_8 & -\tfrac{1}{2}\eta_1 \end{pmatrix} & (k_{17} &= 3) \\
 &= \bigl(-5-4\eta_2-\eta_4-2\eta_8\bigr)A+\bigl(2-\eta_1+\eta_2-\eta_3\bigr)AB && \\
\rho(T_{18}) &= \begin{pmatrix} \tfrac{1}{2}\eta_1 & -2+\eta_1-\eta_2+\eta_3 \\ -\tfrac{1}{4}\eta_1+\tfrac{1}{2}\eta_4-\tfrac{1}{4}\eta_8 & \tfrac{1}{2}\eta_1 \end{pmatrix} & (k_{18} &= 3) \\
 &= \bigl(4+3\eta_2-\eta_3+\eta_8\bigr)I+\bigl(-2+\eta_1-\eta_2+\eta_3\bigr)B && \\
\rho(T_{19}) &= \begin{pmatrix} \tfrac{1}{2}\eta_3 & 5+4\eta_2+\eta_4+2\eta_8 \\ -\tfrac{1}{4}\eta_4+\tfrac{1}{4}\eta_8 & -\tfrac{1}{2}\eta_3 \end{pmatrix} & (k_{19} &= 3) \\
 &= \bigl(-7+3\eta_1-4\eta_2+4\eta_3+\eta_4\bigr)A+\bigl(5+4\eta_2+\eta_4+2\eta_8\bigr)AB && \\
\rho(T_{20}) &= \begin{pmatrix} \tfrac{1}{2}\eta_3 & -5-4\eta_2-\eta_4-2\eta_8 \\ -\tfrac{1}{4}\eta_4+\tfrac{1}{4}\eta_8 & \tfrac{1}{2}\eta_3 \end{pmatrix} & (k_{20} &= 3) \\
 &= \bigl(10+7\eta_2+2\eta_4+3\eta_8\bigr)I+\bigl(-5-4\eta_2-\eta_4-2\eta_8\bigr)B && \\
\rho(T_{21}) &= \begin{pmatrix} \tfrac{1}{2}\eta_3 & 5+4\eta_2+\eta_4+2\eta_8 \\ \tfrac{1}{4}\eta_4-\tfrac{1}{4}\eta_8 & \tfrac{1}{2}\eta_3 \end{pmatrix} & (k_{21} &= 3) \\
 &= \bigl(-7+3\eta_1-4\eta_2+4\eta_3+\eta_4\bigr)I+\bigl(5+4\eta_2+\eta_4+2\eta_8\bigr)B && \\
\rho(T_{22}) &= \begin{pmatrix} \tfrac{1}{2}\eta_3 & -5-4\eta_2-\eta_4-2\eta_8 \\ \tfrac{1}{4}\eta_4-\tfrac{1}{4}\eta_8 & -\tfrac{1}{2}\eta_3 \end{pmatrix} & (k_{22} &= 3) \\
 &= \bigl(10+7\eta_2+2\eta_4+3\eta_8\bigr)A+\bigl(-5-4\eta_2-\eta_4-2\eta_8\bigr)AB && .
\end{align*}

\subsubsection*{Block 5: $M_2(\C)$ on $\chi_{11}$ (degree 6)}
\noindent\textbf{Generators.}
\[
A = \rho(T_{2}) = \begin{pmatrix} 1 & 0 \\ 0 & -1 \end{pmatrix}, \qquad
B = \rho(T_{3}) = \begin{pmatrix} \tfrac{1}{2}\eta_4 & 1 \\ -\tfrac{3}{2}-\tfrac{1}{2}\eta_3+\tfrac{1}{2}\eta_4+\tfrac{1}{4}\eta_8 & \tfrac{1}{2}\eta_4 \end{pmatrix} .
\]
\textbf{Complete relations} (monomial basis $I,\;A,\;B,\;AB$ of $M_2(\C)$):
\begin{gather*}
A^2-I = 0 ,\\
B^2-\eta_4\,B+\bigl(3+\eta_3\bigr)I = 0 ,\\
BA = \eta_4\,A-AB ,\\
ABA = \eta_4I-B ,\\
ABB = \bigl(-3-\eta_3\bigr)A+\eta_4\,AB .
\end{gather*}
\textbf{All orbital images} ($\rho(T_i)$, valency $k_i$, word in the generators):
\begin{align*}
\rho(T_{1}) &= \begin{pmatrix} 1 & 0 \\ 0 & 1 \end{pmatrix} = I & (k_{1} &= 1) \\
\rho(T_{2}) &= \begin{pmatrix} 1 & 0 \\ 0 & -1 \end{pmatrix} = A & (k_{2} &= 1) \\
\rho(T_{3}) &= \begin{pmatrix} \tfrac{1}{2}\eta_4 & 1 \\ -\tfrac{3}{2}-\tfrac{1}{2}\eta_3+\tfrac{1}{2}\eta_4+\tfrac{1}{4}\eta_8 & \tfrac{1}{2}\eta_4 \end{pmatrix} = B & (k_{3} &= 3) \\
\rho(T_{4}) &= \begin{pmatrix} \tfrac{1}{2}\eta_4 & -1 \\ -\tfrac{3}{2}-\tfrac{1}{2}\eta_3+\tfrac{1}{2}\eta_4+\tfrac{1}{4}\eta_8 & -\tfrac{1}{2}\eta_4 \end{pmatrix} = \eta_4\,A-AB & (k_{4} &= 3) \\
\rho(T_{5}) &= \begin{pmatrix} \tfrac{1}{2}\eta_8 & 2+\eta_4 \\ -\tfrac{1}{4}\eta_1+\tfrac{1}{2}\eta_3-\tfrac{1}{4}\eta_4 & \tfrac{1}{2}\eta_8 \end{pmatrix} & (k_{5} &= 3) \\
 &= \bigl(-3-\eta_3-2\eta_4\bigr)I+\bigl(2+\eta_4\bigr)B && \\
\rho(T_{6}) &= \begin{pmatrix} \tfrac{1}{2}\eta_8 & -2-\eta_4 \\ -\tfrac{1}{4}\eta_1+\tfrac{1}{2}\eta_3-\tfrac{1}{4}\eta_4 & -\tfrac{1}{2}\eta_8 \end{pmatrix} & (k_{6} &= 3) \\
 &= \bigl(2-\eta_1-\eta_2+\eta_4\bigr)A+\bigl(-2-\eta_4\bigr)AB && \\
\rho(T_{7}) &= \begin{pmatrix} \tfrac{1}{2}\eta_4 & 1 \\ \tfrac{3}{2}+\tfrac{1}{2}\eta_3-\tfrac{1}{2}\eta_4-\tfrac{1}{4}\eta_8 & -\tfrac{1}{2}\eta_4 \end{pmatrix} = AB & (k_{7} &= 3) \\
\rho(T_{8}) &= \begin{pmatrix} \tfrac{1}{2}\eta_4 & -1 \\ \tfrac{3}{2}+\tfrac{1}{2}\eta_3-\tfrac{1}{2}\eta_4-\tfrac{1}{4}\eta_8 & \tfrac{1}{2}\eta_4 \end{pmatrix} = \eta_4I-B & (k_{8} &= 3) \\
\rho(T_{9}) &= \begin{pmatrix} \tfrac{1}{2}\eta_2 & 2-\eta_1-\eta_2+\eta_4 \\ \tfrac{1}{4}\eta_2+\tfrac{1}{4}\eta_3-\tfrac{1}{2}\eta_8 & \tfrac{1}{2}\eta_2 \end{pmatrix} & (k_{9} &= 3) \\
 &= \bigl(-5-2\eta_3-4\eta_4-\eta_8\bigr)I+\bigl(2-\eta_1-\eta_2+\eta_4\bigr)B && \\
\rho(T_{10}) &= \begin{pmatrix} \tfrac{1}{2}\eta_2 & -2+\eta_1+\eta_2-\eta_4 \\ \tfrac{1}{4}\eta_2+\tfrac{1}{4}\eta_3-\tfrac{1}{2}\eta_8 & -\tfrac{1}{2}\eta_2 \end{pmatrix} & (k_{10} &= 3) \\
 &= \bigl(4-\eta_1+\eta_3+3\eta_4\bigr)A+\bigl(-2+\eta_1+\eta_2-\eta_4\bigr)AB && \\
\rho(T_{11}) &= \begin{pmatrix} \tfrac{1}{2}\eta_8 & 2+\eta_4 \\ \tfrac{1}{4}\eta_1-\tfrac{1}{2}\eta_3+\tfrac{1}{4}\eta_4 & -\tfrac{1}{2}\eta_8 \end{pmatrix} & (k_{11} &= 3) \\
 &= \bigl(-3-\eta_3-2\eta_4\bigr)A+\bigl(2+\eta_4\bigr)AB && \\
\rho(T_{12}) &= \begin{pmatrix} \tfrac{1}{2}\eta_8 & -2-\eta_4 \\ \tfrac{1}{4}\eta_1-\tfrac{1}{2}\eta_3+\tfrac{1}{4}\eta_4 & \tfrac{1}{2}\eta_8 \end{pmatrix} & (k_{12} &= 3) \\
 &= \bigl(2-\eta_1-\eta_2+\eta_4\bigr)I+\bigl(-2-\eta_4\bigr)B && \\
\rho(T_{13}) &= \begin{pmatrix} \tfrac{1}{2}\eta_3 & 2-\eta_1+\eta_4 \\ -\tfrac{1}{4}\eta_1+\tfrac{1}{4}\eta_2 & -\tfrac{1}{2}\eta_3 \end{pmatrix} & (k_{13} &= 3) \\
 &= \bigl(-3+\eta_1+\eta_2-2\eta_4\bigr)A+\bigl(2-\eta_1+\eta_4\bigr)AB && \\
\rho(T_{14}) &= \begin{pmatrix} \tfrac{1}{2}\eta_3 & -2+\eta_1-\eta_4 \\ -\tfrac{1}{4}\eta_1+\tfrac{1}{4}\eta_2 & \tfrac{1}{2}\eta_3 \end{pmatrix} & (k_{14} &= 3) \\
 &= \bigl(4+2\eta_3+3\eta_4+\eta_8\bigr)I+\bigl(-2+\eta_1-\eta_4\bigr)B && \\
\rho(T_{15}) &= \begin{pmatrix} \tfrac{1}{2}\eta_3 & 2-\eta_1+\eta_4 \\ \tfrac{1}{4}\eta_1-\tfrac{1}{4}\eta_2 & \tfrac{1}{2}\eta_3 \end{pmatrix} & (k_{15} &= 3) \\
 &= \bigl(-3+\eta_1+\eta_2-2\eta_4\bigr)I+\bigl(2-\eta_1+\eta_4\bigr)B && \\
\rho(T_{16}) &= \begin{pmatrix} \tfrac{1}{2}\eta_3 & -2+\eta_1-\eta_4 \\ \tfrac{1}{4}\eta_1-\tfrac{1}{4}\eta_2 & -\tfrac{1}{2}\eta_3 \end{pmatrix} & (k_{16} &= 3) \\
 &= \bigl(4+2\eta_3+3\eta_4+\eta_8\bigr)A+\bigl(-2+\eta_1-\eta_4\bigr)AB && \\
\rho(T_{17}) &= \begin{pmatrix} \tfrac{1}{2}\eta_2 & 2-\eta_1-\eta_2+\eta_4 \\ -\tfrac{1}{4}\eta_2-\tfrac{1}{4}\eta_3+\tfrac{1}{2}\eta_8 & -\tfrac{1}{2}\eta_2 \end{pmatrix} & (k_{17} &= 3) \\
 &= \bigl(-5-2\eta_3-4\eta_4-\eta_8\bigr)A+\bigl(2-\eta_1-\eta_2+\eta_4\bigr)AB && \\
\rho(T_{18}) &= \begin{pmatrix} \tfrac{1}{2}\eta_2 & -2+\eta_1+\eta_2-\eta_4 \\ -\tfrac{1}{4}\eta_2-\tfrac{1}{4}\eta_3+\tfrac{1}{2}\eta_8 & \tfrac{1}{2}\eta_2 \end{pmatrix} & (k_{18} &= 3) \\
 &= \bigl(4-\eta_1+\eta_3+3\eta_4\bigr)I+\bigl(-2+\eta_1+\eta_2-\eta_4\bigr)B && \\
\rho(T_{19}) &= \begin{pmatrix} \tfrac{1}{2}\eta_1 & 5+2\eta_3+4\eta_4+\eta_8 \\ \tfrac{1}{4}\eta_3-\tfrac{1}{4}\eta_8 & -\tfrac{1}{2}\eta_1 \end{pmatrix} & (k_{19} &= 3) \\
 &= \bigl(-8+3\eta_1+2\eta_2-\eta_3-5\eta_4\bigr)A+\bigl(5+2\eta_3+4\eta_4+\eta_8\bigr)AB && \\
\rho(T_{20}) &= \begin{pmatrix} \tfrac{1}{2}\eta_1 & -5-2\eta_3-4\eta_4-\eta_8 \\ \tfrac{1}{4}\eta_3-\tfrac{1}{4}\eta_8 & \tfrac{1}{2}\eta_1 \end{pmatrix} & (k_{20} &= 3) \\
 &= \bigl(10+3\eta_3+7\eta_4+2\eta_8\bigr)I+\bigl(-5-2\eta_3-4\eta_4-\eta_8\bigr)B && \\
\rho(T_{21}) &= \begin{pmatrix} \tfrac{1}{2}\eta_1 & 5+2\eta_3+4\eta_4+\eta_8 \\ -\tfrac{1}{4}\eta_3+\tfrac{1}{4}\eta_8 & \tfrac{1}{2}\eta_1 \end{pmatrix} & (k_{21} &= 3) \\
 &= \bigl(-8+3\eta_1+2\eta_2-\eta_3-5\eta_4\bigr)I+\bigl(5+2\eta_3+4\eta_4+\eta_8\bigr)B && \\
\rho(T_{22}) &= \begin{pmatrix} \tfrac{1}{2}\eta_1 & -5-2\eta_3-4\eta_4-\eta_8 \\ -\tfrac{1}{4}\eta_3+\tfrac{1}{4}\eta_8 & -\tfrac{1}{2}\eta_1 \end{pmatrix} & (k_{22} &= 3) \\
 &= \bigl(10+3\eta_3+7\eta_4+2\eta_8\bigr)A+\bigl(-5-2\eta_3-4\eta_4-\eta_8\bigr)AB && .
\end{align*}



\begin{remark}[A uniform law for the Frobenius family]
In all three graphs F26A, CR(19), CR(31) the traceless valency-3 generator
satisfies
\[
T^2=(3+\eta)\,I
\]
where $3$ is the valency and $\eta$ is a sextic Gauss period of
$\Q(\zeta_p)$. For $p=13$ this is $3+\frac{-1\mp\sqrt{13}}{2}
=\frac{5\mp\sqrt{13}}{2}$, recovering the F26A relation. The ``why-$N$''
constant of the whole family is \emph{valency plus a Gauss period}.
\end{remark}

%──────────────────────────────────────────────────────────────────────
\section{Klein graph F056A: $\bigl(\mathrm{PSL}(2,7),\,C_3\bigr)$}

\textbf{Graph.} F056A, the 56-vertex cubic arc-transitive graph of girth 6
on the cosets $\mathrm{PSL}(2,7)/C_3$; $G=\mathrm{PSL}(2,7)$ ($|G|=168$),
$H=C_3$.

\textbf{Algebra.} Rank $20$,
$\HH\cong\C^3\oplus M_2(\C)\oplus M_3(\C)\oplus M_2(\C)$
(blocks on $\chi_4,\chi_5,\chi_6$ of degrees $6,7,8$).

\textbf{Highlights.} Block~1 is generated by the $\alpha$-diagonal orbital
($A^2+A+2I=0$; $2=\alpha\bar\alpha$ is the prime that splits in
$\Q(\sqrt{-7})$). Block~2, the integral $M_3$, carries the irreducible
characteristic polynomial $x^3+x-6$ --- the first genuinely cubic
eigenvalue in the catalogue, echoed by $\mathrm{PSL}(2,13)$'s $x^3+x-26$
below. In Block~3 the order-3 matrix
$\bigl(\begin{smallmatrix}1/2&1\\-7/4&-3/2\end{smallmatrix}\bigr)$ appears
as $\rho(T_{10})$ --- see the Remark after the Klein $\{7,3\}$ section.

% auto-generated by hecke_catalog_full.g — do not edit
\paragraph{One-dimensional characters.}
\[
\begin{array}{c|ccc}
 & \chi_{1}\,(1) & \chi_{2}\,(3) & \chi_{3}\,(3) \\\hline
T_{1} & 1 & 1 & 1 \\
T_{2} & 1 & -1 & -1 \\
T_{3} & 3 & \bar\alpha & \alpha \\
T_{4} & 3 & \alpha & \bar\alpha \\
T_{5} & 3 & -\bar\alpha & -\alpha \\
T_{6} & 3 & -\alpha & -\bar\alpha \\
T_{7} & 3 & -\alpha & -\bar\alpha \\
T_{8} & 3 & 1 & 1 \\
T_{9} & 3 & -1 & -1 \\
T_{10} & 3 & 2 & 2 \\
T_{11} & 3 & -2 & -2 \\
T_{12} & 3 & \bar\alpha & \alpha \\
T_{13} & 3 & 1+2\alpha & -1-2\alpha \\
T_{14} & 3 & \alpha & \bar\alpha \\
T_{15} & 3 & 1 & 1 \\
T_{16} & 3 & 2 & 2 \\
T_{17} & 3 & -\bar\alpha & -\alpha \\
T_{18} & 3 & -1 & -1 \\
T_{19} & 3 & -1-2\alpha & 1+2\alpha \\
T_{20} & 3 & -2 & -2 \\
\end{array}
\]

\subsubsection*{Block 1: $M_2(\C)$ on $\chi_{4}$ (degree 6)}
\noindent\textbf{Generators.}
\[
A = \rho(T_{3}) = \begin{pmatrix} \alpha & 0 \\ 0 & \bar\alpha \end{pmatrix}, \qquad
B = \rho(T_{8}) = \begin{pmatrix} 1 & 2 \\ 2 & 1 \end{pmatrix} .
\]
\textbf{Complete relations} (monomial basis $I,\;A,\;B,\;AB$ of $M_2(\C)$):
\begin{gather*}
A^2+A+2I = 0 ,\\
B^2-2\,B-3I = 0 ,\\
BA = I+2\,A-B-AB ,\\
ABA = -4I-A+2\,B ,\\
ABB = 3\,A+2\,AB .
\end{gather*}
\textbf{All orbital images} ($\rho(T_i)$, valency $k_i$, word in the generators):
\begin{align*}
\rho(T_{1}) &= \begin{pmatrix} 1 & 0 \\ 0 & 1 \end{pmatrix} = I & (k_{1} &= 1) \\
\rho(T_{2}) &= \begin{pmatrix} 1 & 0 \\ 0 & 1 \end{pmatrix} = I & (k_{2} &= 1) \\
\rho(T_{3}) &= \begin{pmatrix} \alpha & 0 \\ 0 & \bar\alpha \end{pmatrix} = A & (k_{3} &= 3) \\
\rho(T_{4}) &= \begin{pmatrix} \bar\alpha & 0 \\ 0 & \alpha \end{pmatrix} = -I-A & (k_{4} &= 3) \\
\rho(T_{5}) &= \begin{pmatrix} \alpha & 0 \\ 0 & \bar\alpha \end{pmatrix} = A & (k_{5} &= 3) \\
\rho(T_{6}) &= \begin{pmatrix} \bar\alpha & 0 \\ 0 & \alpha \end{pmatrix} = -I-A & (k_{6} &= 3) \\
\rho(T_{7}) &= \begin{pmatrix} \bar\alpha & 0 \\ 0 & \alpha \end{pmatrix} = -I-A & (k_{7} &= 3) \\
\rho(T_{8}) &= \begin{pmatrix} 1 & 2 \\ 2 & 1 \end{pmatrix} = B & (k_{8} &= 3) \\
\rho(T_{9}) &= \begin{pmatrix} 1 & -1+\alpha \\ -2-\alpha & 1 \end{pmatrix} = \tfrac{3}{2}I-\tfrac{1}{2}\,A-\tfrac{1}{2}\,B+\tfrac{1}{2}\,AB & (k_{9} &= 3) \\
\rho(T_{10}) &= \begin{pmatrix} 0 & \bar\alpha \\ \alpha & 0 \end{pmatrix} = \tfrac{1}{2}I+\tfrac{1}{2}\,A-\tfrac{1}{2}\,B-\tfrac{1}{2}\,AB & (k_{10} &= 3) \\
\rho(T_{11}) &= \begin{pmatrix} 0 & \bar\alpha \\ \alpha & 0 \end{pmatrix} = \tfrac{1}{2}I+\tfrac{1}{2}\,A-\tfrac{1}{2}\,B-\tfrac{1}{2}\,AB & (k_{11} &= 3) \\
\rho(T_{12}) &= \begin{pmatrix} \alpha & 0 \\ 0 & \bar\alpha \end{pmatrix} = A & (k_{12} &= 3) \\
\rho(T_{13}) &= \begin{pmatrix} -1 & -\bar\alpha \\ -\alpha & -1 \end{pmatrix} = -\tfrac{3}{2}I-\tfrac{1}{2}\,A+\tfrac{1}{2}\,B+\tfrac{1}{2}\,AB & (k_{13} &= 3) \\
\rho(T_{14}) &= \begin{pmatrix} \bar\alpha & 0 \\ 0 & \alpha \end{pmatrix} = -I-A & (k_{14} &= 3) \\
\rho(T_{15}) &= \begin{pmatrix} 1 & -1+\alpha \\ -2-\alpha & 1 \end{pmatrix} = \tfrac{3}{2}I-\tfrac{1}{2}\,A-\tfrac{1}{2}\,B+\tfrac{1}{2}\,AB & (k_{15} &= 3) \\
\rho(T_{16}) &= \begin{pmatrix} 0 & \bar\alpha \\ \alpha & 0 \end{pmatrix} = \tfrac{1}{2}I+\tfrac{1}{2}\,A-\tfrac{1}{2}\,B-\tfrac{1}{2}\,AB & (k_{16} &= 3) \\
\rho(T_{17}) &= \begin{pmatrix} \alpha & 0 \\ 0 & \bar\alpha \end{pmatrix} = A & (k_{17} &= 3) \\
\rho(T_{18}) &= \begin{pmatrix} 1 & 2 \\ 2 & 1 \end{pmatrix} = B & (k_{18} &= 3) \\
\rho(T_{19}) &= \begin{pmatrix} -1 & -\bar\alpha \\ -\alpha & -1 \end{pmatrix} = -\tfrac{3}{2}I-\tfrac{1}{2}\,A+\tfrac{1}{2}\,B+\tfrac{1}{2}\,AB & (k_{19} &= 3) \\
\rho(T_{20}) &= \begin{pmatrix} 0 & \bar\alpha \\ \alpha & 0 \end{pmatrix} = \tfrac{1}{2}I+\tfrac{1}{2}\,A-\tfrac{1}{2}\,B-\tfrac{1}{2}\,AB & (k_{20} &= 3) .
\end{align*}

\subsubsection*{Block 2: $M_3(\C)$ on $\chi_{5}$ (degree 7)}
\noindent\textbf{Generators.}
\[
A = \rho(T_{2}) = \begin{pmatrix} 1 & 0 & 0 \\ 0 & -1 & 0 \\ 0 & 0 & -1 \end{pmatrix}, \qquad
B = \rho(T_{3}) = \begin{pmatrix} 1 & 1 & 0 \\ 2 & 0 & 2 \\ 4 & 0 & 1 \end{pmatrix} .
\]
\textbf{Complete relations} (monomial basis $I,\;A,\;B,\;AB,\;BA,\;BB,\;ABB,\;BAB,\;BABB$ of $M_3(\C)$):
\begin{gather*}
A^2-I = 0 ,\\
B^3-2\,B^2-B-6I = 0 ,\\
ABA = 2I+2\,A-B-AB-BA ,\\
BAA = B ,\\
BBA = I+A+2\,B+2\,AB+2\,BA-BB-ABB-BAB ,\\
BBB = 6I+B+2\,BB ,\\
ABBA = 5I+5\,A-2\,B-2\,AB-2\,BA+BAB ,\\
ABBB = 6\,A+AB+2\,ABB ,\\
BABA = -I-A-2\,AB+ABB ,\\
BABBA = -8I-8\,A+B-4\,AB+BA+2\,ABB+2\,BAB-BABB ,\\
BABBB = 6\,BA+BAB+2\,BABB .
\end{gather*}
\textbf{All orbital images} ($\rho(T_i)$, valency $k_i$, word in the generators):
\begin{align*}
\rho(T_{1}) &= \begin{pmatrix} 1 & 0 & 0 \\ 0 & 1 & 0 \\ 0 & 0 & 1 \end{pmatrix} = I & (k_{1} &= 1) \\
\rho(T_{2}) &= \begin{pmatrix} 1 & 0 & 0 \\ 0 & -1 & 0 \\ 0 & 0 & -1 \end{pmatrix} = A & (k_{2} &= 1) \\
\rho(T_{3}) &= \begin{pmatrix} 1 & 1 & 0 \\ 2 & 0 & 2 \\ 4 & 0 & 1 \end{pmatrix} = B & (k_{3} &= 3) \\
\rho(T_{4}) &= \begin{pmatrix} 1 & 0 & 1 \\ 4 & 1 & 0 \\ 2 & 2 & 0 \end{pmatrix} = -\tfrac{1}{2}\,B+\tfrac{1}{2}\,BB & (k_{4} &= 3) \\
\rho(T_{5}) &= \begin{pmatrix} 1 & 1 & 0 \\ -2 & 0 & -2 \\ -4 & 0 & -1 \end{pmatrix} = AB & (k_{5} &= 3) \\
\rho(T_{6}) &= \begin{pmatrix} 1 & 0 & 1 \\ -4 & -1 & 0 \\ -2 & -2 & 0 \end{pmatrix} = -\tfrac{1}{2}\,AB+\tfrac{1}{2}\,ABB & (k_{6} &= 3) \\
\rho(T_{7}) &= \begin{pmatrix} 1 & 0 & -1 \\ 4 & -1 & 0 \\ 2 & -2 & 0 \end{pmatrix} & (k_{7} &= 3) \\
 &= \tfrac{1}{2}I+\tfrac{1}{2}\,A+B+AB+\tfrac{1}{2}\,BA-\tfrac{1}{2}\,BB-\tfrac{1}{2}\,ABB-\tfrac{1}{2}\,BAB && \\
\rho(T_{8}) &= \begin{pmatrix} -1 & 0 & 0 \\ 0 & 1 & 0 \\ 0 & 2 & -3 \end{pmatrix} & (k_{8} &= 3) \\
 &= -\tfrac{7}{4}I-\tfrac{1}{4}\,A-\tfrac{1}{2}\,B-AB-\tfrac{1}{4}\,BA+\tfrac{1}{4}\,BB+\tfrac{1}{4}\,ABB+\tfrac{1}{2}\,BAB-\tfrac{1}{4}\,BABB && \\
\rho(T_{9}) &= \begin{pmatrix} -1 & 0 & 0 \\ 0 & 3 & -2 \\ 0 & 0 & -1 \end{pmatrix} & (k_{9} &= 3) \\
 &= -I-\tfrac{5}{2}\,A-\tfrac{1}{4}\,AB+\tfrac{1}{2}\,BA+\tfrac{1}{4}\,ABB+\tfrac{1}{4}\,BAB-\tfrac{1}{4}\,BABB && \\
\rho(T_{10}) &= \begin{pmatrix} -1 & -1 & 1 \\ 2 & 0 & 0 \\ -2 & 0 & 0 \end{pmatrix} & (k_{10} &= 3) \\
 &= -\tfrac{1}{4}I-\tfrac{1}{4}\,A-B-\tfrac{1}{4}\,AB-\tfrac{1}{4}\,BA+\tfrac{1}{4}\,BB-\tfrac{1}{4}\,BAB && \\
\rho(T_{11}) &= \begin{pmatrix} -1 & 1 & -1 \\ 2 & 0 & 0 \\ -2 & 0 & 0 \end{pmatrix} & (k_{11} &= 3) \\
 &= -\tfrac{1}{4}I-\tfrac{1}{4}\,A+\tfrac{1}{2}\,B+\tfrac{5}{4}\,AB-\tfrac{1}{4}\,BA-\tfrac{1}{4}\,BB-\tfrac{1}{2}\,ABB-\tfrac{1}{4}\,BAB && \\
\rho(T_{12}) &= \begin{pmatrix} 1 & -1 & 0 \\ -2 & 0 & 2 \\ -4 & 0 & 1 \end{pmatrix} = 2I+2\,A-B-AB-BA & (k_{12} &= 3) \\
\rho(T_{13}) &= \begin{pmatrix} -1 & 0 & 0 \\ 0 & -1 & 2 \\ 0 & -2 & 1 \end{pmatrix} & (k_{13} &= 3) \\
 &= -\tfrac{1}{4}I-\tfrac{1}{4}\,A+\tfrac{1}{2}\,B-\tfrac{1}{4}\,AB-\tfrac{1}{4}\,BA-\tfrac{1}{4}\,BB-\tfrac{1}{4}\,BAB && \\
\rho(T_{14}) &= \begin{pmatrix} 1 & 0 & -1 \\ -4 & 1 & 0 \\ -2 & 2 & 0 \end{pmatrix} = \tfrac{3}{2}I+\tfrac{3}{2}\,A-\tfrac{1}{2}\,B-\tfrac{1}{2}\,AB-\tfrac{1}{2}\,BA+\tfrac{1}{2}\,BAB & (k_{14} &= 3) \\
\rho(T_{15}) &= \begin{pmatrix} -1 & 0 & 0 \\ 0 & -3 & 2 \\ 0 & 0 & 1 \end{pmatrix} & (k_{15} &= 3) \\
 &= \tfrac{3}{2}\,A+\tfrac{1}{4}\,AB-\tfrac{1}{2}\,BA-\tfrac{1}{4}\,ABB-\tfrac{1}{4}\,BAB+\tfrac{1}{4}\,BABB && \\
\rho(T_{16}) &= \begin{pmatrix} -1 & 1 & -1 \\ -2 & 0 & 0 \\ 2 & 0 & 0 \end{pmatrix} & (k_{16} &= 3) \\
 &= -\tfrac{3}{4}I-\tfrac{3}{4}\,A+B+\tfrac{1}{4}\,AB+\tfrac{1}{4}\,BA-\tfrac{1}{4}\,BB+\tfrac{1}{4}\,BAB && \\
\rho(T_{17}) &= \begin{pmatrix} 1 & -1 & 0 \\ 2 & 0 & -2 \\ 4 & 0 & -1 \end{pmatrix} = BA & (k_{17} &= 3) \\
\rho(T_{18}) &= \begin{pmatrix} -1 & 0 & 0 \\ 0 & -1 & 0 \\ 0 & -2 & 3 \end{pmatrix} & (k_{18} &= 3) \\
 &= \tfrac{3}{4}I-\tfrac{3}{4}\,A+\tfrac{1}{2}\,B+AB+\tfrac{1}{4}\,BA-\tfrac{1}{4}\,BB-\tfrac{1}{4}\,ABB-\tfrac{1}{2}\,BAB+\tfrac{1}{4}\,BABB && \\
\rho(T_{19}) &= \begin{pmatrix} -1 & 0 & 0 \\ 0 & 1 & -2 \\ 0 & 2 & -1 \end{pmatrix} & (k_{19} &= 3) \\
 &= -\tfrac{3}{4}I-\tfrac{3}{4}\,A-\tfrac{1}{2}\,B+\tfrac{1}{4}\,AB+\tfrac{1}{4}\,BA+\tfrac{1}{4}\,BB+\tfrac{1}{4}\,BAB && \\
\rho(T_{20}) &= \begin{pmatrix} -1 & -1 & 1 \\ -2 & 0 & 0 \\ 2 & 0 & 0 \end{pmatrix} & (k_{20} &= 3) \\
 &= -\tfrac{3}{4}I-\tfrac{3}{4}\,A-\tfrac{1}{2}\,B-\tfrac{5}{4}\,AB+\tfrac{1}{4}\,BA+\tfrac{1}{4}\,BB+\tfrac{1}{2}\,ABB+\tfrac{1}{4}\,BAB && .
\end{align*}

\subsubsection*{Block 3: $M_2(\C)$ on $\chi_{6}$ (degree 8)}
\noindent\textbf{Generators.}
\[
A = \rho(T_{2}) = \begin{pmatrix} 1 & 0 \\ 0 & -1 \end{pmatrix}, \qquad
B = \rho(T_{3}) = \begin{pmatrix} -\tfrac{1}{2} & 1 \\ -\tfrac{7}{4} & -\tfrac{1}{2} \end{pmatrix} .
\]
\textbf{Complete relations} (monomial basis $I,\;A,\;B,\;AB$ of $M_2(\C)$):
\begin{gather*}
A^2-I = 0 ,\\
B^2+B+2I = 0 ,\\
BA = -A-AB ,\\
ABA = -I-B ,\\
ABB = -2\,A-AB .
\end{gather*}
\textbf{All orbital images} ($\rho(T_i)$, valency $k_i$, word in the generators):
\begin{align*}
\rho(T_{1}) &= \begin{pmatrix} 1 & 0 \\ 0 & 1 \end{pmatrix} = I & (k_{1} &= 1) \\
\rho(T_{2}) &= \begin{pmatrix} 1 & 0 \\ 0 & -1 \end{pmatrix} = A & (k_{2} &= 1) \\
\rho(T_{3}) &= \begin{pmatrix} -\tfrac{1}{2} & 1 \\ -\tfrac{7}{4} & -\tfrac{1}{2} \end{pmatrix} = B & (k_{3} &= 3) \\
\rho(T_{4}) &= \begin{pmatrix} -\tfrac{1}{2} & -1 \\ \tfrac{7}{4} & -\tfrac{1}{2} \end{pmatrix} = -I-B & (k_{4} &= 3) \\
\rho(T_{5}) &= \begin{pmatrix} -\tfrac{1}{2} & 1 \\ \tfrac{7}{4} & \tfrac{1}{2} \end{pmatrix} = AB & (k_{5} &= 3) \\
\rho(T_{6}) &= \begin{pmatrix} -\tfrac{1}{2} & -1 \\ -\tfrac{7}{4} & \tfrac{1}{2} \end{pmatrix} = -A-AB & (k_{6} &= 3) \\
\rho(T_{7}) &= \begin{pmatrix} -\tfrac{1}{2} & 1 \\ \tfrac{7}{4} & \tfrac{1}{2} \end{pmatrix} = AB & (k_{7} &= 3) \\
\rho(T_{8}) &= \begin{pmatrix} -1 & 0 \\ 0 & 1 \end{pmatrix} = -A & (k_{8} &= 3) \\
\rho(T_{9}) &= \begin{pmatrix} -1 & 0 \\ 0 & -1 \end{pmatrix} = -I & (k_{9} &= 3) \\
\rho(T_{10}) &= \begin{pmatrix} \tfrac{1}{2} & 1 \\ -\tfrac{7}{4} & -\tfrac{3}{2} \end{pmatrix} = A+B & (k_{10} &= 3) \\
\rho(T_{11}) &= \begin{pmatrix} \tfrac{1}{2} & -1 \\ -\tfrac{7}{4} & \tfrac{3}{2} \end{pmatrix} = I-A-AB & (k_{11} &= 3) \\
\rho(T_{12}) &= \begin{pmatrix} -\tfrac{1}{2} & -1 \\ \tfrac{7}{4} & -\tfrac{1}{2} \end{pmatrix} = -I-B & (k_{12} &= 3) \\
\rho(T_{13}) &= \begin{pmatrix} 2 & 0 \\ 0 & 0 \end{pmatrix} = I+A & (k_{13} &= 3) \\
\rho(T_{14}) &= \begin{pmatrix} -\tfrac{1}{2} & 1 \\ -\tfrac{7}{4} & -\tfrac{1}{2} \end{pmatrix} = B & (k_{14} &= 3) \\
\rho(T_{15}) &= \begin{pmatrix} -1 & 0 \\ 0 & 1 \end{pmatrix} = -A & (k_{15} &= 3) \\
\rho(T_{16}) &= \begin{pmatrix} \tfrac{1}{2} & -1 \\ \tfrac{7}{4} & -\tfrac{3}{2} \end{pmatrix} = -I+A-B & (k_{16} &= 3) \\
\rho(T_{17}) &= \begin{pmatrix} -\tfrac{1}{2} & -1 \\ -\tfrac{7}{4} & \tfrac{1}{2} \end{pmatrix} = -A-AB & (k_{17} &= 3) \\
\rho(T_{18}) &= \begin{pmatrix} -1 & 0 \\ 0 & -1 \end{pmatrix} = -I & (k_{18} &= 3) \\
\rho(T_{19}) &= \begin{pmatrix} 2 & 0 \\ 0 & 0 \end{pmatrix} = I+A & (k_{19} &= 3) \\
\rho(T_{20}) &= \begin{pmatrix} \tfrac{1}{2} & 1 \\ \tfrac{7}{4} & \tfrac{3}{2} \end{pmatrix} = I+AB & (k_{20} &= 3) .
\end{align*}



%──────────────────────────────────────────────────────────────────────
\section{Klein $\{3,7\}$ graph: $\bigl(\mathrm{PGL}(2,7),\,D_7\bigr)$}

\textbf{Graph.} The 1-skeleton of the $\{3,7\}$ tiling of the Klein
quartic: 24 vertices, valency 7 (the dual of the $\{7,3\}$ map).
$G=\operatorname{Aut}\Gamma\cong\mathrm{PGL}(2,7)$ ($|G|=336$,
GAP-verified), with vertex stabilizer $H\cong D_7$ (dihedral of order
14).

\textbf{Algebra.} Rank $6$, $\HH\cong\C^2\oplus M_2(\C)$.

\textbf{Highlights.} The stabilizer fixes three vertices, and the two
non-trivial valency-1 orbitals act as the $\omega$-diagonal pair
($A^2+A+I=0$, so $\rho(T_2)$ has order 3). The symmetric valency-7
orbital satisfies
\[
B^2=7I
\]
--- the ``why-$N$'' constant equals the valency on the nose, the cleanest
instance in the catalogue --- and the chiral valency-7 pair is exactly
$\rho(T_5)=BA$, $\rho(T_6)=AB$, with $AB+BA=-B$.

% auto-generated by hecke_catalog_full.g — do not edit
\paragraph{One-dimensional characters.}
\[
\begin{array}{c|cc}
 & \chi_{1}\,(1) & \chi_{7}\,(7) \\\hline
T_{1} & 1 & 1 \\
T_{2} & 1 & 1 \\
T_{3} & 1 & 1 \\
T_{4} & 7 & -1 \\
T_{5} & 7 & -1 \\
T_{6} & 7 & -1 \\
\end{array}
\]

\subsubsection*{Block 1: $M_2(\C)$ on $\chi_{9}$ (degree 8)}
\noindent\textbf{Generators.}
\[
A = \rho(T_{2}) = \begin{pmatrix} \omega & 0 \\ 0 & \bar\omega \end{pmatrix}, \qquad
B = \rho(T_{4}) = \begin{pmatrix} 0 & \sqrt7 \\ \sqrt7 & 0 \end{pmatrix} .
\]
\textbf{Complete relations} (monomial basis $I,\;A,\;B,\;AB$ of $M_2(\C)$):
\begin{gather*}
A^2+A+I = 0 ,\\
B^2-7I = 0 ,\\
BA = -B-AB ,\\
ABA = B ,\\
ABB = 7\,A .
\end{gather*}
\textbf{All orbital images} ($\rho(T_i)$, valency $k_i$, word in the generators):
\begin{align*}
\rho(T_{1}) &= \begin{pmatrix} 1 & 0 \\ 0 & 1 \end{pmatrix} = I & (k_{1} &= 1) \\
\rho(T_{2}) &= \begin{pmatrix} \omega & 0 \\ 0 & \bar\omega \end{pmatrix} = A & (k_{2} &= 1) \\
\rho(T_{3}) &= \begin{pmatrix} \bar\omega & 0 \\ 0 & \omega \end{pmatrix} = -I-A & (k_{3} &= 1) \\
\rho(T_{4}) &= \begin{pmatrix} 0 & \sqrt7 \\ \sqrt7 & 0 \end{pmatrix} = B & (k_{4} &= 7) \\
\rho(T_{5}) &= \begin{pmatrix} 0 & \bar\omega\sqrt7 \\ \omega\sqrt7 & 0 \end{pmatrix} = -B-AB & (k_{5} &= 7) \\
\rho(T_{6}) &= \begin{pmatrix} 0 & \omega\sqrt7 \\ \bar\omega\sqrt7 & 0 \end{pmatrix} = AB & (k_{6} &= 7) .
\end{align*}



%──────────────────────────────────────────────────────────────────────
\section{Klein $\{7,3\}$ graph: $\bigl(\mathrm{PGL}(2,7),\,S_3\bigr)$}

\textbf{Graph.} The 1-skeleton of the $\{7,3\}$ tiling of the Klein quartic
(56 vertices, cubic, girth 7): the \emph{same} coset space
$\mathrm{PSL}(2,7)/C_3$ as F056A, but a girth-7 orbital. The full group is
$G=\operatorname{Aut}\Gamma\cong\mathrm{PGL}(2,7)$ ($|G|=336$; GAP-verified,
and visible in the character data: the degree-6 constituents take values in
$\Q(\zeta_8)$, so $G$ has elements of order 8 --- ruling out
$\mathrm{PSL}(2,7)\times C_2$), with vertex stabilizer $H\cong S_3$.

\textbf{Algebra.} Rank $13$,
$\HH\cong\C^5\oplus M_2(\C)\oplus M_2(\C)$.

\textbf{Highlights.} Block~1 is integral with
$B^2+B+4I=0$ and the companion orbital relations producing eigenvalues
$\frac{-1\pm\sqrt{17}}{2}$; Block~2 contains an orbital of order 3 and one
with square $8I$.

% auto-generated by hecke_catalog_full.g — do not edit
\paragraph{One-dimensional characters.}
\[
\begin{array}{c|ccccc}
 & \chi_{1}\,(1) & \chi_{3}\,(6) & \chi_{4}\,(6) & \chi_{5}\,(6) & \chi_{6}\,(7) \\\hline
T_{1} & 1 & 1 & 1 & 1 & 1 \\
T_{2} & 1 & -1 & 1 & 1 & -1 \\
T_{3} & 3 & 2 & -\sqrt2 & \sqrt2 & 0 \\
T_{4} & 3 & -2 & -\sqrt2 & \sqrt2 & 0 \\
T_{5} & 3 & -2 & -\sqrt2 & \sqrt2 & 0 \\
T_{6} & 3 & 2 & -\sqrt2 & \sqrt2 & 0 \\
T_{7} & 6 & -1 & -1 & -1 & 3 \\
T_{8} & 6 & -1 & -1 & -1 & 3 \\
T_{9} & 6 & 0 & -2+2\sqrt2 & -2-2\sqrt2 & 0 \\
T_{10} & 6 & -2 & 2+\sqrt2 & 2-\sqrt2 & 0 \\
T_{11} & 6 & 1 & -1 & -1 & -3 \\
T_{12} & 6 & 2 & 2+\sqrt2 & 2-\sqrt2 & 0 \\
T_{13} & 6 & 1 & -1 & -1 & -3 \\
\end{array}
\]

\subsubsection*{Block 1: $M_2(\C)$ on $\chi_{7}$ (degree 7)}
\noindent\textbf{Generators.}
\[
A = \rho(T_{2}) = \begin{pmatrix} 1 & 0 \\ 0 & -1 \end{pmatrix}, \qquad
B = \rho(T_{3}) = \begin{pmatrix} -1 & 1 \\ -4 & 0 \end{pmatrix} .
\]
\textbf{Complete relations} (monomial basis $I,\;A,\;B,\;AB$ of $M_2(\C)$):
\begin{gather*}
A^2-I = 0 ,\\
B^2+B+4I = 0 ,\\
BA = -I-A-AB ,\\
ABA = -I-A-B ,\\
ABB = -4\,A-AB .
\end{gather*}
\textbf{All orbital images} ($\rho(T_i)$, valency $k_i$, word in the generators):
\begin{align*}
\rho(T_{1}) &= \begin{pmatrix} 1 & 0 \\ 0 & 1 \end{pmatrix} = I & (k_{1} &= 1) \\
\rho(T_{2}) &= \begin{pmatrix} 1 & 0 \\ 0 & -1 \end{pmatrix} = A & (k_{2} &= 1) \\
\rho(T_{3}) &= \begin{pmatrix} -1 & 1 \\ -4 & 0 \end{pmatrix} = B & (k_{3} &= 3) \\
\rho(T_{4}) &= \begin{pmatrix} -1 & -1 \\ -4 & 0 \end{pmatrix} = -I-A-AB & (k_{4} &= 3) \\
\rho(T_{5}) &= \begin{pmatrix} -1 & 1 \\ 4 & 0 \end{pmatrix} = AB & (k_{5} &= 3) \\
\rho(T_{6}) &= \begin{pmatrix} -1 & -1 \\ 4 & 0 \end{pmatrix} = -I-A-B & (k_{6} &= 3) \\
\rho(T_{7}) &= \begin{pmatrix} 2 & 1 \\ -4 & -1 \end{pmatrix} = I+2\,A+B & (k_{7} &= 6) \\
\rho(T_{8}) &= \begin{pmatrix} 2 & -1 \\ 4 & -1 \end{pmatrix} = A-B & (k_{8} &= 6) \\
\rho(T_{9}) &= \begin{pmatrix} -2 & 0 \\ 0 & 0 \end{pmatrix} = -I-A & (k_{9} &= 6) \\
\rho(T_{10}) &= \begin{pmatrix} -2 & 0 \\ 0 & 4 \end{pmatrix} = I-3\,A & (k_{10} &= 6) \\
\rho(T_{11}) &= \begin{pmatrix} 2 & -1 \\ -4 & 1 \end{pmatrix} = I-AB & (k_{11} &= 6) \\
\rho(T_{12}) &= \begin{pmatrix} -2 & 0 \\ 0 & -4 \end{pmatrix} = -3I+A & (k_{12} &= 6) \\
\rho(T_{13}) &= \begin{pmatrix} 2 & 1 \\ 4 & 1 \end{pmatrix} = 2I+A+AB & (k_{13} &= 6) .
\end{align*}

\subsubsection*{Block 2: $M_2(\C)$ on $\chi_{9}$ (degree 8)}
\noindent\textbf{Generators.}
\[
A = \rho(T_{2}) = \begin{pmatrix} 1 & 0 \\ 0 & -1 \end{pmatrix}, \qquad
B = \rho(T_{3}) = \begin{pmatrix} \tfrac{1}{2} & 1 \\ -\tfrac{7}{4} & -\tfrac{3}{2} \end{pmatrix} .
\]
\textbf{Complete relations} (monomial basis $I,\;A,\;B,\;AB$ of $M_2(\C)$):
\begin{gather*}
A^2-I = 0 ,\\
B^2+B+I = 0 ,\\
BA = 2I-A-AB ,\\
ABA = -I+2\,A-B ,\\
ABB = -A-AB .
\end{gather*}
\textbf{All orbital images} ($\rho(T_i)$, valency $k_i$, word in the generators):
\begin{align*}
\rho(T_{1}) &= \begin{pmatrix} 1 & 0 \\ 0 & 1 \end{pmatrix} = I & (k_{1} &= 1) \\
\rho(T_{2}) &= \begin{pmatrix} 1 & 0 \\ 0 & -1 \end{pmatrix} = A & (k_{2} &= 1) \\
\rho(T_{3}) &= \begin{pmatrix} \tfrac{1}{2} & 1 \\ -\tfrac{7}{4} & -\tfrac{3}{2} \end{pmatrix} = B & (k_{3} &= 3) \\
\rho(T_{4}) &= \begin{pmatrix} \tfrac{1}{2} & -1 \\ -\tfrac{7}{4} & \tfrac{3}{2} \end{pmatrix} = 2I-A-AB & (k_{4} &= 3) \\
\rho(T_{5}) &= \begin{pmatrix} \tfrac{1}{2} & 1 \\ \tfrac{7}{4} & \tfrac{3}{2} \end{pmatrix} = AB & (k_{5} &= 3) \\
\rho(T_{6}) &= \begin{pmatrix} \tfrac{1}{2} & -1 \\ \tfrac{7}{4} & -\tfrac{3}{2} \end{pmatrix} = -I+2\,A-B & (k_{6} &= 3) \\
\rho(T_{7}) &= \begin{pmatrix} -1 & -2 \\ \tfrac{7}{2} & -1 \end{pmatrix} = -2I+2\,A-2\,B & (k_{7} &= 6) \\
\rho(T_{8}) &= \begin{pmatrix} -1 & 2 \\ -\tfrac{7}{2} & -1 \end{pmatrix} = -2\,A+2\,B & (k_{8} &= 6) \\
\rho(T_{9}) &= \begin{pmatrix} 4 & 0 \\ 0 & 0 \end{pmatrix} = 2I+2\,A & (k_{9} &= 6) \\
\rho(T_{10}) &= \begin{pmatrix} -2 & 0 \\ 0 & -2 \end{pmatrix} = -2I & (k_{10} &= 6) \\
\rho(T_{11}) &= \begin{pmatrix} -1 & 2 \\ \tfrac{7}{2} & 1 \end{pmatrix} = -2I+2\,AB & (k_{11} &= 6) \\
\rho(T_{12}) &= \begin{pmatrix} -2 & 0 \\ 0 & 2 \end{pmatrix} = -2\,A & (k_{12} &= 6) \\
\rho(T_{13}) &= \begin{pmatrix} -1 & -2 \\ -\tfrac{7}{2} & 1 \end{pmatrix} = 2I-2\,A-2\,AB & (k_{13} &= 6) .
\end{align*}



\begin{remark}
The order-3 matrix
$\bigl(\begin{smallmatrix}1/2&1\\-7/4&-3/2\end{smallmatrix}\bigr)$
is \emph{literally the same} as $\rho(T_{10})$ in F056A's Block~3 --- same
coset geometry, different orbital, identical $2\times2$ block.
\end{remark}

%──────────────────────────────────────────────────────────────────────
\section{Zhou-3 (F182A): $\bigl(\mathrm{PSL}(2,13),\,S_3\bigr)$}

\textbf{Graph.} The 182-vertex cubic coset graph of girth 7;
$G=\mathrm{PSL}(2,13)$ ($|G|=1092$), $H=S_3$.

\textbf{Algebra.} Rank $34$,
$\HH\cong\C^4\oplus M_2(\C)^{\oplus3}\oplus M_3(\C)^{\oplus2}$.

\textbf{Highlights.} The three $M_2$ blocks (degrees $12,12,12$) are Galois
conjugate over $\Q(\zeta_7)^+$ via $c_k\mapsto c_{2k}$. Block~4 is the
integral Steinberg $M_3$ with $B^3=-I$ and an orbital of characteristic
polynomial $x^3+x-26$, echoing $\mathrm{PSL}(2,7)$'s $x^3+x-6$. Block~5
carries $B^3=5I$ --- the first \emph{cubic} analogue of Meinhold's
$B^2=13I$. The $\pm\sqrt{13}$ of $\Q(\zeta_{13})$ appears in the
1-dimensional characters (quadratic Gauss periods
$\frac{-1\pm\sqrt{13}}{2}$), a different 13 from Meinhold's.

% auto-generated by hecke_catalog_full.g — do not edit
\paragraph{One-dimensional characters.}
\[
\begin{array}{c|cccc}
 & \chi_{1}\,(1) & \chi_{2}\,(7) & \chi_{3}\,(7) & \chi_{8}\,(14) \\\hline
T_{1} & 1 & 1 & 1 & 1 \\
T_{2} & 1 & -1 & -1 & -1 \\
T_{3} & 3 & \tfrac{1}{2}+\tfrac{1}{2}\sqrt{13} & \tfrac{1}{2}-\tfrac{1}{2}\sqrt{13} & 1 \\
T_{4} & 3 & -\tfrac{1}{2}-\tfrac{1}{2}\sqrt{13} & -\tfrac{1}{2}+\tfrac{1}{2}\sqrt{13} & -1 \\
T_{5} & 3 & \tfrac{1}{2}+\tfrac{1}{2}\sqrt{13} & \tfrac{1}{2}-\tfrac{1}{2}\sqrt{13} & 1 \\
T_{6} & 3 & -\tfrac{1}{2}-\tfrac{1}{2}\sqrt{13} & -\tfrac{1}{2}+\tfrac{1}{2}\sqrt{13} & -1 \\
T_{7} & 6 & \tfrac{1}{2}+\tfrac{1}{2}\sqrt{13} & \tfrac{1}{2}-\tfrac{1}{2}\sqrt{13} & -2 \\
T_{8} & 6 & 0 & 0 & 0 \\
T_{9} & 6 & -\tfrac{1}{2}-\tfrac{1}{2}\sqrt{13} & -\tfrac{1}{2}+\tfrac{1}{2}\sqrt{13} & 2 \\
T_{10} & 6 & \tfrac{5}{2}-\tfrac{1}{2}\sqrt{13} & \tfrac{5}{2}+\tfrac{1}{2}\sqrt{13} & -1 \\
T_{11} & 6 & -\tfrac{5}{2}+\tfrac{1}{2}\sqrt{13} & -\tfrac{5}{2}-\tfrac{1}{2}\sqrt{13} & 1 \\
T_{12} & 6 & 0 & 0 & 0 \\
T_{13} & 6 & -\tfrac{1}{2}-\tfrac{1}{2}\sqrt{13} & -\tfrac{1}{2}+\tfrac{1}{2}\sqrt{13} & 2 \\
T_{14} & 6 & \tfrac{1}{2}+\tfrac{1}{2}\sqrt{13} & \tfrac{1}{2}-\tfrac{1}{2}\sqrt{13} & -2 \\
T_{15} & 6 & -\tfrac{5}{2}+\tfrac{1}{2}\sqrt{13} & -\tfrac{5}{2}-\tfrac{1}{2}\sqrt{13} & 1 \\
T_{16} & 6 & \tfrac{5}{2}-\tfrac{1}{2}\sqrt{13} & \tfrac{5}{2}+\tfrac{1}{2}\sqrt{13} & -1 \\
T_{17} & 6 & -\tfrac{1}{2}-\tfrac{1}{2}\sqrt{13} & -\tfrac{1}{2}+\tfrac{1}{2}\sqrt{13} & 2 \\
T_{18} & 6 & \tfrac{1}{2}+\tfrac{1}{2}\sqrt{13} & \tfrac{1}{2}-\tfrac{1}{2}\sqrt{13} & -2 \\
T_{19} & 6 & \tfrac{1}{2}+\tfrac{1}{2}\sqrt{13} & \tfrac{1}{2}-\tfrac{1}{2}\sqrt{13} & -2 \\
T_{20} & 6 & -\tfrac{1}{2}-\tfrac{1}{2}\sqrt{13} & -\tfrac{1}{2}+\tfrac{1}{2}\sqrt{13} & 2 \\
T_{21} & 6 & 3 & 3 & 0 \\
T_{22} & 6 & -3 & -3 & 0 \\
T_{23} & 6 & 0 & 0 & 0 \\
T_{24} & 6 & \tfrac{1}{2}+\tfrac{1}{2}\sqrt{13} & \tfrac{1}{2}-\tfrac{1}{2}\sqrt{13} & 1 \\
T_{25} & 6 & -\tfrac{1}{2}-\tfrac{1}{2}\sqrt{13} & -\tfrac{1}{2}+\tfrac{1}{2}\sqrt{13} & -1 \\
T_{26} & 6 & 3 & 3 & 0 \\
T_{27} & 6 & 0 & 0 & 3 \\
T_{28} & 6 & 0 & 0 & -3 \\
T_{29} & 6 & 0 & 0 & 0 \\
T_{30} & 6 & \tfrac{1}{2}+\tfrac{1}{2}\sqrt{13} & \tfrac{1}{2}-\tfrac{1}{2}\sqrt{13} & 1 \\
T_{31} & 6 & -\tfrac{1}{2}-\tfrac{1}{2}\sqrt{13} & -\tfrac{1}{2}+\tfrac{1}{2}\sqrt{13} & -1 \\
T_{32} & 6 & 0 & 0 & 0 \\
T_{33} & 6 & -3 & -3 & 0 \\
T_{34} & 6 & 0 & 0 & 0 \\
\end{array}
\]

\subsubsection*{Block 1: $M_2(\C)$ on $\chi_{4}$ (degree 12)}
\noindent\textbf{Generators.}
\[
A = \rho(T_{2}) = \begin{pmatrix} 1 & 0 \\ 0 & -1 \end{pmatrix}, \qquad
B = \rho(T_{3}) = \begin{pmatrix} -\tfrac{1}{2}+\tfrac{1}{2}c_3 & 1 \\ -\tfrac{11}{4}-\tfrac{5}{4}c_1-c_2 & -\tfrac{1}{2}c_1+\tfrac{1}{2}c_2 \end{pmatrix} .
\]
\textbf{Complete relations} (monomial basis $I,\;A,\;B,\;AB$ of $M_2(\C)$):
\begin{gather*}
A^2-I = 0 ,\\
B^2+\bigl(1+c_1\bigr)B+\bigl(3+2c_1+c_2\bigr)I = 0 ,\\
BA = \bigl(-1-c_2\bigr)I+\bigl(-1-c_1\bigr)A-AB ,\\
ABA = \bigl(-1-c_1\bigr)I+\bigl(-1-c_2\bigr)A-B ,\\
ABB = \bigl(-3-2c_1-c_2\bigr)A+\bigl(-1-c_1\bigr)AB .
\end{gather*}
\textbf{All orbital images} ($\rho(T_i)$, valency $k_i$, word in the generators):
\begin{align*}
\rho(T_{1}) &= \begin{pmatrix} 1 & 0 \\ 0 & 1 \end{pmatrix} = I & (k_{1} &= 1) \\
\rho(T_{2}) &= \begin{pmatrix} 1 & 0 \\ 0 & -1 \end{pmatrix} = A & (k_{2} &= 1) \\
\rho(T_{3}) &= \begin{pmatrix} -\tfrac{1}{2}+\tfrac{1}{2}c_3 & 1 \\ -\tfrac{11}{4}-\tfrac{5}{4}c_1-c_2 & -\tfrac{1}{2}c_1+\tfrac{1}{2}c_2 \end{pmatrix} = B & (k_{3} &= 3) \\
\rho(T_{4}) &= \begin{pmatrix} -\tfrac{1}{2}+\tfrac{1}{2}c_3 & -1 \\ -\tfrac{11}{4}-\tfrac{5}{4}c_1-c_2 & \tfrac{1}{2}c_1-\tfrac{1}{2}c_2 \end{pmatrix} & (k_{4} &= 3) \\
 &= \bigl(-1-c_2\bigr)I+\bigl(-1-c_1\bigr)A-AB && \\
\rho(T_{5}) &= \begin{pmatrix} -\tfrac{1}{2}+\tfrac{1}{2}c_3 & -1 \\ \tfrac{11}{4}+\tfrac{5}{4}c_1+c_2 & -\tfrac{1}{2}c_1+\tfrac{1}{2}c_2 \end{pmatrix} & (k_{5} &= 3) \\
 &= \bigl(-1-c_1\bigr)I+\bigl(-1-c_2\bigr)A-B && \\
\rho(T_{6}) &= \begin{pmatrix} -\tfrac{1}{2}+\tfrac{1}{2}c_3 & 1 \\ \tfrac{11}{4}+\tfrac{5}{4}c_1+c_2 & \tfrac{1}{2}c_1-\tfrac{1}{2}c_2 \end{pmatrix} = AB & (k_{6} &= 3) \\
\rho(T_{7}) &= \begin{pmatrix} -\tfrac{1}{2}c_1+c_2 & c_2 \\ -\tfrac{1}{4}-c_1+\tfrac{1}{2}c_2 & \tfrac{1}{2}c_1+c_3 \end{pmatrix} = -c_1I+c_2\,A+c_2\,AB & (k_{7} &= 6) \\
\rho(T_{8}) &= \begin{pmatrix} -2c_2 & 0 \\ 0 & 0 \end{pmatrix} = -c_2I-c_2\,A & (k_{8} &= 6) \\
\rho(T_{9}) &= \begin{pmatrix} -\tfrac{1}{2}c_1+c_2 & -c_2 \\ -\tfrac{1}{4}-c_1+\tfrac{1}{2}c_2 & -\tfrac{1}{2}c_1-c_3 \end{pmatrix} = \bigl(1+c_2\bigr)I-A-c_2\,B & (k_{9} &= 6) \\
\rho(T_{10}) &= \begin{pmatrix} -\tfrac{1}{2} & -c_1-2c_2 \\ 1+\tfrac{3}{4}c_1+\tfrac{5}{4}c_2 & -\tfrac{1}{2}-c_3 \end{pmatrix} = -A+\bigl(-c_1-2c_2\bigr)B & (k_{10} &= 6) \\
\rho(T_{11}) &= \begin{pmatrix} -\tfrac{1}{2} & c_1+2c_2 \\ 1+\tfrac{3}{4}c_1+\tfrac{5}{4}c_2 & \tfrac{1}{2}+c_3 \end{pmatrix} & (k_{11} &= 6) \\
 &= \bigl(-c_1-c_2\bigr)I+\bigl(c_1+c_2\bigr)A+\bigl(c_1+2c_2\bigr)AB && \\
\rho(T_{12}) &= \begin{pmatrix} -2c_2 & 0 \\ 0 & 0 \end{pmatrix} = -c_2I-c_2\,A & (k_{12} &= 6) \\
\rho(T_{13}) &= \begin{pmatrix} -\tfrac{1}{2}c_1+c_2 & c_2 \\ \tfrac{1}{4}+c_1-\tfrac{1}{2}c_2 & -\tfrac{1}{2}c_1-c_3 \end{pmatrix} = c_2I-c_1\,A+c_2\,B & (k_{13} &= 6) \\
\rho(T_{14}) &= \begin{pmatrix} -\tfrac{1}{2}c_1+c_2 & -c_2 \\ \tfrac{1}{4}+c_1-\tfrac{1}{2}c_2 & \tfrac{1}{2}c_1+c_3 \end{pmatrix} = -I+\bigl(1+c_2\bigr)A-c_2\,AB & (k_{14} &= 6) \\
\rho(T_{15}) &= \begin{pmatrix} -\tfrac{1}{2} & -c_1-2c_2 \\ -1-\tfrac{3}{4}c_1-\tfrac{5}{4}c_2 & \tfrac{1}{2}+c_3 \end{pmatrix} = -I+\bigl(-c_1-2c_2\bigr)AB & (k_{15} &= 6) \\
\rho(T_{16}) &= \begin{pmatrix} -\tfrac{1}{2} & c_1+2c_2 \\ -1-\tfrac{3}{4}c_1-\tfrac{5}{4}c_2 & -\tfrac{1}{2}-c_3 \end{pmatrix} & (k_{16} &= 6) \\
 &= \bigl(c_1+c_2\bigr)I+\bigl(-c_1-c_2\bigr)A+\bigl(c_1+2c_2\bigr)B && \\
\rho(T_{17}) &= \begin{pmatrix} 1+2c_1+\tfrac{3}{2}c_2 & -1-c_2 \\ \tfrac{5}{2}+\tfrac{1}{4}c_1+\tfrac{3}{2}c_2 & -\tfrac{1}{2}c_2+c_3 \end{pmatrix} & (k_{17} &= 6) \\
 &= \bigl(2c_1+c_2\bigr)A+\bigl(-1-c_2\bigr)B && \\
\rho(T_{18}) &= \begin{pmatrix} 1+2c_1+\tfrac{3}{2}c_2 & 1+c_2 \\ \tfrac{5}{2}+\tfrac{1}{4}c_1+\tfrac{3}{2}c_2 & \tfrac{1}{2}c_2-c_3 \end{pmatrix} & (k_{18} &= 6) \\
 &= \bigl(-c_1-2c_3\bigr)I+c_1\,A+\bigl(1+c_2\bigr)AB && \\
\rho(T_{19}) &= \begin{pmatrix} 1+2c_1+\tfrac{3}{2}c_2 & -1-c_2 \\ -\tfrac{5}{2}-\tfrac{1}{4}c_1-\tfrac{3}{2}c_2 & \tfrac{1}{2}c_2-c_3 \end{pmatrix} & (k_{19} &= 6) \\
 &= \bigl(2c_1+c_2\bigr)I+\bigl(-1-c_2\bigr)AB && \\
\rho(T_{20}) &= \begin{pmatrix} 1+2c_1+\tfrac{3}{2}c_2 & 1+c_2 \\ -\tfrac{5}{2}-\tfrac{1}{4}c_1-\tfrac{3}{2}c_2 & -\tfrac{1}{2}c_2+c_3 \end{pmatrix} & (k_{20} &= 6) \\
 &= c_1I+\bigl(-c_1-2c_3\bigr)A+\bigl(1+c_2\bigr)B && \\
\rho(T_{21}) &= \begin{pmatrix} -\tfrac{1}{2} & -c_1 \\ \tfrac{3}{2}+\tfrac{11}{4}c_1+\tfrac{1}{4}c_2 & -\tfrac{1}{2}+c_3 \end{pmatrix} & (k_{21} &= 6) \\
 &= \bigl(-1+c_3\bigr)I+\bigl(1+c_2\bigr)A-c_1\,B && \\
\rho(T_{22}) &= \begin{pmatrix} -\tfrac{1}{2} & c_1 \\ \tfrac{3}{2}+\tfrac{11}{4}c_1+\tfrac{1}{4}c_2 & \tfrac{1}{2}-c_3 \end{pmatrix} = c_1I+c_1\,AB & (k_{22} &= 6) \\
\rho(T_{23}) &= \begin{pmatrix} 1-c_1 & 2c_1+2c_2 \\ 0 & 0 \end{pmatrix} = I+A+\bigl(c_1+c_2\bigr)B+\bigl(c_1+c_2\bigr)AB & (k_{23} &= 6) \\
\rho(T_{24}) &= \begin{pmatrix} -c_2-\tfrac{1}{2}c_3 & 1-c_1-2c_2 \\ -\tfrac{7}{4}-\tfrac{1}{2}c_1+\tfrac{1}{4}c_2 & \tfrac{1}{2}+\tfrac{3}{2}c_1-\tfrac{1}{2}c_2 \end{pmatrix} & (k_{24} &= 6) \\
 &= \bigl(-2c_2-c_3\bigr)I+c_2\,A+\bigl(1-c_1-2c_2\bigr)B && \\
\rho(T_{25}) &= \begin{pmatrix} -c_2-\tfrac{1}{2}c_3 & -1+c_1+2c_2 \\ -\tfrac{7}{4}-\tfrac{1}{2}c_1+\tfrac{1}{4}c_2 & -\tfrac{1}{2}-\tfrac{3}{2}c_1+\tfrac{1}{2}c_2 \end{pmatrix} & (k_{25} &= 6) \\
 &= \bigl(-c_1-c_2\bigr)I+c_1\,A+\bigl(-1+c_1+2c_2\bigr)AB && \\
\rho(T_{26}) &= \begin{pmatrix} -\tfrac{1}{2} & c_1 \\ -\tfrac{3}{2}-\tfrac{11}{4}c_1-\tfrac{1}{4}c_2 & -\tfrac{1}{2}+c_3 \end{pmatrix} = c_1\,A+c_1\,B & (k_{26} &= 6) \\
\rho(T_{27}) &= \begin{pmatrix} -1+c_3 & 0 \\ 0 & -3-c_3 \end{pmatrix} = -2I+\bigl(-c_1-c_2\bigr)A & (k_{27} &= 6) \\
\rho(T_{28}) &= \begin{pmatrix} -1+c_3 & 0 \\ 0 & 3+c_3 \end{pmatrix} = \bigl(-c_1-c_2\bigr)I-2\,A & (k_{28} &= 6) \\
\rho(T_{29}) &= \begin{pmatrix} 1-c_1 & 0 \\ -\tfrac{5}{2}-\tfrac{7}{2}c_1-\tfrac{3}{2}c_2 & 0 \end{pmatrix} & (k_{29} &= 6) \\
 &= \bigl(1+c_2\bigr)I+\bigl(-c_1-c_2\bigr)A+\bigl(c_1+c_2\bigr)B+\bigl(-c_1-c_2\bigr)AB && \\
\rho(T_{30}) &= \begin{pmatrix} -c_2-\tfrac{1}{2}c_3 & -1+c_1+2c_2 \\ \tfrac{7}{4}+\tfrac{1}{2}c_1-\tfrac{1}{4}c_2 & \tfrac{1}{2}+\tfrac{3}{2}c_1-\tfrac{1}{2}c_2 \end{pmatrix} & (k_{30} &= 6) \\
 &= c_1I+\bigl(-c_1-c_2\bigr)A+\bigl(-1+c_1+2c_2\bigr)B && \\
\rho(T_{31}) &= \begin{pmatrix} -c_2-\tfrac{1}{2}c_3 & 1-c_1-2c_2 \\ \tfrac{7}{4}+\tfrac{1}{2}c_1-\tfrac{1}{4}c_2 & -\tfrac{1}{2}-\tfrac{3}{2}c_1+\tfrac{1}{2}c_2 \end{pmatrix} & (k_{31} &= 6) \\
 &= c_2I+\bigl(-2c_2-c_3\bigr)A+\bigl(1-c_1-2c_2\bigr)AB && \\
\rho(T_{32}) &= \begin{pmatrix} 1-c_1 & -2c_1-2c_2 \\ 0 & 0 \end{pmatrix} & (k_{32} &= 6) \\
 &= -c_1I-c_1\,A+\bigl(-c_1-c_2\bigr)B+\bigl(-c_1-c_2\bigr)AB && \\
\rho(T_{33}) &= \begin{pmatrix} -\tfrac{1}{2} & -c_1 \\ -\tfrac{3}{2}-\tfrac{11}{4}c_1-\tfrac{1}{4}c_2 & \tfrac{1}{2}-c_3 \end{pmatrix} & (k_{33} &= 6) \\
 &= \bigl(1+c_2\bigr)I+\bigl(-1+c_3\bigr)A-c_1\,AB && \\
\rho(T_{34}) &= \begin{pmatrix} 1-c_1 & 0 \\ \tfrac{5}{2}+\tfrac{7}{2}c_1+\tfrac{3}{2}c_2 & 0 \end{pmatrix} & (k_{34} &= 6) \\
 &= \bigl(-c_1-c_2\bigr)I+\bigl(1+c_2\bigr)A+\bigl(-c_1-c_2\bigr)B+\bigl(c_1+c_2\bigr)AB && .
\end{align*}

\subsubsection*{Block 2: $M_2(\C)$ on $\chi_{5}$ (degree 12)}
\noindent\textbf{Generators.}
\[
A = \rho(T_{2}) = \begin{pmatrix} 1 & 0 \\ 0 & -1 \end{pmatrix}, \qquad
B = \rho(T_{3}) = \begin{pmatrix} -\tfrac{1}{2}+\tfrac{1}{2}c_2 & 1 \\ -\tfrac{3}{2}+\tfrac{1}{4}c_1+\tfrac{5}{4}c_2 & \tfrac{1}{2}c_1-\tfrac{1}{2}c_3 \end{pmatrix} .
\]
\textbf{Complete relations} (monomial basis $I,\;A,\;B,\;AB$ of $M_2(\C)$):
\begin{gather*}
A^2-I = 0 ,\\
B^2+\bigl(-c_1-c_2\bigr)B+\bigl(1-c_1-2c_2\bigr)I = 0 ,\\
BA = \bigl(-1-c_1\bigr)I+\bigl(c_1+c_2\bigr)A-AB ,\\
ABA = \bigl(c_1+c_2\bigr)I+\bigl(-1-c_1\bigr)A-B ,\\
ABB = \bigl(-1+c_1+2c_2\bigr)A+\bigl(c_1+c_2\bigr)AB .
\end{gather*}
\textbf{All orbital images} ($\rho(T_i)$, valency $k_i$, word in the generators):
\begin{align*}
\rho(T_{1}) &= \begin{pmatrix} 1 & 0 \\ 0 & 1 \end{pmatrix} = I & (k_{1} &= 1) \\
\rho(T_{2}) &= \begin{pmatrix} 1 & 0 \\ 0 & -1 \end{pmatrix} = A & (k_{2} &= 1) \\
\rho(T_{3}) &= \begin{pmatrix} -\tfrac{1}{2}+\tfrac{1}{2}c_2 & 1 \\ -\tfrac{3}{2}+\tfrac{1}{4}c_1+\tfrac{5}{4}c_2 & \tfrac{1}{2}c_1-\tfrac{1}{2}c_3 \end{pmatrix} = B & (k_{3} &= 3) \\
\rho(T_{4}) &= \begin{pmatrix} -\tfrac{1}{2}+\tfrac{1}{2}c_2 & -1 \\ -\tfrac{3}{2}+\tfrac{1}{4}c_1+\tfrac{5}{4}c_2 & -\tfrac{1}{2}c_1+\tfrac{1}{2}c_3 \end{pmatrix} & (k_{4} &= 3) \\
 &= \bigl(-1-c_1\bigr)I+\bigl(c_1+c_2\bigr)A-AB && \\
\rho(T_{5}) &= \begin{pmatrix} -\tfrac{1}{2}+\tfrac{1}{2}c_2 & -1 \\ \tfrac{3}{2}-\tfrac{1}{4}c_1-\tfrac{5}{4}c_2 & \tfrac{1}{2}c_1-\tfrac{1}{2}c_3 \end{pmatrix} & (k_{5} &= 3) \\
 &= \bigl(c_1+c_2\bigr)I+\bigl(-1-c_1\bigr)A-B && \\
\rho(T_{6}) &= \begin{pmatrix} -\tfrac{1}{2}+\tfrac{1}{2}c_2 & 1 \\ \tfrac{3}{2}-\tfrac{1}{4}c_1-\tfrac{5}{4}c_2 & -\tfrac{1}{2}c_1+\tfrac{1}{2}c_3 \end{pmatrix} = AB & (k_{6} &= 3) \\
\rho(T_{7}) &= \begin{pmatrix} c_1-\tfrac{1}{2}c_3 & c_1 \\ \tfrac{3}{4}+\tfrac{3}{2}c_1+c_2 & c_2+\tfrac{1}{2}c_3 \end{pmatrix} = -c_3I+c_1\,A+c_1\,AB & (k_{7} &= 6) \\
\rho(T_{8}) &= \begin{pmatrix} -2c_1 & 0 \\ 0 & 0 \end{pmatrix} = -c_1I-c_1\,A & (k_{8} &= 6) \\
\rho(T_{9}) &= \begin{pmatrix} c_1-\tfrac{1}{2}c_3 & -c_1 \\ \tfrac{3}{4}+\tfrac{3}{2}c_1+c_2 & -c_2-\tfrac{1}{2}c_3 \end{pmatrix} = \bigl(1+c_1\bigr)I-A-c_1\,B & (k_{9} &= 6) \\
\rho(T_{10}) &= \begin{pmatrix} -\tfrac{1}{2} & -2c_1-c_3 \\ \tfrac{1}{4}+\tfrac{1}{2}c_1-\tfrac{3}{4}c_2 & -\tfrac{1}{2}-c_2 \end{pmatrix} = -A+\bigl(-2c_1-c_3\bigr)B & (k_{10} &= 6) \\
\rho(T_{11}) &= \begin{pmatrix} -\tfrac{1}{2} & 2c_1+c_3 \\ \tfrac{1}{4}+\tfrac{1}{2}c_1-\tfrac{3}{4}c_2 & \tfrac{1}{2}+c_2 \end{pmatrix} & (k_{11} &= 6) \\
 &= \bigl(1+c_2\bigr)I+\bigl(-1-c_2\bigr)A+\bigl(2c_1+c_3\bigr)AB && \\
\rho(T_{12}) &= \begin{pmatrix} -2c_1 & 0 \\ 0 & 0 \end{pmatrix} = -c_1I-c_1\,A & (k_{12} &= 6) \\
\rho(T_{13}) &= \begin{pmatrix} c_1-\tfrac{1}{2}c_3 & c_1 \\ -\tfrac{3}{4}-\tfrac{3}{2}c_1-c_2 & -c_2-\tfrac{1}{2}c_3 \end{pmatrix} = c_1I-c_3\,A+c_1\,B & (k_{13} &= 6) \\
\rho(T_{14}) &= \begin{pmatrix} c_1-\tfrac{1}{2}c_3 & -c_1 \\ -\tfrac{3}{4}-\tfrac{3}{2}c_1-c_2 & c_2+\tfrac{1}{2}c_3 \end{pmatrix} = -I+\bigl(1+c_1\bigr)A-c_1\,AB & (k_{14} &= 6) \\
\rho(T_{15}) &= \begin{pmatrix} -\tfrac{1}{2} & -2c_1-c_3 \\ -\tfrac{1}{4}-\tfrac{1}{2}c_1+\tfrac{3}{4}c_2 & \tfrac{1}{2}+c_2 \end{pmatrix} = -I+\bigl(-2c_1-c_3\bigr)AB & (k_{15} &= 6) \\
\rho(T_{16}) &= \begin{pmatrix} -\tfrac{1}{2} & 2c_1+c_3 \\ -\tfrac{1}{4}-\tfrac{1}{2}c_1+\tfrac{3}{4}c_2 & -\tfrac{1}{2}-c_2 \end{pmatrix} & (k_{16} &= 6) \\
 &= \bigl(-1-c_2\bigr)I+\bigl(1+c_2\bigr)A+\bigl(2c_1+c_3\bigr)B && \\
\rho(T_{17}) &= \begin{pmatrix} -1-\tfrac{1}{2}c_1-2c_2 & -1-c_1 \\ \tfrac{9}{4}+\tfrac{5}{4}c_1-\tfrac{1}{4}c_2 & -\tfrac{1}{2}c_1+c_2 \end{pmatrix} & (k_{17} &= 6) \\
 &= \bigl(c_1+2c_3\bigr)A+\bigl(-1-c_1\bigr)B && \\
\rho(T_{18}) &= \begin{pmatrix} -1-\tfrac{1}{2}c_1-2c_2 & 1+c_1 \\ \tfrac{9}{4}+\tfrac{5}{4}c_1-\tfrac{1}{4}c_2 & \tfrac{1}{2}c_1-c_2 \end{pmatrix} & (k_{18} &= 6) \\
 &= \bigl(-2c_2-c_3\bigr)I+c_3\,A+\bigl(1+c_1\bigr)AB && \\
\rho(T_{19}) &= \begin{pmatrix} -1-\tfrac{1}{2}c_1-2c_2 & -1-c_1 \\ -\tfrac{9}{4}-\tfrac{5}{4}c_1+\tfrac{1}{4}c_2 & \tfrac{1}{2}c_1-c_2 \end{pmatrix} & (k_{19} &= 6) \\
 &= \bigl(c_1+2c_3\bigr)I+\bigl(-1-c_1\bigr)AB && \\
\rho(T_{20}) &= \begin{pmatrix} -1-\tfrac{1}{2}c_1-2c_2 & 1+c_1 \\ -\tfrac{9}{4}-\tfrac{5}{4}c_1+\tfrac{1}{4}c_2 & -\tfrac{1}{2}c_1+c_2 \end{pmatrix} & (k_{20} &= 6) \\
 &= c_3I+\bigl(-2c_2-c_3\bigr)A+\bigl(1+c_1\bigr)B && \\
\rho(T_{21}) &= \begin{pmatrix} -\tfrac{1}{2} & -c_3 \\ -\tfrac{5}{4}-\tfrac{5}{2}c_1-\tfrac{11}{4}c_2 & -\tfrac{1}{2}+c_2 \end{pmatrix} & (k_{21} &= 6) \\
 &= \bigl(-1+c_2\bigr)I+\bigl(1+c_1\bigr)A-c_3\,B && \\
\rho(T_{22}) &= \begin{pmatrix} -\tfrac{1}{2} & c_3 \\ -\tfrac{5}{4}-\tfrac{5}{2}c_1-\tfrac{11}{4}c_2 & \tfrac{1}{2}-c_2 \end{pmatrix} = c_3I+c_3\,AB & (k_{22} &= 6) \\
\rho(T_{23}) &= \begin{pmatrix} 1-c_3 & -2-2c_2 \\ 0 & 0 \end{pmatrix} = I+A+\bigl(-1-c_2\bigr)B+\bigl(-1-c_2\bigr)AB & (k_{23} &= 6) \\
\rho(T_{24}) &= \begin{pmatrix} -c_1-\tfrac{1}{2}c_2 & 2-c_1+c_2 \\ -\tfrac{5}{4}+\tfrac{3}{4}c_1+\tfrac{1}{2}c_2 & -1-2c_1-\tfrac{3}{2}c_2 \end{pmatrix} & (k_{24} &= 6) \\
 &= \bigl(-2c_1-c_2\bigr)I+c_1\,A+\bigl(2-c_1+c_2\bigr)B && \\
\rho(T_{25}) &= \begin{pmatrix} -c_1-\tfrac{1}{2}c_2 & -2+c_1-c_2 \\ -\tfrac{5}{4}+\tfrac{3}{4}c_1+\tfrac{1}{2}c_2 & 1+2c_1+\tfrac{3}{2}c_2 \end{pmatrix} & (k_{25} &= 6) \\
 &= \bigl(1+c_2\bigr)I+c_3\,A+\bigl(-2+c_1-c_2\bigr)AB && \\
\rho(T_{26}) &= \begin{pmatrix} -\tfrac{1}{2} & c_3 \\ \tfrac{5}{4}+\tfrac{5}{2}c_1+\tfrac{11}{4}c_2 & -\tfrac{1}{2}+c_2 \end{pmatrix} = c_3\,A+c_3\,B & (k_{26} &= 6) \\
\rho(T_{27}) &= \begin{pmatrix} -1+c_2 & 0 \\ 0 & -3-c_2 \end{pmatrix} = -2I+\bigl(1+c_2\bigr)A & (k_{27} &= 6) \\
\rho(T_{28}) &= \begin{pmatrix} -1+c_2 & 0 \\ 0 & 3+c_2 \end{pmatrix} = \bigl(1+c_2\bigr)I-2\,A & (k_{28} &= 6) \\
\rho(T_{29}) &= \begin{pmatrix} 1-c_3 & 0 \\ 1+2c_1+\tfrac{7}{2}c_2 & 0 \end{pmatrix} & (k_{29} &= 6) \\
 &= \bigl(1+c_1\bigr)I+\bigl(1+c_2\bigr)A+\bigl(-1-c_2\bigr)B+\bigl(1+c_2\bigr)AB && \\
\rho(T_{30}) &= \begin{pmatrix} -c_1-\tfrac{1}{2}c_2 & -2+c_1-c_2 \\ \tfrac{5}{4}-\tfrac{3}{4}c_1-\tfrac{1}{2}c_2 & -1-2c_1-\tfrac{3}{2}c_2 \end{pmatrix} & (k_{30} &= 6) \\
 &= c_3I+\bigl(1+c_2\bigr)A+\bigl(-2+c_1-c_2\bigr)B && \\
\rho(T_{31}) &= \begin{pmatrix} -c_1-\tfrac{1}{2}c_2 & 2-c_1+c_2 \\ \tfrac{5}{4}-\tfrac{3}{4}c_1-\tfrac{1}{2}c_2 & 1+2c_1+\tfrac{3}{2}c_2 \end{pmatrix} & (k_{31} &= 6) \\
 &= c_1I+\bigl(-2c_1-c_2\bigr)A+\bigl(2-c_1+c_2\bigr)AB && \\
\rho(T_{32}) &= \begin{pmatrix} 1-c_3 & 2+2c_2 \\ 0 & 0 \end{pmatrix} = -c_3I-c_3\,A+\bigl(1+c_2\bigr)B+\bigl(1+c_2\bigr)AB & (k_{32} &= 6) \\
\rho(T_{33}) &= \begin{pmatrix} -\tfrac{1}{2} & -c_3 \\ \tfrac{5}{4}+\tfrac{5}{2}c_1+\tfrac{11}{4}c_2 & \tfrac{1}{2}-c_2 \end{pmatrix} & (k_{33} &= 6) \\
 &= \bigl(1+c_1\bigr)I+\bigl(-1+c_2\bigr)A-c_3\,AB && \\
\rho(T_{34}) &= \begin{pmatrix} 1-c_3 & 0 \\ -1-2c_1-\tfrac{7}{2}c_2 & 0 \end{pmatrix} & (k_{34} &= 6) \\
 &= \bigl(1+c_2\bigr)I+\bigl(1+c_1\bigr)A+\bigl(1+c_2\bigr)B+\bigl(-1-c_2\bigr)AB && .
\end{align*}

\subsubsection*{Block 3: $M_2(\C)$ on $\chi_{6}$ (degree 12)}
\noindent\textbf{Generators.}
\[
A = \rho(T_{2}) = \begin{pmatrix} 1 & 0 \\ 0 & -1 \end{pmatrix}, \qquad
B = \rho(T_{3}) = \begin{pmatrix} -\tfrac{1}{2}+\tfrac{1}{2}c_1 & 1 \\ -\tfrac{7}{4}+c_1-\tfrac{1}{4}c_2 & -\tfrac{1}{2}c_2+\tfrac{1}{2}c_3 \end{pmatrix} .
\]
\textbf{Complete relations} (monomial basis $I,\;A,\;B,\;AB$ of $M_2(\C)$):
\begin{gather*}
A^2-I = 0 ,\\
B^2+\bigl(1+c_2\bigr)B+\bigl(2-c_1+c_2\bigr)I = 0 ,\\
BA = \bigl(c_1+c_2\bigr)I+\bigl(-1-c_2\bigr)A-AB ,\\
ABA = \bigl(-1-c_2\bigr)I+\bigl(c_1+c_2\bigr)A-B ,\\
ABB = \bigl(-2+c_1-c_2\bigr)A+\bigl(-1-c_2\bigr)AB .
\end{gather*}
\textbf{All orbital images} ($\rho(T_i)$, valency $k_i$, word in the generators):
\begin{align*}
\rho(T_{1}) &= \begin{pmatrix} 1 & 0 \\ 0 & 1 \end{pmatrix} = I & (k_{1} &= 1) \\
\rho(T_{2}) &= \begin{pmatrix} 1 & 0 \\ 0 & -1 \end{pmatrix} = A & (k_{2} &= 1) \\
\rho(T_{3}) &= \begin{pmatrix} -\tfrac{1}{2}+\tfrac{1}{2}c_1 & 1 \\ -\tfrac{7}{4}+c_1-\tfrac{1}{4}c_2 & -\tfrac{1}{2}c_2+\tfrac{1}{2}c_3 \end{pmatrix} = B & (k_{3} &= 3) \\
\rho(T_{4}) &= \begin{pmatrix} -\tfrac{1}{2}+\tfrac{1}{2}c_1 & -1 \\ -\tfrac{7}{4}+c_1-\tfrac{1}{4}c_2 & \tfrac{1}{2}c_2-\tfrac{1}{2}c_3 \end{pmatrix} & (k_{4} &= 3) \\
 &= \bigl(c_1+c_2\bigr)I+\bigl(-1-c_2\bigr)A-AB && \\
\rho(T_{5}) &= \begin{pmatrix} -\tfrac{1}{2}+\tfrac{1}{2}c_1 & -1 \\ \tfrac{7}{4}-c_1+\tfrac{1}{4}c_2 & -\tfrac{1}{2}c_2+\tfrac{1}{2}c_3 \end{pmatrix} & (k_{5} &= 3) \\
 &= \bigl(-1-c_2\bigr)I+\bigl(c_1+c_2\bigr)A-B && \\
\rho(T_{6}) &= \begin{pmatrix} -\tfrac{1}{2}+\tfrac{1}{2}c_1 & 1 \\ \tfrac{7}{4}-c_1+\tfrac{1}{4}c_2 & \tfrac{1}{2}c_2-\tfrac{1}{2}c_3 \end{pmatrix} = AB & (k_{6} &= 3) \\
\rho(T_{7}) &= \begin{pmatrix} -\tfrac{1}{2}c_2+c_3 & c_3 \\ -\tfrac{3}{4}-\tfrac{1}{2}c_1-\tfrac{3}{2}c_2 & c_1+\tfrac{1}{2}c_2 \end{pmatrix} = -c_2I+c_3\,A+c_3\,AB & (k_{7} &= 6) \\
\rho(T_{8}) &= \begin{pmatrix} -2c_3 & 0 \\ 0 & 0 \end{pmatrix} = -c_3I-c_3\,A & (k_{8} &= 6) \\
\rho(T_{9}) &= \begin{pmatrix} -\tfrac{1}{2}c_2+c_3 & -c_3 \\ -\tfrac{3}{4}-\tfrac{1}{2}c_1-\tfrac{3}{2}c_2 & -c_1-\tfrac{1}{2}c_2 \end{pmatrix} & (k_{9} &= 6) \\
 &= \bigl(-c_1-c_2\bigr)I-A-c_3\,B && \\
\rho(T_{10}) &= \begin{pmatrix} -\tfrac{1}{2} & -c_2-2c_3 \\ -\tfrac{1}{4}-\tfrac{5}{4}c_1-\tfrac{1}{2}c_2 & -\tfrac{1}{2}-c_1 \end{pmatrix} = -A+\bigl(-c_2-2c_3\bigr)B & (k_{10} &= 6) \\
\rho(T_{11}) &= \begin{pmatrix} -\tfrac{1}{2} & c_2+2c_3 \\ -\tfrac{1}{4}-\tfrac{5}{4}c_1-\tfrac{1}{2}c_2 & \tfrac{1}{2}+c_1 \end{pmatrix} & (k_{11} &= 6) \\
 &= \bigl(1+c_1\bigr)I+\bigl(-1-c_1\bigr)A+\bigl(c_2+2c_3\bigr)AB && \\
\rho(T_{12}) &= \begin{pmatrix} -2c_3 & 0 \\ 0 & 0 \end{pmatrix} = -c_3I-c_3\,A & (k_{12} &= 6) \\
\rho(T_{13}) &= \begin{pmatrix} -\tfrac{1}{2}c_2+c_3 & c_3 \\ \tfrac{3}{4}+\tfrac{1}{2}c_1+\tfrac{3}{2}c_2 & -c_1-\tfrac{1}{2}c_2 \end{pmatrix} = c_3I-c_2\,A+c_3\,B & (k_{13} &= 6) \\
\rho(T_{14}) &= \begin{pmatrix} -\tfrac{1}{2}c_2+c_3 & -c_3 \\ \tfrac{3}{4}+\tfrac{1}{2}c_1+\tfrac{3}{2}c_2 & c_1+\tfrac{1}{2}c_2 \end{pmatrix} & (k_{14} &= 6) \\
 &= -I+\bigl(-c_1-c_2\bigr)A-c_3\,AB && \\
\rho(T_{15}) &= \begin{pmatrix} -\tfrac{1}{2} & -c_2-2c_3 \\ \tfrac{1}{4}+\tfrac{5}{4}c_1+\tfrac{1}{2}c_2 & \tfrac{1}{2}+c_1 \end{pmatrix} = -I+\bigl(-c_2-2c_3\bigr)AB & (k_{15} &= 6) \\
\rho(T_{16}) &= \begin{pmatrix} -\tfrac{1}{2} & c_2+2c_3 \\ \tfrac{1}{4}+\tfrac{5}{4}c_1+\tfrac{1}{2}c_2 & -\tfrac{1}{2}-c_1 \end{pmatrix} & (k_{16} &= 6) \\
 &= \bigl(-1-c_1\bigr)I+\bigl(1+c_1\bigr)A+\bigl(c_2+2c_3\bigr)B && \\
\rho(T_{17}) &= \begin{pmatrix} -\tfrac{1}{2}-\tfrac{3}{2}c_1+\tfrac{1}{2}c_2 & c_1+c_2 \\ 1-\tfrac{3}{2}c_1-\tfrac{5}{4}c_2 & c_1-\tfrac{1}{2}c_3 \end{pmatrix} & (k_{17} &= 6) \\
 &= \bigl(2c_2+c_3\bigr)A+\bigl(c_1+c_2\bigr)B && \\
\rho(T_{18}) &= \begin{pmatrix} -\tfrac{1}{2}-\tfrac{3}{2}c_1+\tfrac{1}{2}c_2 & -c_1-c_2 \\ 1-\tfrac{3}{2}c_1-\tfrac{5}{4}c_2 & -c_1+\tfrac{1}{2}c_3 \end{pmatrix} & (k_{18} &= 6) \\
 &= \bigl(-2c_1-c_2\bigr)I+c_2\,A+\bigl(-c_1-c_2\bigr)AB && \\
\rho(T_{19}) &= \begin{pmatrix} -\tfrac{1}{2}-\tfrac{3}{2}c_1+\tfrac{1}{2}c_2 & c_1+c_2 \\ -1+\tfrac{3}{2}c_1+\tfrac{5}{4}c_2 & -c_1+\tfrac{1}{2}c_3 \end{pmatrix} & (k_{19} &= 6) \\
 &= \bigl(2c_2+c_3\bigr)I+\bigl(c_1+c_2\bigr)AB && \\
\rho(T_{20}) &= \begin{pmatrix} -\tfrac{1}{2}-\tfrac{3}{2}c_1+\tfrac{1}{2}c_2 & -c_1-c_2 \\ -1+\tfrac{3}{2}c_1+\tfrac{5}{4}c_2 & c_1-\tfrac{1}{2}c_3 \end{pmatrix} & (k_{20} &= 6) \\
 &= c_2I+\bigl(-2c_1-c_2\bigr)A+\bigl(-c_1-c_2\bigr)B && \\
\rho(T_{21}) &= \begin{pmatrix} -\tfrac{1}{2} & -c_2 \\ \tfrac{5}{4}-\tfrac{1}{4}c_1+\tfrac{5}{2}c_2 & -\tfrac{1}{2}+c_1 \end{pmatrix} & (k_{21} &= 6) \\
 &= \bigl(-1+c_1\bigr)I+\bigl(-c_1-c_2\bigr)A-c_2\,B && \\
\rho(T_{22}) &= \begin{pmatrix} -\tfrac{1}{2} & c_2 \\ \tfrac{5}{4}-\tfrac{1}{4}c_1+\tfrac{5}{2}c_2 & \tfrac{1}{2}-c_1 \end{pmatrix} = c_2I+c_2\,AB & (k_{22} &= 6) \\
\rho(T_{23}) &= \begin{pmatrix} 1-c_2 & -2-2c_1 \\ 0 & 0 \end{pmatrix} = I+A+\bigl(-1-c_1\bigr)B+\bigl(-1-c_1\bigr)AB & (k_{23} &= 6) \\
\rho(T_{24}) &= \begin{pmatrix} -\tfrac{1}{2}c_1-c_3 & 3+2c_1+c_2 \\ -2-\tfrac{1}{4}c_1-\tfrac{3}{4}c_2 & 1+\tfrac{1}{2}c_1+2c_2 \end{pmatrix} & (k_{24} &= 6) \\
 &= \bigl(-c_1-2c_3\bigr)I+c_3\,A+\bigl(3+2c_1+c_2\bigr)B && \\
\rho(T_{25}) &= \begin{pmatrix} -\tfrac{1}{2}c_1-c_3 & -3-2c_1-c_2 \\ -2-\tfrac{1}{4}c_1-\tfrac{3}{4}c_2 & -1-\tfrac{1}{2}c_1-2c_2 \end{pmatrix} & (k_{25} &= 6) \\
 &= \bigl(1+c_1\bigr)I+c_2\,A+\bigl(-3-2c_1-c_2\bigr)AB && \\
\rho(T_{26}) &= \begin{pmatrix} -\tfrac{1}{2} & c_2 \\ -\tfrac{5}{4}+\tfrac{1}{4}c_1-\tfrac{5}{2}c_2 & -\tfrac{1}{2}+c_1 \end{pmatrix} = c_2\,A+c_2\,B & (k_{26} &= 6) \\
\rho(T_{27}) &= \begin{pmatrix} -1+c_1 & 0 \\ 0 & -3-c_1 \end{pmatrix} = -2I+\bigl(1+c_1\bigr)A & (k_{27} &= 6) \\
\rho(T_{28}) &= \begin{pmatrix} -1+c_1 & 0 \\ 0 & 3+c_1 \end{pmatrix} = \bigl(1+c_1\bigr)I-2\,A & (k_{28} &= 6) \\
\rho(T_{29}) &= \begin{pmatrix} 1-c_2 & 0 \\ -1+\tfrac{3}{2}c_1-2c_2 & 0 \end{pmatrix} & (k_{29} &= 6) \\
 &= \bigl(-c_1-c_2\bigr)I+\bigl(1+c_1\bigr)A+\bigl(-1-c_1\bigr)B+\bigl(1+c_1\bigr)AB && \\
\rho(T_{30}) &= \begin{pmatrix} -\tfrac{1}{2}c_1-c_3 & -3-2c_1-c_2 \\ 2+\tfrac{1}{4}c_1+\tfrac{3}{4}c_2 & 1+\tfrac{1}{2}c_1+2c_2 \end{pmatrix} & (k_{30} &= 6) \\
 &= c_2I+\bigl(1+c_1\bigr)A+\bigl(-3-2c_1-c_2\bigr)B && \\
\rho(T_{31}) &= \begin{pmatrix} -\tfrac{1}{2}c_1-c_3 & 3+2c_1+c_2 \\ 2+\tfrac{1}{4}c_1+\tfrac{3}{4}c_2 & -1-\tfrac{1}{2}c_1-2c_2 \end{pmatrix} & (k_{31} &= 6) \\
 &= c_3I+\bigl(-c_1-2c_3\bigr)A+\bigl(3+2c_1+c_2\bigr)AB && \\
\rho(T_{32}) &= \begin{pmatrix} 1-c_2 & 2+2c_1 \\ 0 & 0 \end{pmatrix} = -c_2I-c_2\,A+\bigl(1+c_1\bigr)B+\bigl(1+c_1\bigr)AB & (k_{32} &= 6) \\
\rho(T_{33}) &= \begin{pmatrix} -\tfrac{1}{2} & -c_2 \\ -\tfrac{5}{4}+\tfrac{1}{4}c_1-\tfrac{5}{2}c_2 & \tfrac{1}{2}-c_1 \end{pmatrix} & (k_{33} &= 6) \\
 &= \bigl(-c_1-c_2\bigr)I+\bigl(-1+c_1\bigr)A-c_2\,AB && \\
\rho(T_{34}) &= \begin{pmatrix} 1-c_2 & 0 \\ 1-\tfrac{3}{2}c_1+2c_2 & 0 \end{pmatrix} & (k_{34} &= 6) \\
 &= \bigl(1+c_1\bigr)I+\bigl(-c_1-c_2\bigr)A+\bigl(1+c_1\bigr)B+\bigl(-1-c_1\bigr)AB && .
\end{align*}

\subsubsection*{Block 4: $M_3(\C)$ on $\chi_{7}$ (degree 13)}
\noindent\textbf{Generators.}
\[
A = \rho(T_{2}) = \begin{pmatrix} 1 & 0 & 0 \\ 0 & 1 & 0 \\ 0 & 0 & -1 \end{pmatrix}, \qquad
B = \rho(T_{3}) = \begin{pmatrix} 2 & -1 & -2 \\ 2 & -2 & -3 \\ 1 & 0 & 0 \end{pmatrix} .
\]
\textbf{Complete relations} (monomial basis $I,\;A,\;B,\;AB,\;BA,\;BB,\;ABB,\;BAB,\;BABB$ of $M_3(\C)$):
\begin{gather*}
A^2-I = 0 ,\\
B^3+I = 0 ,\\
ABA = -B+AB+BA ,\\
BAA = B ,\\
BBA = 4I+BB-ABB-BAB ,\\
BBB = -I ,\\
ABBA = 4\,A-BAB ,\\
ABBB = -A ,\\
BABA = 4I-ABB ,\\
BABBA = I-A-4\,B+4\,BA+BABB ,\\
BABBB = -BA .
\end{gather*}
\textbf{All orbital images} ($\rho(T_i)$, valency $k_i$, word in the generators):
\begin{align*}
\rho(T_{1}) &= \begin{pmatrix} 1 & 0 & 0 \\ 0 & 1 & 0 \\ 0 & 0 & 1 \end{pmatrix} = I & (k_{1} &= 1) \\
\rho(T_{2}) &= \begin{pmatrix} 1 & 0 & 0 \\ 0 & 1 & 0 \\ 0 & 0 & -1 \end{pmatrix} = A & (k_{2} &= 1) \\
\rho(T_{3}) &= \begin{pmatrix} 2 & -1 & -2 \\ 2 & -2 & -3 \\ 1 & 0 & 0 \end{pmatrix} = B & (k_{3} &= 3) \\
\rho(T_{4}) &= \begin{pmatrix} 2 & -1 & 2 \\ 2 & -2 & 3 \\ 1 & 0 & 0 \end{pmatrix} = BA & (k_{4} &= 3) \\
\rho(T_{5}) &= \begin{pmatrix} 2 & -1 & 2 \\ 2 & -2 & 3 \\ -1 & 0 & 0 \end{pmatrix} = -B+AB+BA & (k_{5} &= 3) \\
\rho(T_{6}) &= \begin{pmatrix} 2 & -1 & -2 \\ 2 & -2 & -3 \\ -1 & 0 & 0 \end{pmatrix} = AB & (k_{6} &= 3) \\
\rho(T_{7}) &= \begin{pmatrix} 1 & 0 & 3 \\ 3 & -1 & 8 \\ 0 & 1 & 3 \end{pmatrix} = -I-2\,A-AB+BA-BB+ABB+BAB & (k_{7} &= 6) \\
\rho(T_{8}) &= \begin{pmatrix} 6 & -2 & 0 \\ 14 & -4 & 0 \\ 0 & 0 & 0 \end{pmatrix} = -A-2\,B-AB+BA-BB+BAB+BABB & (k_{8} &= 6) \\
\rho(T_{9}) &= \begin{pmatrix} 1 & 0 & -3 \\ 3 & -1 & -8 \\ 0 & 1 & -3 \end{pmatrix} = -2I+3\,A+2\,B-AB-BA-BB & (k_{9} &= 6) \\
\rho(T_{10}) &= \begin{pmatrix} 1 & -1 & -2 \\ -4 & 1 & -3 \\ 1 & 0 & 1 \end{pmatrix} = 2I+A+B+BB-BAB & (k_{10} &= 6) \\
\rho(T_{11}) &= \begin{pmatrix} 1 & -1 & 2 \\ -4 & 1 & 3 \\ 1 & 0 & -1 \end{pmatrix} = I+2\,A+BA+BB-BAB & (k_{11} &= 6) \\
\rho(T_{12}) &= \begin{pmatrix} 0 & 2 & 0 \\ -2 & 2 & 0 \\ 0 & 0 & 0 \end{pmatrix} = -I+2\,B+AB-BA-BB+BAB-BABB & (k_{12} &= 6) \\
\rho(T_{13}) &= \begin{pmatrix} 1 & 0 & 3 \\ 3 & -1 & 8 \\ 0 & -1 & -3 \end{pmatrix} = -2I-A-2\,B+AB+BA+BAB & (k_{13} &= 6) \\
\rho(T_{14}) &= \begin{pmatrix} 1 & 0 & -3 \\ 3 & -1 & -8 \\ 0 & -1 & 3 \end{pmatrix} = 3I-2\,A+AB-BA-ABB & (k_{14} &= 6) \\
\rho(T_{15}) &= \begin{pmatrix} 1 & -1 & -2 \\ -4 & 1 & -3 \\ -1 & 0 & -1 \end{pmatrix} = I+2\,A+AB+BB-BAB & (k_{15} &= 6) \\
\rho(T_{16}) &= \begin{pmatrix} 1 & -1 & 2 \\ -4 & 1 & 3 \\ -1 & 0 & 1 \end{pmatrix} = 2I+A-B+AB+BA+BB-BAB & (k_{16} &= 6) \\
\rho(T_{17}) &= \begin{pmatrix} 1 & 0 & -1 \\ 3 & -1 & 2 \\ 2 & -1 & 1 \end{pmatrix} = -3\,A+BAB & (k_{17} &= 6) \\
\rho(T_{18}) &= \begin{pmatrix} 1 & 0 & 1 \\ 3 & -1 & -2 \\ 2 & -1 & -1 \end{pmatrix} = I-ABB & (k_{18} &= 6) \\
\rho(T_{19}) &= \begin{pmatrix} 1 & 0 & -1 \\ 3 & -1 & 2 \\ -2 & 1 & -1 \end{pmatrix} = -3I-BB+ABB+BAB & (k_{19} &= 6) \\
\rho(T_{20}) &= \begin{pmatrix} 1 & 0 & 1 \\ 3 & -1 & -2 \\ -2 & 1 & 1 \end{pmatrix} = A-BB & (k_{20} &= 6) \\
\rho(T_{21}) &= \begin{pmatrix} -2 & 1 & -3 \\ -5 & 4 & -1 \\ 3 & -1 & -2 \end{pmatrix} = B-AB-BA+BB & (k_{21} &= 6) \\
\rho(T_{22}) &= \begin{pmatrix} -2 & 1 & 3 \\ -5 & 4 & 1 \\ 3 & -1 & 2 \end{pmatrix} = 4I-AB+BB-ABB-BAB & (k_{22} &= 6) \\
\rho(T_{23}) &= \begin{pmatrix} 2 & -2 & -2 \\ 8 & -4 & 4 \\ 0 & 0 & 0 \end{pmatrix} & (k_{23} &= 6) \\
 &= -2I-3\,A-2\,B-AB+BA+ABB+BAB+BABB && \\
\rho(T_{24}) &= \begin{pmatrix} -3 & 0 & 3 \\ -3 & -1 & 8 \\ 0 & -1 & -1 \end{pmatrix} = -I-2\,A-2\,B+AB+BA+BB & (k_{24} &= 6) \\
\rho(T_{25}) &= \begin{pmatrix} -3 & 0 & -3 \\ -3 & -1 & -8 \\ 0 & -1 & 1 \end{pmatrix} = 2I-A+AB-BA+BB-ABB-BAB & (k_{25} &= 6) \\
\rho(T_{26}) &= \begin{pmatrix} -2 & 1 & 3 \\ -5 & 4 & 1 \\ -3 & 1 & -2 \end{pmatrix} = 4\,A-B-BAB & (k_{26} &= 6) \\
\rho(T_{27}) &= \begin{pmatrix} -2 & 2 & 0 \\ 2 & 2 & 0 \\ 0 & 0 & 2 \end{pmatrix} = -2\,A-AB-BA-BB+BAB & (k_{27} &= 6) \\
\rho(T_{28}) &= \begin{pmatrix} -2 & 2 & 0 \\ 2 & 2 & 0 \\ 0 & 0 & -2 \end{pmatrix} = -2I-AB-BA-BB+BAB & (k_{28} &= 6) \\
\rho(T_{29}) &= \begin{pmatrix} -4 & 2 & 0 \\ -8 & 2 & 0 \\ 4 & -2 & 0 \end{pmatrix} = I-2\,A+2\,B+AB-BA+BB-ABB-BABB & (k_{29} &= 6) \\
\rho(T_{30}) &= \begin{pmatrix} -3 & 0 & -3 \\ -3 & -1 & -8 \\ 0 & 1 & -1 \end{pmatrix} = -I+2\,A+2\,B-AB-BA-BAB & (k_{30} &= 6) \\
\rho(T_{31}) &= \begin{pmatrix} -3 & 0 & 3 \\ -3 & -1 & 8 \\ 0 & 1 & 1 \end{pmatrix} = -2I-A-AB+BA+ABB & (k_{31} &= 6) \\
\rho(T_{32}) &= \begin{pmatrix} 2 & -2 & 2 \\ 8 & -4 & -4 \\ 0 & 0 & 0 \end{pmatrix} = 2I+A-2\,B-AB+BA-ABB-BAB+BABB & (k_{32} &= 6) \\
\rho(T_{33}) &= \begin{pmatrix} -2 & 1 & -3 \\ -5 & 4 & -1 \\ -3 & 1 & 2 \end{pmatrix} = -BA+ABB & (k_{33} &= 6) \\
\rho(T_{34}) &= \begin{pmatrix} -4 & 2 & 0 \\ -8 & 2 & 0 \\ -4 & 2 & 0 \end{pmatrix} = -3I+2\,A+2\,B+AB-BA-BB+ABB-BABB & (k_{34} &= 6) .
\end{align*}

\subsubsection*{Block 5: $M_3(\C)$ on $\chi_{9}$ (degree 14)}
\noindent\textbf{Generators.}
\[
A = \rho(T_{2}) = \begin{pmatrix} 1 & 0 & 0 \\ 0 & 1 & 0 \\ 0 & 0 & -1 \end{pmatrix}, \qquad
B = \rho(T_{3}) = \begin{pmatrix} \tfrac{1}{2} & \tfrac{1}{2} & -\tfrac{11}{4} \\ 2 & 1 & 3 \\ 1 & 0 & -\tfrac{3}{2} \end{pmatrix} .
\]
\textbf{Complete relations} (monomial basis $I,\;A,\;B,\;AB,\;BA,\;BB,\;ABB,\;BAB,\;BABB$ of $M_3(\C)$):
\begin{gather*}
A^2-I = 0 ,\\
B^3-5I = 0 ,\\
ABA = -3I+3\,A-B+AB+BA ,\\
BAA = B ,\\
BBA = I+3\,B+BB-ABB-BAB ,\\
BBB = 5I ,\\
ABBA = A+3\,B-BAB ,\\
ABBB = 5\,A ,\\
BABA = I+3\,BA-ABB ,\\
BABBA = -5I+5\,A-B+BA+BABB ,\\
BABBB = 5\,BA .
\end{gather*}
\textbf{All orbital images} ($\rho(T_i)$, valency $k_i$, word in the generators):
\begin{align*}
\rho(T_{1}) &= \begin{pmatrix} 1 & 0 & 0 \\ 0 & 1 & 0 \\ 0 & 0 & 1 \end{pmatrix} = I & (k_{1} &= 1) \\
\rho(T_{2}) &= \begin{pmatrix} 1 & 0 & 0 \\ 0 & 1 & 0 \\ 0 & 0 & -1 \end{pmatrix} = A & (k_{2} &= 1) \\
\rho(T_{3}) &= \begin{pmatrix} \tfrac{1}{2} & \tfrac{1}{2} & -\tfrac{11}{4} \\ 2 & 1 & 3 \\ 1 & 0 & -\tfrac{3}{2} \end{pmatrix} = B & (k_{3} &= 3) \\
\rho(T_{4}) &= \begin{pmatrix} \tfrac{1}{2} & \tfrac{1}{2} & \tfrac{11}{4} \\ 2 & 1 & -3 \\ 1 & 0 & \tfrac{3}{2} \end{pmatrix} = BA & (k_{4} &= 3) \\
\rho(T_{5}) &= \begin{pmatrix} \tfrac{1}{2} & \tfrac{1}{2} & \tfrac{11}{4} \\ 2 & 1 & -3 \\ -1 & 0 & -\tfrac{3}{2} \end{pmatrix} = -3I+3\,A-B+AB+BA & (k_{5} &= 3) \\
\rho(T_{6}) &= \begin{pmatrix} \tfrac{1}{2} & \tfrac{1}{2} & -\tfrac{11}{4} \\ 2 & 1 & 3 \\ -1 & 0 & \tfrac{3}{2} \end{pmatrix} = AB & (k_{6} &= 3) \\
\rho(T_{7}) &= \begin{pmatrix} -2 & -\tfrac{3}{2} & -3 \\ 0 & 2 & 8 \\ 0 & 1 & 0 \end{pmatrix} & (k_{7} &= 6) \\
 &= 8I-2\,A+6\,B-4\,AB-2\,BA+2\,BB-2\,ABB-2\,BAB && \\
\rho(T_{8}) &= \begin{pmatrix} -6 & -\tfrac{7}{2} & 0 \\ 8 & 2 & 0 \\ 0 & 0 & 0 \end{pmatrix} & (k_{8} &= 6) \\
 &= 9I-7\,A+7\,B-AB-2\,BA+5\,BB-3\,ABB-2\,BAB-2\,BABB && \\
\rho(T_{9}) &= \begin{pmatrix} -2 & -\tfrac{3}{2} & 3 \\ 0 & 2 & -8 \\ 0 & 1 & 0 \end{pmatrix} = 10I-6\,A+2\,B-4\,AB-4\,BA+2\,BB & (k_{9} &= 6) \\
\rho(T_{10}) &= \begin{pmatrix} 1 & -\tfrac{1}{4} & 1 \\ -4 & -2 & 6 \\ -2 & -\tfrac{3}{2} & -2 \end{pmatrix} & (k_{10} &= 6) \\
 &= -10I+10\,A-2\,B+3\,AB+3\,BA-2\,BB-BAB && \\
\rho(T_{11}) &= \begin{pmatrix} 1 & -\tfrac{1}{4} & -1 \\ -4 & -2 & -6 \\ -2 & -\tfrac{3}{2} & 2 \end{pmatrix} & (k_{11} &= 6) \\
 &= -2I-A-6\,B+3\,AB-2\,BA-2\,BB+3\,ABB+2\,BAB && \\
\rho(T_{12}) &= \begin{pmatrix} 0 & 2 & 0 \\ -8 & -4 & 0 \\ 0 & 0 & 0 \end{pmatrix} & (k_{12} &= 6) \\
 &= -4I+3\,A-4\,B-2\,AB-BA-4\,BB+3\,ABB+BAB+2\,BABB && \\
\rho(T_{13}) &= \begin{pmatrix} -2 & -\tfrac{3}{2} & -3 \\ 0 & 2 & 8 \\ 0 & -1 & 0 \end{pmatrix} = 4I+2\,A+4\,B-2\,AB-2\,BA-2\,BAB & (k_{13} &= 6) \\
\rho(T_{14}) &= \begin{pmatrix} -2 & -\tfrac{3}{2} & 3 \\ 0 & 2 & -8 \\ 0 & -1 & 0 \end{pmatrix} = 6I-2\,A-2\,AB-4\,BA+2\,ABB & (k_{14} &= 6) \\
\rho(T_{15}) &= \begin{pmatrix} 1 & -\tfrac{1}{4} & 1 \\ -4 & -2 & 6 \\ 2 & \tfrac{3}{2} & 2 \end{pmatrix} = I-A+3\,B-2\,AB+3\,BA+BB-3\,ABB-BAB & (k_{15} &= 6) \\
\rho(T_{16}) &= \begin{pmatrix} 1 & -\tfrac{1}{4} & -1 \\ -4 & -2 & -6 \\ 2 & \tfrac{3}{2} & -2 \end{pmatrix} = 5I-8\,A-B-2\,AB-2\,BA+BB+2\,BAB & (k_{16} &= 6) \\
\rho(T_{17}) &= \begin{pmatrix} 1 & \tfrac{3}{4} & -4 \\ 0 & -1 & 2 \\ 2 & \tfrac{1}{2} & -2 \end{pmatrix} = -3\,A+BAB & (k_{17} &= 6) \\
\rho(T_{18}) &= \begin{pmatrix} 1 & \tfrac{3}{4} & 4 \\ 0 & -1 & -2 \\ 2 & \tfrac{1}{2} & 2 \end{pmatrix} = -2I+3\,BA-ABB & (k_{18} &= 6) \\
\rho(T_{19}) &= \begin{pmatrix} 1 & \tfrac{3}{4} & -4 \\ 0 & -1 & 2 \\ -2 & -\tfrac{1}{2} & 2 \end{pmatrix} = -3I-3\,B+3\,AB-BB+ABB+BAB & (k_{19} &= 6) \\
\rho(T_{20}) &= \begin{pmatrix} 1 & \tfrac{3}{4} & 4 \\ 0 & -1 & -2 \\ -2 & -\tfrac{1}{2} & -2 \end{pmatrix} = -9I+7\,A-3\,B+3\,AB+3\,BA-BB & (k_{20} &= 6) \\
\rho(T_{21}) &= \begin{pmatrix} -2 & \tfrac{1}{4} & \tfrac{3}{2} \\ 4 & 1 & -4 \\ 0 & \tfrac{1}{2} & 1 \end{pmatrix} = 3I-3\,A+B-AB-BA+BB & (k_{21} &= 6) \\
\rho(T_{22}) &= \begin{pmatrix} -2 & \tfrac{1}{4} & -\tfrac{3}{2} \\ 4 & 1 & 4 \\ 0 & \tfrac{1}{2} & -1 \end{pmatrix} = I+3\,B-AB+BB-ABB-BAB & (k_{22} &= 6) \\
\rho(T_{23}) &= \begin{pmatrix} 2 & \tfrac{5}{2} & 4 \\ -4 & -4 & -2 \\ 0 & 0 & 0 \end{pmatrix} & (k_{23} &= 6) \\
 &= -11I+6\,A-5\,B+2\,AB+4\,BA-3\,BB+ABB+BAB+BABB && \\
\rho(T_{24}) &= \begin{pmatrix} 0 & -\tfrac{3}{4} & \tfrac{3}{2} \\ 0 & 2 & -4 \\ 0 & \tfrac{1}{2} & -1 \end{pmatrix} = 5I-2\,A+B-2\,AB-2\,BA+BB & (k_{24} &= 6) \\
\rho(T_{25}) &= \begin{pmatrix} 0 & -\tfrac{3}{4} & -\tfrac{3}{2} \\ 0 & 2 & 4 \\ 0 & \tfrac{1}{2} & 1 \end{pmatrix} = 5I-A+3\,B-2\,AB-BA+BB-ABB-BAB & (k_{25} &= 6) \\
\rho(T_{26}) &= \begin{pmatrix} -2 & \tfrac{1}{4} & -\tfrac{3}{2} \\ 4 & 1 & 4 \\ 0 & -\tfrac{1}{2} & 1 \end{pmatrix} = A+2\,B-BAB & (k_{26} &= 6) \\
\rho(T_{27}) &= \begin{pmatrix} 4 & \tfrac{1}{2} & 0 \\ -4 & -1 & 0 \\ 0 & 0 & 2 \end{pmatrix} = -3I+A-3\,B+2\,AB+2\,BA-BB+BAB & (k_{27} &= 6) \\
\rho(T_{28}) &= \begin{pmatrix} 4 & \tfrac{1}{2} & 0 \\ -4 & -1 & 0 \\ 0 & 0 & -2 \end{pmatrix} = -5I+3\,A-3\,B+2\,AB+2\,BA-BB+BAB & (k_{28} &= 6) \\
\rho(T_{29}) &= \begin{pmatrix} -1 & -\tfrac{1}{4} & 0 \\ 4 & -1 & 0 \\ -2 & -\tfrac{1}{2} & 0 \end{pmatrix} & (k_{29} &= 6) \\
 &= -5I+A-B+4\,AB+2\,BA+BB-ABB-BABB && \\
\rho(T_{30}) &= \begin{pmatrix} 0 & -\tfrac{3}{4} & -\tfrac{3}{2} \\ 0 & 2 & 4 \\ 0 & -\tfrac{1}{2} & -1 \end{pmatrix} = 2I+2\,A+2\,B-AB-BA-BAB & (k_{30} &= 6) \\
\rho(T_{31}) &= \begin{pmatrix} 0 & -\tfrac{3}{4} & \tfrac{3}{2} \\ 0 & 2 & -4 \\ 0 & -\tfrac{1}{2} & 1 \end{pmatrix} = 4I-A-AB-2\,BA+ABB & (k_{31} &= 6) \\
\rho(T_{32}) &= \begin{pmatrix} 2 & \tfrac{5}{2} & -4 \\ -4 & -4 & 2 \\ 0 & 0 & 0 \end{pmatrix} & (k_{32} &= 6) \\
 &= -7I+A-5\,B+2\,AB+BA-3\,BB+2\,ABB+2\,BAB+BABB && \\
\rho(T_{33}) &= \begin{pmatrix} -2 & \tfrac{1}{4} & \tfrac{3}{2} \\ 4 & 1 & -4 \\ 0 & -\tfrac{1}{2} & -1 \end{pmatrix} = -BA+ABB & (k_{33} &= 6) \\
\rho(T_{34}) &= \begin{pmatrix} -1 & -\tfrac{1}{4} & 0 \\ 4 & -1 & 0 \\ 2 & \tfrac{1}{2} & 0 \end{pmatrix} & (k_{34} &= 6) \\
 &= -4\,A+2\,B+AB+2\,BA+2\,BB-2\,ABB-BABB && .
\end{align*}



%──────────────────────────────────────────────────────────────────────
\section{Meinhold-72: $\bigl(G_2(2),\,\mathrm{PSL}(3,2)\bigr)$}

\textbf{Graph.} The Meinhold-72 orbital graph: $G=G_2(2)=\mathrm{PSU}(3,3){:}2$
($|G|=12096$) acting on the $72$ cosets of $H=\mathrm{PSL}(3,2)$ (the $36+36$
chiral Fano structures).

\textbf{Algebra.} Rank $8$, $\HH\cong\C^2\oplus M_2(\C)\oplus\C^2$.

\textbf{Highlights.} The block is generated by the mirror involution
$A=\rho(T_2)$ and the directed valency-7 orbital $B=\rho(T_3)$ with
\[
A^2=I,\qquad B^2+4B+13I=0 ;
\]
the roots $-2\pm3i$ of $x^2+4x+13$ are Gaussian integers of norm $13$.
Equivalently, the symmetric generator $\rho(T_4)=-4A-AB$ satisfies
$\rho(T_4)^2=13I$ and $A\rho(T_4)+\rho(T_4)A=-4I$.
Why 13? $13=7+2\cdot3$: $NN^{\mathsf T}=7I+2A_{\mathrm{SRG}}$ for the
$\mathrm{SRG}(36,21,12,12)$ on the chiral Fano structures, and
$3=\sqrt{k-\mu}$; the $-4$ is the $U_3(3)$-graph eigenvalue. (Full account:
\texttt{why13.pdf}.)

% auto-generated by hecke_catalog_full.g — do not edit
\paragraph{One-dimensional characters.}
\[
\begin{array}{c|cccc}
 & \chi_{1}\,(1) & \chi_{2}\,(1) & \chi_{10}\,(21) & \chi_{11}\,(21) \\\hline
T_{1} & 1 & 1 & 1 & 1 \\
T_{2} & 1 & -1 & -1 & 1 \\
T_{3} & 7 & 7 & 1 & 1 \\
T_{4} & 7 & -7 & -1 & 1 \\
T_{5} & 7 & 7 & 1 & 1 \\
T_{6} & 7 & -7 & -1 & 1 \\
T_{7} & 21 & -21 & 3 & -3 \\
T_{8} & 21 & 21 & -3 & -3 \\
\end{array}
\]

\subsubsection*{Block 1: $M_2(\C)$ on $\chi_{9}$ (degree 14)}
\noindent\textbf{Generators.}
\[
A = \rho(T_{2}) = \begin{pmatrix} 1 & 0 \\ 0 & -1 \end{pmatrix}, \qquad
B = \rho(T_{3}) = \begin{pmatrix} -2 & 1 \\ -9 & -2 \end{pmatrix} .
\]
\textbf{Complete relations} (monomial basis $I,\;A,\;B,\;AB$ of $M_2(\C)$):
\begin{gather*}
A^2-I = 0 ,\\
B^2+4\,B+13I = 0 ,\\
BA = -4\,A-AB ,\\
ABA = -4I-B ,\\
ABB = -13\,A-4\,AB .
\end{gather*}
\textbf{All orbital images} ($\rho(T_i)$, valency $k_i$, word in the generators):
\begin{align*}
\rho(T_{1}) &= \begin{pmatrix} 1 & 0 \\ 0 & 1 \end{pmatrix} = I & (k_{1} &= 1) \\
\rho(T_{2}) &= \begin{pmatrix} 1 & 0 \\ 0 & -1 \end{pmatrix} = A & (k_{2} &= 1) \\
\rho(T_{3}) &= \begin{pmatrix} -2 & 1 \\ -9 & -2 \end{pmatrix} = B & (k_{3} &= 7) \\
\rho(T_{4}) &= \begin{pmatrix} -2 & -1 \\ -9 & 2 \end{pmatrix} = -4\,A-AB & (k_{4} &= 7) \\
\rho(T_{5}) &= \begin{pmatrix} -2 & -1 \\ 9 & -2 \end{pmatrix} = -4I-B & (k_{5} &= 7) \\
\rho(T_{6}) &= \begin{pmatrix} -2 & 1 \\ 9 & 2 \end{pmatrix} = AB & (k_{6} &= 7) \\
\rho(T_{7}) &= \begin{pmatrix} 3 & 0 \\ 0 & -3 \end{pmatrix} = 3\,A & (k_{7} &= 21) \\
\rho(T_{8}) &= \begin{pmatrix} 3 & 0 \\ 0 & 3 \end{pmatrix} = 3I & (k_{8} &= 21) .
\end{align*}



%──────────────────────────────────────────────────────────────────────
\section*{Provenance}
\addcontentsline{toc}{section}{Provenance}
{\sloppy
All data blocks: GAP, exact cyclotomic arithmetic, by isotypic restriction
and spin-up of eigenvectors --- single script
\texttt{hecke\_catalog\_full.g} in \texttt{\~{}/git/hecke/}, console audit
trail \texttt{hecke\_catalog\_full.out}, LaTeX fragments
\texttt{catalog-data/*.tex} (machine-generated; do not edit). Methods and
proofs: \texttt{hecke-irreps.pdf}; the Meinhold story: \texttt{why13.pdf}.
The non-Gelfand pairs $\mathrm{PSL}(2,19)/S_3$ and $\mathrm{PSL}(2,31)/S_4$
are deliberately out of scope.\par}

\end{document}
